% concatenated from JD-main-current.tex
\documentclass{amsart}

\usepackage[utf8]{inputenc}
\usepackage{amsmath}
\usepackage{amssymb}
\usepackage{amsthm}
\usepackage{amscd}
\usepackage{amsfonts}
\usepackage{graphicx}
\usepackage{fancyhdr}
\usepackage{mathrsfs}
\usepackage{verbatim}
\usepackage{array}
\usepackage{booktabs}
\usepackage{enumitem}
\usepackage{xr-hyper}
\usepackage[colorlinks=true, linkcolor=blue, citecolor=blue, urlcolor=blue]{hyperref}
\usepackage{tikz-cd}
\usepackage[all,cmtip]{xy}
\usepackage{wasysym}
\usepackage{commath}
\usepackage{mathtools}
\usepackage{regexpatch,fancyvrb,xparse}
\usepackage{rotating}
\usepackage[margin=1.0in]{geometry}
\usepackage{csquotes}

\title{Pinned Jordan Decomposition of Characters and Depth-Zero Hecke Algebras}
\subjclass{20C33}
\keywords{finite reductive groups, Jordan decomposition of characters, Harish-Chandra induction}

\date{}

\theoremstyle{plain}
\numberwithin{equation}{section}
\newtheorem{theorem}{Theorem}[section]
\newtheorem{corollary}[theorem]{Corollary}
\newtheorem{claim}[theorem]{claim}
\newtheorem{lemma}[theorem]{Lemma}
\newtheorem{proposition}[theorem]{Proposition}
\theoremstyle{definition}
\newtheorem{definition}[theorem]{Definition}
\newtheorem{hypothesis}[theorem]{Hypothesis}
\newtheorem{finalremark}[theorem]{Final Remark}
\newtheorem{remark}[theorem]{Remark}
\newtheorem{conjecture}[theorem]{Conjecture}
\newtheorem{example}[theorem]{Example}
\newtheorem{question}{Question}
\topmargin-2cm
\newtheorem{Fact}[theorem]{Fact}

\newtheorem{notation}[theorem]{Notation}

\newcommand{\Sim}{\operatorname{Sim}}
\renewcommand{\L}{\mathcal{L}}
\renewcommand{\dim}{\operatorname{dim}}

\newcommand{\op}{\operatorname{op}}
\newcommand{\Ad}{\operatorname{Ad}}
\newcommand{\cP}{\mathcal{P}}
\newcommand{\cE}{\mathcal{E}}
\newcommand{\ab}{\operatorname{ab}}
\newcommand{\Z}{\mathbb{Z}}
\newcommand{\Q}{\mathbb{Q}}
\newcommand{\id}{\mathrm{id}}
\newcommand{\Out}{\operatorname{Out}}
\newcommand{\ff}{\mathfrak{f}}

\newcommand{\Ind}{\operatorname{Ind}}
\newcommand{\End}{\operatorname{End}}
\newcommand{\Ext}{\operatorname{Ext}}
\newcommand{\Qp}{\mathbb{Q}_{p}}
\newcommand{\Ql}{\mathbb{Q}_{\ell}}
\newcommand{\lz}{\ell\mathbb{Z}}
\newcommand{\Qlcl}{\overline{\mathbb{Q}}_{\ell}}
\newcommand{\Zp}{\mathbb{Z}_{p}}
\newcommand{\Ga}{\mathbb{G}_{a}}
\newcommand{\Gap}{\mathbb{G}_{a,\perf}}
\newcommand{\cpu}{\mathfrak{cpu}_{k}^{\circ}}
\newcommand{\Hom}{\operatorname{Hom}}
\newcommand{\perf}{\operatorname{perf}}
\newcommand{\spec}{\operatorname{Spec}}
\newcommand{\CH}{\operatorname{CH}}
\newcommand{\red}{\operatorname{red}}
\newcommand{\Res}{\operatorname{Res}}
\renewcommand{\l}{\ell}
\newcommand{\Pic}{\operatorname{Pic}}
\newcommand{\EqBr}{\operatorname{Eqbr}}
\newcommand{\Frob}{\operatorname{Frob}}
\newcommand{\Irr}{\operatorname{Irr}}
\newcommand{\Fun}{\operatorname{Fun}}
\newcommand{\ACh}{\operatorname{ACh}}
\newcommand{\Sh}{\operatorname{Sh}}
\newcommand{\CS}{\operatorname{CS}}
\newcommand{\fq}{\mathbb{F}_{q}}
\newcommand{\Id}{\operatorname{Id}}
\newcommand{\infl}{\operatorname{Inf}}
\newcommand{\Uch}{\operatorname{Uch}}
\newcommand{\J}{\mathbb{J}}
\newcommand{\LF}{L^{F}}
\newcommand{\GF}{G^{F}}
\newcommand{\Stab}{\operatorname{Stab}}
\newcommand{\sgn}{\operatorname{sgn}}
\newcommand{\ad}{\operatorname{ad}}
\newcommand{\der}{\operatorname{der}}
\newcommand{\cusp}{\operatorname{cusp}}
\newcommand{\fact}{\operatorname{Fact}}
\newcommand{\GL}{\operatorname{GL}}
\newcommand{\U}{\operatorname{U}}
\newcommand{\SL}{\operatorname{SL}}
\newcommand{\SO}{\operatorname{SO}}
\newcommand{\Sp}{\operatorname{Sp}}
\newcommand{\tr}{\operatorname{Tr}}
\newcommand{\Aut}{\operatorname{Aut}}
\DeclareMathOperator{\Reg}{Reg}

\newcommand{\rW}{\mathrm{W}}
\newcommand{\R}{\mathrm{R}}
\newcommand{\cA}{\mathcal{A}}
\newcommand{\cW}{\mathcal{W}}
\newcommand{\cH}{\mathcal{H}}
\newcommand{\cG}{\mathcal{G}}
\newcommand{\cR}{\mathcal{R}}

\newcommand{\bG}{\mathbf{G}}
\newcommand{\bL}{\mathbf{L}}
\newcommand{\bP}{\mathbf{P}}
\newcommand{\bM}{\mathbf{M}}
\newcommand{\bT}{\mathbf{T}}
\newcommand{\bB}{\mathbf{B}}
\newcommand{\bH}{\mathbf{H}}
\newcommand{\bX}{\mathbf{X}}
\newcommand{\bK}{\mathbf{K}}
\newcommand{\bU}{\mathbf{U}}
\newcommand{\bS}{\mathbf{S}}

\newcommand{\A}{\mathsf{A}}
\newcommand{\B}{\mathsf{B}}
\newcommand{\C}{\mathbb{C}}
\newcommand{\D}{\mathsf{D}}
\newcommand{\fE}{\mathsf{E}}
\newcommand{\F}{\mathbb{F}}
\newcommand{\G}{\mathsf{G}}
\newcommand{\triv}{1}

\begin{document}

\author{Prashant Arote}
\email{prashantarote123@gmail.com}
\author{Manish Mishra}
\email{manish@iiserpune.ac.in}
\address{
Department of Mathematics \\
Indian Institute for Science Education and Research \\
Dr.\ Homi Bhabha Road, Pashan \\
Pune 411 008 \\
India
}
\thanks{The second named author was partially supported by SERB Core Research Grant
(CRG/2022/000415) and  ANRF MATRICS Grant ANRF/ARGM/2025/000365/MTR.}
\begin{abstract}
We construct a pinned canonical Jordan decomposition of characters for finite
reductive groups in situations where the dual centralizers may be disconnected.
For a connected reductive group \(\bG\) over a finite field, equipped with a
pinning, and for a semisimple element \(s\in G^*\), we construct a uniquely
determined bijection
\[
\J_s:\cE(G,s)\xrightarrow{\sim}\Uch\bigl(C_{\bG^*}(s)^{F^*}\bigr).
\]
This refines Lusztig's orbit-valued Jordan decomposition for groups with
disconnected centre, and is characterized by compatibility with the
Deligne--Lusztig character formula and with Harish--Chandra series.  We then
extend the construction to the class of possibly disconnected reductive groups
with abelian component group whose rational components admit pinning-preserving
representatives.  For this class the correct unconditional statement is an
enriched disconnected Jordan decomposition: the target is the connected
unipotent Jordan datum together with the source Clifford class and the
corresponding projective Clifford label.  When the transported source Clifford
classes agree with the ordinary Clifford classes on the dual-centralizer side,
this enriched target specializes to the usual set of unipotent characters of the
corresponding disconnected dual centralizer.

The main technical input is a pinned-normalized construction of preferred
extensions of cuspidal unipotent characters to their inertia groups.  The
extension is not forced by abstract Clifford theory alone: after the pinning
fixes the relevant component action, Lusztig's preferred-extension convention
removes the remaining scalar ambiguity.  The construction then uses Clifford
theory, relative Weyl group comparison, Malle's matching, and connected and
disconnected forms of Howlett--Lehrer theory to extend the cuspidal Jordan
decomposition functorially to all Harish--Chandra series.

As an application, we give a pinned canonical form of the finite-field input in
Ohara's comparison of depth-zero Hecke algebra parameters with the unipotent
case.  More precisely, for a depth-zero Bernstein type
\((K_{x_0},\rho_{x_0})\), the affine Hecke algebra factor in the
AFMO--Morris presentation is canonically identified, after fixing the same
pinning, with the affine Hecke algebra factor attached to Ohara's corresponding
unipotent type.  
\end{abstract}

\maketitle

\tableofcontents
\section{Introduction}

The Jordan decomposition of characters is one of the central organizing
principles in the representation theory of finite reductive groups.  Let
\(\bG\) be a connected reductive group over \(\F_q\), with Frobenius \(F\), and
write \(G=\bG^F\).  If \((\bG^*,F^*)\) is dual to \((\bG,F)\), then
Deligne--Lusztig theory partitions \(\Irr(G)\) into Lusztig series
\[
\Irr(G)=\bigsqcup_{(s)}\cE(G,s),
\]
where \(s\) runs over semisimple conjugacy classes in \(G^*=\bG^{*F^*}\).
The Jordan decomposition problem asks for a natural parametrization of
\(\cE(G,s)\) by unipotent characters of the centralizer
\[
H=C_{\bG^*}(s)^{F^*}.
\]
Thus one seeks bijections
\[
J_s:\cE(G,s)\xrightarrow{\sim}\Uch(H),
\]
compatible with the character formulae coming from Deligne--Lusztig induction.

When \(Z(\bG)\) is connected, the centralizer \(C_{\bG^*}(s)\) is connected.
In this case Lusztig constructed the Jordan decomposition of characters
\cite{Lus84}, and Digne--Michel proved a uniqueness theorem for a normalized
family of such bijections \cite{DM90}.  Their normalization includes
compatibility with central quotients, products, and certain induction
properties, and it gives a canonical way to remove the ambiguity present in
Lusztig's original parametrization.

For groups with disconnected centre the situation is subtler.  In this case
\(C_{\bG^*}(s)\) need not be connected.  Lusztig proved in \cite{Lus88} that
one still has a natural orbit-valued Jordan decomposition
\[
\cE(G,s)\longrightarrow
\Uch\bigl(C_{\bG^*}(s)^{\circ F^*}\bigr)/
\bigl(C_{\bG^*}(s)^{F^*}/C_{\bG^*}(s)^{\circ F^*}\bigr).
\]
The fibres on the \(G\)-side are orbits under diagonal automorphisms, and the
orbits on the dual side are orbits of the component group of the centralizer.
This result is fundamental, but it does not by itself produce a bijection with
the irreducible unipotent characters of the disconnected finite group
\(C_{\bG^*}(s)^{F^*}\).  Passing from orbit data to actual irreducible
characters requires Clifford-theoretic choices, namely choices of extensions of
unipotent characters from the identity component to their stabilizers.

The canonical endoscopic Jordan decomposition of Lusztig and Yun
\cite{LusztigYun2020,LusztigYunCorrigendum2021} gives a different and more
geometric treatment of the same general problem.  Their construction uses
monodromic Hecke categories and character sheaves, and relates a fixed
semisimple block to unipotent objects for a corresponding endoscopic group.  The
present paper has a more representation-theoretic and application-driven aim:
we retain the dual-centralizer language used in depth-zero local Langlands
applications, fix a pinning to normalize the connected Jordan decomposition,
and record the necessary Clifford-theoretic data when parahoric finite
reductive quotients are disconnected.  Thus the two constructions have
different targets and different normalizations, and the results below should be
viewed as a pinned finite-field package tailored to the depth-zero applications,
not as a substitute for the Lusztig--Yun endoscopic construction.

The purpose of this paper is to make these choices canonical in a precise
pinned-normalized sense.  After a pinning has been fixed, the action of the
component group on the based root datum is fixed.  The remaining ambiguity is
the familiar Clifford-theoretic scalar ambiguity in extending invariant
characters; it is removed by Lusztig's preferred-extension convention, and by
Malle's compatible matching on the cuspidal side.  Thus the construction below
involves no further auxiliary choices.  We do not mean that the abstract
disconnected group alone forces a unique extension before these normalizations
are specified.  With this convention, we construct a pinned Jordan decomposition
of characters which refines Lusztig's orbit-valued decomposition and is
compatible with Harish--Chandra induction.

Our first main result concerns connected reductive groups, but allows the
dual centralizer \(C_{\bG^*}(s)\) to be disconnected.

\medskip
\noindent
\textbf{The connected-group theorem.}
Let \(\bG\) be connected reductive over \(\F_q\), equipped with a pinning
\(\cP\), and let \(s\in G^*\) be semisimple.  Put
\[
\bH=C_{\bG^*}(s),\qquad H=\bH^{F^*}.
\]
Theorem~\ref{thm:JD-connected-groups} constructs a uniquely determined
bijection
\[
\J_s:\cE(G,s)\xrightarrow{\sim}\Uch(H),
\]
depending only on the pinning, such that for every \(F^*\)-stable maximal torus
\(\bT^*\subseteq \bG^*\) containing \(s\) one has
\[
\bigl\langle R_{\bT^*}^{\bG}(s),\rho\bigr\rangle_G
=
\epsilon_{\bG}\epsilon_{\bH}
\bigl\langle R_{\bT^*}^{\bH}(1),\J_s(\rho)\bigr\rangle_H.
\]
Moreover, this bijection is compatible with Harish--Chandra series.  If
\((L,\tau)\) is a cuspidal pair with \(\tau\in\cE(L,s)\), and if
\[
H_L=C_{\bL^*}(s)^{F^*},
\qquad
u_\tau=\J_s^L(\tau)\in\Uch(H_L)_{\cusp},
\]
then \(\J_s\) restricts to a canonical bijection
\[
\Irr(G,(L,\tau))
\xrightarrow{\sim}
\Irr(H,(H_L,u_\tau)).
\]
Thus the Jordan decomposition is not merely a set-theoretic parametrization of
Lusztig series; it is compatible with the Harish--Chandra decomposition inside
each Lusztig series.

Our second main result extends this construction to a class of disconnected
finite reductive groups, but in an enriched form.  Let \(\bG\) be a possibly
disconnected reductive group over \(\F_q\), fix an \(F\)-stable pinning \(\cP\) of
\(\bG^\circ\), and assume that the finite component group is abelian and
satisfies the rational pinned-component condition of
Hypothesis~\ref{hyp:rational-pinned-component-condition}.  For
\[
s\in(\bG^*)^{\circ F^*}
\]
semisimple, we define a disconnected Lusztig series \(\cE(G,s)\) by restriction
to the identity component.  If
\[
H=C_{\bG^*}(s)^{F^*},
\]
then Theorem~\ref{thm:disc-JD-bijection} gives a canonical enriched bijection
\[
J^{\mathrm{enh}}_{G,s}:\cE(G,s)
\xrightarrow{\sim}
\Uch^{\mathrm{enh}}_{G,s}(H).
\]
On the cuspidal part, an irreducible character \(\rho\) lying above
\(\rho^\circ\in\cE(G^\circ,s)_{\cusp}\) is sent to
\[
\left[
\J_s^{G^\circ}(\rho^\circ),
[\alpha_{\rho^\circ}],
E_\rho
\right],
\]
where \([\alpha_{\rho^\circ}]\) and \(E_\rho\) are the source Clifford class and
source Clifford label.  On a general Harish--Chandra series the construction
preserves the corresponding relative Howlett--Lehrer label.  Thus the theorem
no longer requires an unproved equality between source and target Clifford
cocycles; that equality becomes precisely the extra condition under which the
enriched target reduces to the ordinary set \(\Uch(H)\).

Although the main application carried out in this paper concerns depth-zero
Hecke algebras, the disconnected counterpart of Jordan decomposition also has a
more direct Langlands-theoretic application.  In the second named author's
construction of a pinned local Langlands correspondence for depth-zero
supercuspidal representations of a connected reductive group over a
non-archimedean local field \(F_{\mathrm{loc}}\)
\cite{Mishra2026PinnedLLCDepthZero}, the finite representation attached to a
depth-zero type at a vertex \(x\) lives naturally on the full parahoric quotient
\[
G(F_{\mathrm{loc}})_x/G(F_{\mathrm{loc}})_{x,0+},
\]
which is in general a possibly disconnected finite reductive group.  The
pinned disconnected Jordan decomposition developed here supplies the finite
unipotent label on the dual-centralizer side.  In that construction it is the
finite step which separates the finite representation in a depth-zero type from
the toral part of the parameter and makes the passage to a cuspidal enhanced
unramified parameter canonical.  The inverse construction uses the same
finite-field bijection in the opposite direction.  Thus the enriched
Clifford-theoretic bookkeeping in Theorem~\ref{thm:disc-JD-bijection} is not
merely a technical refinement of Lusztig's orbit-valued decomposition: it is the
finite mechanism which makes the depth-zero supercuspidal LLC reversible after a
pinned normalization has been fixed.

The third main result is an application to Hecke algebras attached to
depth-zero Bernstein blocks of \(p\)-adic groups.  Let \(F_{\mathrm{loc}}\) be
a non-archimedean local field with residue field \(\mathfrak f\), and let
\(G\) be a connected reductive group over \(F_{\mathrm{loc}}\).  For a depth-zero
\(\mathfrak s\)-type \((K,\rho)=(K_{x_0},\rho_{x_0})\), the Hecke algebra
\(\cH_{\mathfrak s}(G)=\mathcal H(G(F),(K,\rho))\) admits the AFMO--Morris
presentation as an affine Hecke algebra crossed with a finite twisted group
algebra; see \cite{AFMO24a}.  Ohara's construction \cite{ohara2025} associates
to the same depth-zero datum a unipotent type \((K_{\theta,x_0},u_{x_0})\) for a
related group \(G_\theta\).  Ohara proves that the affine Hecke algebra factor in
\(\cH_{\mathfrak s}(G)\) is isomorphic to the affine Hecke algebra factor
appearing in the Hecke algebra of this unipotent type.  Theorem~\ref{thm:padic-depth0-to-unipotent-hecke}
records the pinned version of this affine-factor comparison.  Equivalently, one
may rewrite \(\cH_{\mathfrak s}(G)\) as its original source-side twisted group
algebra crossed with the corresponding unipotent affine Hecke factor.  We do
not claim a general isomorphism from \(\cH_{\mathfrak s}(G)\) to the full Hecke
algebra of the unipotent type, because the finite stabilizer group and its
cocycle in the crossed-product presentation need not agree with the unipotent
ones.  The point of the present paper is that the finite-field Jordan
decomposition used in the comparison of the affine parameters is no longer an
auxiliary non-canonical choice: it is the pinned canonical bijection constructed
here.

We now describe the strategy of the proofs.

The construction of Theorem~\ref{thm:JD-connected-groups} begins with the
cuspidal case.  Lusztig's orbit-valued map gives, for each cuspidal
\(\rho\in\cE(G,s)\), an orbit
\[
O\subseteq \Uch(H^\circ)_{\cusp},
\qquad H^\circ=(C_{\bG^*}(s)^\circ)^{F^*}.
\]
For \(u^\circ\in O\), let \(I_H(u^\circ)\) be its stabilizer in \(H\).  Since
the relevant component groups are abelian, Clifford theory says that the
irreducible characters of \(H\) lying above \(O\) are obtained from extensions
of \(u^\circ\) to \(I_H(u^\circ)\), twisted by characters of
\(I_H(u^\circ)/H^\circ\), and then induced to \(H\).  The main difficulty is
to choose these extensions canonically.

This is done using Lusztig's preferred-extension convention for Weyl-group
representations, together with the interpretation of unipotent characters via
generalized Springer theory and character sheaves.  We first handle cyclic
component groups, then glue the resulting preferred extensions for abelian
component groups, and finally pass to products and repeated quasi-simple
factors by a wreath-product argument.  This gives preferred extensions of
cuspidal unipotent characters to their inertia groups.

On the \(G\)-side, the fibre of Lusztig's orbit-valued map over \(O\) is an
orbit under the diagonal automorphism group
\[
A=G_{\ad}/G.
\]
Lusztig's Clifford-theoretic construction gives an orthogonality relation
between the stabilizer on the \(G\)-side and the stabilizer on the unipotent
side.  The pinning supplies the missing normalization: using semisimple
characters and Malle's results on automorphisms of cuspidal characters
\cite{Malle2017}, one attaches to each cuspidal \(\rho\) a canonical class in
the appropriate quotient of \(A\).  Under the orthogonality isomorphism this
class gives a character of the unipotent stabilizer, and hence selects a
specific irreducible character of \(H\) above the orbit \(O\).  This produces
the pinned cuspidal bijection
\[
\cE(G,s)_{\cusp}\xrightarrow{\sim}\Uch(H)_{\cusp}.
\]

The passage from cuspidal characters to all of \(\cE(G,s)\) uses
Harish--Chandra theory.  For a cuspidal pair \((L,\tau)\) with
\(\tau\in\cE(L,s)\), the cuspidal construction gives a unipotent cuspidal
character
\[
u_\tau\in\Uch(C_{\bL^*}(s)^{F^*}).
\]
A key point is that the relative Weyl groups on the two sides agree:
\[
W_\tau \cong W_{u_\tau}.
\]
This follows from the Weyl-equivariance of the pinned cuspidal correspondence.
Howlett--Lehrer theory identifies the constituents of
\(R_{\bL}^{\bG}(\tau)\) with \(\Irr(W_\tau)\).  On the dual centralizer side,
we use a disconnected version of Howlett--Lehrer theory, in stabilizer form,
to identify the constituents of
\(R_{H_L}^{H}(u_\tau)\) with \(\Irr(W_{u_\tau})\).  The isomorphism
\(W_\tau\cong W_{u_\tau}\) then gives the desired bijection of Harish--Chandra
series.  Taking the disjoint union over all cuspidal pairs gives
Theorem~\ref{thm:JD-connected-groups}.

The proof of Theorem~\ref{thm:disc-JD-bijection} follows the same philosophy,
but with an additional Clifford-theoretic layer coming from the disconnected
group \(G\) itself.  In the disconnected part of the paper we impose the
rational pinned-component condition: each rational component has a representative
whose conjugation action preserves the chosen pinning of \(G^\circ\).  Thus the
semidirect-product model determined by the pinning records the actual component
action on characters of the finite identity component, up to inner conjugacy by
\(G^\circ\).  For a cuspidal constituent
\[
\rho^\circ\in\cE(G^\circ,s)_{\cusp},
\]
the connected-group theorem gives a cuspidal unipotent correspondent
\[
u_0\in\Uch(C_{G^{\circ *}}(s)^{F^*})_{\cusp}.
\]
The pinned equivariance of the connected correspondence, applied to these
pinning-preserving rational representatives, implies that the stabilizers of
\(\rho^\circ\) and \(u_0\) in the component group agree.  We do not assert that
the source and target Clifford-label cocycles agree.  Instead, the source Clifford
class \([\alpha_{\rho^\circ}]\) and the corresponding projective label \(E_\rho\)
are retained as part of the enriched unipotent datum.

The full disconnected construction is then obtained by Harish--Chandra
induction.  Disconnected Lusztig series are shown to be unions of
Harish--Chandra series.  For a cuspidal pair \((L,\tau)\), the enriched cuspidal
datum records \(\J_s^{L^\circ}(\tau^\circ)\), the source Clifford class of
\(\tau^\circ\), and the Clifford label of \(\tau\).  The disconnected
Howlett--Lehrer/corner analysis then supplies the relative label of each
constituent of \(R_L^G(\tau)\).  Combining the enriched cuspidal datum with this
relative label gives the full enriched bijection of
Theorem~\ref{thm:disc-JD-bijection}.

We close with a guide to the organization of the paper.  The early sections
recall Deligne--Lusztig theory, Lusztig series, the Digne--Michel
normalization, Lusztig's disconnected-centre construction, and the
Howlett--Lehrer parametrization of Harish--Chandra series.  The technical heart
of the paper is the construction of preferred extensions of cuspidal unipotent
characters and the resulting cuspidal pinned Jordan decomposition.  The
connected-group theorem is then obtained by combining this cuspidal
construction with relative Weyl group comparison and Howlett--Lehrer theory.
The later sections extend the construction to disconnected reductive groups
with abelian component group satisfying the rational pinned-component condition,
in the enriched sense explained above.  Readers mainly interested in the
Hecke-algebra application to \(p\)-adic groups may read
Theorems~\ref{thm:JD-connected-groups} and \ref{thm:disc-JD-bijection} as the
finite-field input and then proceed directly to
\S\ref{sec:padic-depth0-to-unipotent-hecke}.  For the complementary
Langlands-theoretic application of the disconnected theorem to depth-zero
supercuspidal representations, see \cite{Mishra2026PinnedLLCDepthZero}.

\section{Notation and conventions}
\label{sec:notation}

Let \(\fq\) be a finite field with \(q=p^r\) elements, and fix a prime
\(\ell\neq p\).  We fix, once and for all, an isomorphism
\[
\mathbb C \cong \overline{\mathbb Q}_{\ell}.
\]
All characters will be viewed over either field through this fixed
isomorphism.

Unless otherwise stated, \(\bG\) denotes a connected reductive algebraic group
defined over \(\fq\), and
\[
F:\bG\longrightarrow \bG
\]
denotes the Frobenius morphism defining its \(\fq\)-rational structure.  We
write
\[
G=\bG^F.
\]
More generally, if \(\bH\) is an \(F\)-stable algebraic subgroup of \(\bG\), we
write \(H=\bH^F\).  On the dual side, rational points are always taken with
respect to the dual Frobenius \(F^*\).

We write \(Z(\bH)\), \(C_{\bH}(x)\), and \(N_{\bH}(\bK)\) for the center of
\(\bH\), the centralizer of an element \(x\in\bH\), and the normalizer of a
subgroup \(\bK\leq\bH\), respectively.  If \(\bS\) is a torus in \(\bH\), then
\[
W_{\bH}(\bS)\coloneqq N_{\bH}(\bS)/C_{\bH}(\bS)
\]
denotes the corresponding Weyl group.  If \(\bS\) is \(F\)-stable, we shall
also use the finite rational Weyl group
\[
W_H(S)\coloneqq N_H(S)/C_H(S),
\qquad S=\bS^F,
\]
when this distinction is relevant.

For a finite group \(X\), we write \(\Irr(X)\) for the set of irreducible
characters of \(X\).  If \(X\) is the group of rational points of a connected
reductive group, we write \(\Uch(X)\) for its set of unipotent characters.  If
\(i:H\to K\) is a homomorphism of finite groups, and if \(i(H)\) is clear from
context, we shall sometimes write \(K/H\) for \(K/i(H)\).  This notation will
only be used when no ambiguity can result.

We now fix notation for root data.  Let \(\bB_0\) be an \(F\)-stable Borel
subgroup of \(\bG\), and let \(\bT_0\subseteq\bB_0\) be an \(F\)-stable maximal
torus.  We write
\[
X(\bT_0)=\Hom(\bT_0,\mathbb G_m),
\qquad
X^\vee(\bT_0)=\Hom(\mathbb G_m,\bT_0)
\]
for the character and cocharacter groups.  They are equipped with the usual
perfect pairing
\[
\langle\ ,\ \rangle:
X(\bT_0)\times X^\vee(\bT_0)\longrightarrow \mathbb Z,
\]
defined by
\[
\lambda(\mu(a))=a^{\langle \lambda,\mu\rangle}
\qquad
(a\in\mathbb G_m).
\]
Let \(\Phi=\Phi(\bG,\bT_0)\) be the corresponding root system and
\(\Phi^\vee\) the coroot system.  The Borel subgroup \(\bB_0\) determines a
positive system \(\Phi^+\subseteq\Phi\), and we denote the associated set of
simple roots by \(\Delta\).  Thus the triple
\[
(\bG,\bB_0,\bT_0)
\]
determines the based root datum
\[
\bigl(
X(\bT_0),\Phi,\Delta,
X^\vee(\bT_0),\Phi^\vee,\Delta^\vee
\bigr).
\]

Let \((\bG^*,F^*)\) be dual to \((\bG,F)\).  Thus there are \(F\)- and
\(F^*\)-stable pairs \((\bB_0,\bT_0)\) and \((\bB_0^*,\bT_0^*)\), together with
an isomorphism
\[
\delta:X(\bT_0)\xrightarrow{\sim} X^\vee(\bT_0^*)
\]
which identifies the root datum of \(\bG\) with the dual root datum of
\(\bG^*\).  In particular,
\[
\delta(\Phi)=\Phi^{*\vee}.
\]
If \(\alpha\in\Phi\), we denote by \(\alpha^*\in\Phi^*\) the root determined by
\[
\delta(\alpha)=\alpha^{*\vee}.
\]
The isomorphism \(\delta\) is assumed to be compatible with Frobenius in the
sense that
\[
\delta(\lambda\circ F|_{\bT_0})
=
F^*|_{\bT_0^*}\circ \delta(\lambda)
\qquad
(\lambda\in X(\bT_0)).
\]
Equivalently, the induced isomorphism of root data commutes with the Frobenius
actions on the character and cocharacter lattices.

We write
\[
W=W_{\bG}(\bT_0),
\qquad
W^*=W_{\bG^*}(\bT_0^*).
\]
The duality above identifies \(W\) with \(W^*\), by sending the simple
reflection associated with \(\alpha\) to the simple reflection associated with
\(\alpha^*\).  This identification is compatible with the Frobenius actions,
and hence restricts to an isomorphism
\[
W^F\xrightarrow{\sim} (W^*)^{F^*}.
\]

Finally, we fix our convention for pinnings.  A pinning of \(\bG\) will be
written
\[
\cP=(\bG,\bB,\bT,\{x_\alpha\}_{\alpha\in\Delta}),
\]
where \(\bB\) is a Borel subgroup, \(\bT\subseteq\bB\) is a maximal torus, and
\(x_\alpha:\mathbb G_a\to \bX_\alpha\) is a choice of root homomorphism for
each simple root \(\alpha\).  When the pinning is \(F\)-stable, the dual group
\(\bG^*\) is equipped with the corresponding \(F^*\)-stable dual pinning
\[
\cP^*=(\bG^*,\bB^*,\bT^*,\{x_{\alpha^*}\}_{\alpha^*\in\Delta^*}).
\]

If \(s\in T^*=\bT^{*F^*}\) is semisimple, we put
\[
\bH\coloneqq C_{\bG^*}(s),
\qquad
H\coloneqq \bH^{F^*},
\qquad
\bH^\circ\coloneqq C_{\bG^*}(s)^\circ,
\qquad
H^\circ\coloneqq (\bH^\circ)^{F^*}.
\]
The pinning \(\cP^*\) determines a natural pinning of the connected reductive
group \(\bH^\circ\), obtained from the root subgroups of \(\bG^*\) whose roots
vanish on \(s\).

We shall also use the following notation for pinned automorphisms.  Define
\[
\A^*
\coloneqq
\Stab_{\Aut(\bG^*,\cP^*)}(s).
\]
Via duality of pinned root data, \(\A^*\) corresponds to a subgroup
\[
\A\subseteq \Aut(\bG,\cP).
\]
When no confusion is possible, we use the same symbol for an algebraic
automorphism and for the induced automorphism of the corresponding finite group
of rational points.

We shall use the following terminology for algebraic groups.  A connected
semisimple algebraic group \(\mathbf X\) over \(\overline{\mathbb F}_p\) will
be called \emph{quasi-simple} if its absolute root system is irreducible,
equivalently if its adjoint quotient
\[
        \mathbf X_{\mathrm{ad}}=\mathbf X/Z(\mathbf X)
\]
is a simple adjoint algebraic group.  Thus a quasi-simple algebraic group need
not be simply connected or adjoint.  When it is adjoint we shall usually say
\emph{simple adjoint}.

If \(F\) is a Frobenius endomorphism and \(\mathbf X\) is connected
semisimple, we shall say that \(\mathbf X\) is \(F\)-almost simple if \(F\)
acts transitively on the set of absolute quasi-simple factors of \(\mathbf X\).
Equivalently, up to isogeny, \(\mathbf X\) is obtained by restriction of
scalars from a quasi-simple group over a finite extension of the ground field.

These terms refer to algebraic groups.  We shall not use the finite-group
theoretic meanings of ``quasisimple'' and ``almost simple'' unless this is
explicitly stated.  In particular, when \(X=\mathbf X^F\), the phrase
``\(\mathbf X\) is quasi-simple'' is a statement about the algebraic group
\(\mathbf X\), not about the abstract finite group \(X\).

We shall also use the following terminology for the disconnected groups that
arise as centralizers of semisimple elements.

\begin{definition}[Semisimple-centralizer type]
\label{def:semisimple-centralizer-type}
A possibly disconnected reductive algebraic group \(\mathbf H\) is said to be
of \emph{semisimple-centralizer type} if there exist a connected reductive
algebraic group \(\mathbf G_0\) and a semisimple element \(s\in \mathbf G_0\)
such that
\[
        \mathbf H=C_{\mathbf G_0}(s).
\]
If a Frobenius endomorphism \(F\) is present, we shall say that \((\mathbf H,F)\)
is of \emph{\(F\)-stable semisimple-centralizer type} if \(\mathbf H\) is
\(F\)-stable in such a realization.
\end{definition}

\begin{lemma}[Digne--Michel]
\label{lem:DM-semisimple-centralizer-splits}
Let \(\mathbf H=C_{\mathbf G_0}(s)\) be of semisimple-centralizer type, with
\(\mathbf G_0\) connected reductive and \(s\in \mathbf G_0\) semisimple.  Then
\[
1\longrightarrow \mathbf H^\circ
\longrightarrow \mathbf H
\longrightarrow \pi_0(\mathbf H)
\longrightarrow 1
\]
splits.  Equivalently, there is a finite subgroup \(A\subset \mathbf H\),
mapping isomorphically onto \(\pi_0(\mathbf H)\), such that
\[
        \mathbf H=\mathbf H^\circ\rtimes A .
\]
If, moreover, \((\mathbf H,F)\) is of \(F\)-stable semisimple-centralizer type
and \(H=\mathbf H^F\), \(H^\circ=(\mathbf H^\circ)^F\), then the finite exact
sequence
\[
1\longrightarrow H^\circ
\longrightarrow H
\longrightarrow (\pi_0(\mathbf H))^F
\longrightarrow 1
\]
also splits.  Finally, after choosing a maximal torus of \(\mathbf H^\circ\)
and a positive system for its root system, the component group is represented
at the Weyl-group level by the subgroup preserving this positive system.
\end{lemma}

\begin{proof}
The algebraic splitting is \cite[Theorem~0.1]{DigneMichel2026Centralizers}, and
the splitting on rational points is \cite[Theorem~0.2]{DigneMichel2026Centralizers}.
The final assertion is the Weyl-group splitting in
\cite[Proposition~1.1]{DigneMichel2026Centralizers}.
\end{proof}

\begin{definition}[Pinned semidirect component model]
\label{def:pinned-semidirect-component-model}
Let \(\mathbf H\) be a possibly disconnected reductive algebraic group with
Frobenius \(F\), let \(H=\mathbf H^F\), and let \(H^\circ=(\mathbf H^\circ)^F\).
Fix a pinning of \(\mathbf H^\circ\).  If \(J\le H\) is a subgroup containing
\(H^\circ\), we shall say that \(J\) has a \emph{pinned semidirect component
model} if there is a finite subgroup \(A_J\le J\), mapping isomorphically onto
\(J/H^\circ\), such that
\[
        J=H^\circ\rtimes A_J,
\]
and the conjugation action of \(A_J\) on \(\mathbf H^\circ\), after passing to
outer automorphism classes, is represented by the corresponding pinned
automorphisms of \(\mathbf H^\circ\).  Thus the semidirect-product notation is
part of the specified model; it is not asserted for an arbitrary disconnected
reductive group.
\end{definition}

\begin{remark}
\label{rmk:DM-supplies-pinned-component-models}
If \((\mathbf H,F)\) is of \(F\)-stable semisimple-centralizer type, then
Lemma~\ref{lem:DM-semisimple-centralizer-splits} supplies such a semidirect
component model for \(H=\mathbf H^F\).  Intersecting the Digne--Michel
complement with the inverse image of a subgroup of \(H/H^\circ\) supplies the
same kind of model for every intermediate subgroup \(J\) with
\(H^\circ\le J\le H\).  This is the justification for the semidirect-product
notation used for centralizers of semisimple elements below.
\end{remark}

\section{Preliminaries: Deligne--Lusztig Theory and Jordan Decomposition of Characters}
\label{sec:DL-theory-JD}

In this section we recall the parts of Deligne--Lusztig theory and Jordan
decomposition of characters that will be used later.  This summary is adapted
from \cite[\textsection 2.3]{AroteMishra22b}.  For further details we refer to
\cite{DelLus76,Lus84,Lus88,Book:DigneandMichel,Book:GeckandMalle}.

Let \(\bG\) be a connected reductive group over \(\fq\), with Frobenius
morphism \(F:\bG\to \bG\), and put \(G=\bG^F\).  Let
\((\bG^*,F^*)\) be dual to \((\bG,F)\), and put
\(G^*=\bG^{*F^*}\).

\begin{proposition}\label{prop:geometric-conjugacy}
Let \(\bT\) be an \(F\)-stable maximal torus of \(\bG\), and let
\(\bT^*\) be an \(F^*\)-stable maximal torus of \(\bG^*\) dual to \(\bT\).
There is a natural bijection
\[
\left\{
\begin{array}{c}
G\text{-conjugacy classes of pairs }(\bT,\theta),\\
\theta\in\Irr(T),\ T=\bT^F
\end{array}
\right\}
\longleftrightarrow
\left\{
\begin{array}{c}
G^*\text{-conjugacy classes of pairs }(\bT^*,s),\\
s\in T^*,\ T^*=\bT^{*F^*}
\end{array}
\right\}.
\]
\end{proposition}

For an \(F\)-stable maximal torus \(\bT\) of \(\bG\) and a character
\(\theta\in\Irr(T)\), Deligne and Lusztig associate a virtual character
\(R_{\bT}^{\bG}(\theta)\) of \(G\), constructed from \(\ell\)-adic cohomology
with compact support, where \(\ell\neq p\).  If \((\bT,\theta)\) corresponds
to \((\bT^*,s)\) under Proposition~\ref{prop:geometric-conjugacy}, we also
write
\[
R_{\bT^*}^{\bG}(s) \coloneqq R_{\bT}^{\bG}(\theta).
\]

\begin{remark}
Suppose that \(\bT\) is contained in an \(F\)-stable Borel subgroup
\(\bB\) of \(\bG\).  Then \(R_{\bT}^{\bG}(\theta)\) is an actual character,
rather than merely a virtual character, and it agrees with ordinary
Harish--Chandra induction from \(B=\bB^F\):
\[
R_{\bT}^{\bG}(\theta)=\Ind_{B}^{G}(\widetilde{\theta}),
\]
where \(\widetilde{\theta}\) denotes the inflation of \(\theta\) to \(B\).
\end{remark}

\begin{definition}[Lusztig series]
Let \(s\in G^*\) be semisimple.  The Lusztig series, or rational series,
associated with \(s\) is
\[
\cE(G,s)
\coloneqq
\left\{
\rho\in\Irr(G)
\ \middle|\
\langle R_{\bT^*}^{\bG}(s),\rho\rangle\neq 0
\text{ for some }F^*\text{-stable maximal torus }
\bT^*\subseteq \bG^*
\text{ with }s\in T^*
\right\}.
\]
\end{definition}

\begin{definition}[Unipotent character]
A character \(\rho\in\Irr(G)\) is called \emph{unipotent} if
\[
\langle R_{\bT^*}^{\bG}(1),\rho\rangle\neq 0
\]
for some \(F^*\)-stable maximal torus \(\bT^*\subseteq \bG^*\).  The set of
unipotent characters of \(G\) is denoted by \(\Uch(G)\).  Thus
\[
\Uch(G)=\cE(G,1).
\]
\end{definition}

\begin{theorem}[Partition into Lusztig series]
Every irreducible character \(\rho\in\Irr(G)\) occurs in some
Deligne--Lusztig virtual character.  More precisely, there exist an
\(F\)-stable maximal torus \(\bT\subseteq\bG\) and a character
\(\theta\in\Irr(T)\) such that
\[
\langle R_{\bT}^{\bG}(\theta),\rho\rangle\neq 0.
\]
Moreover,
\[
\Irr(G)=\bigsqcup_{(s)}\cE(G,s),
\]
where \((s)\) runs over the semisimple conjugacy classes of \(G^*\).
\end{theorem}

\begin{definition}[Sign]
Let \(\bG\) be a connected algebraic group defined over \(\fq\), with
Frobenius morphism \(F\).  Its sign is
\[
\epsilon_{\bG}\coloneqq (-1)^{\operatorname{rk}_{\fq}(\bG)},
\]
where \(\operatorname{rk}_{\fq}(\bG)\) denotes the \(F\)-split rank of \(\bG\).
\end{definition}

For connected reductive groups with connected center, Lusztig
\cite{Lus84} introduced the Jordan decomposition of characters.  It gives a
parametrization of the Lusztig series \(\cE(G,s)\) by unipotent characters of
the finite group \(C_{\bG^*}(s)^{F^*}\).  Lusztig's construction does not, by
itself, determine a unique bijection with all the functorial properties needed
below.  Digne and Michel resolved this ambiguity by proving the following
uniqueness theorem.

\begin{theorem}[{\cite[Theorem 7.1]{DM90}}]
\label{thm:JD-connected-center}
There exists a unique collection of bijections
\[
J_s^G:\cE(G,s)\longrightarrow \Uch(H),
\]
where \(\bG\) runs over connected reductive groups with connected center,
\(s\in G^*\) is semisimple, and
\[
\bH\coloneqq C_{\bG^*}(s),
\qquad
H\coloneqq \bH^{F^*},
\]
satisfying the following properties.

\begin{enumerate}
    \item For every \(F^*\)-stable maximal torus \(\bT^*\leq \bH\) and every
    \(\rho\in\cE(G,s)\), one has
    \[
    \langle R_{\bT^*}^{\bG}(s),\rho\rangle
    =
    \epsilon_{\bG}\epsilon_{\bH}
    \langle R_{\bT^*}^{\bH}(1_{T^*}),J_s^G(\rho)\rangle .
    \]

    \item If \(s=1\) and \(\rho\in\Uch(G)\), then:
    \begin{enumerate}
        \item[(a)] the Frobenius eigenvalues \(\omega_{\rho}\) and
        \(\omega_{J_1^G(\rho)}\) are equal;

        \item[(b)] if \(\rho\) lies in the principal series, then
        \(\rho\) and \(J_1^G(\rho)\) correspond to the same character of the
        Iwahori--Hecke algebra.
    \end{enumerate}

    \item If \(z\in Z(G^*)\), then
    \[
    J_{sz}^G(\rho\otimes\widehat z)=J_s^G(\rho)
    \]
    for every \(\rho\in\cE(G,s)\), where \(\widehat z\) is the linear character
    of \(G\) corresponding to \(z\).

    \item Let \(\bL^*\) be an \(F^*\)-stable Levi subgroup of \(\bG^*\) such
    that \(\bH\leq\bL^*\), and let \(\bL\leq\bG\) be the dual Levi subgroup.
    Then the following diagram commutes:
    \[
    \begin{tikzcd}
    \cE(G,s)\arrow{r}{J_s^G}
    & \Uch(H)  \\
    \cE(L,s) \arrow{u}{R_L^G}\arrow{r}{J_s^L}
    & \Uch(H)\arrow{u}{\Id}.
    \end{tikzcd}
    \]

    \item Suppose that \(\bG\) is of type \(\fE_8\), and that \(\bH\) is of
    type \(\fE_7\A_1\), respectively of type \(\fE_6\A_2\).  Let
    \(\bL\leq\bG\) be a Levi subgroup of type \(\fE_7\), respectively
    \(\fE_6\), whose dual \(\bL^*\) is contained in \(\bH\).  Then the
    following diagram commutes:
    \[
    \begin{tikzcd}
    \mathbb{Z}\cE(G,s)\arrow{r}{J_s^G}
    & \mathbb{Z}\Uch(H)  \\
    \mathbb{Z}\cE(L,s)_{c} \arrow{u}{R_{\bL}^{\bG}}\arrow{r}{J_s^L}
    & \mathbb{Z}\Uch(L^*)_{c}\arrow{u}{R_{\bL^*}^{\bH}},
    \end{tikzcd}
    \]
    where the subscript \(c\) denotes the subgroup spanned by the cuspidal
    part of the corresponding Lusztig series.

    \item Let \(\bT_1\leq Z(\bG)\) be an \(F\)-stable central torus, and let
    \[
    \pi_1:\bG\longrightarrow \bG_1\coloneqq \bG/\bT_1
    \]
    be the natural epimorphism.  Let \(s_1\in G_1^*\) and suppose that
    \(s=\pi_1^*(s_1)\).  Put
    \[
    \bH_1\coloneqq C_{\bG_1^*}(s_1).
    \]
    Then the following diagram commutes:
    \[
    \begin{tikzcd}
    \cE(G,s)\arrow{r}{J_s^G}
    & \Uch(H) \arrow{d}{} \\
    \cE(G_1,s_1) \arrow{u}{}
    \arrow{r}{J_{s_1}^{G_1}}
    & \Uch(H_1),
    \end{tikzcd}
    \]
    where the left vertical map is inflation along \(G\to G_1\), and the right
    vertical map is restriction along the natural embedding \(H_1\hookrightarrow H\).

    \item If \(\bG=\prod_i \bG_i\) is a direct product of \(F\)-stable
    subgroups \(\bG_i\), and \(s=\prod_i s_i\), then
    \[
    J_s^G=\prod_i J_{s_i}^{G_i}.
    \]
\end{enumerate}
\end{theorem}
Lusztig \cite{Lus88} extended the Jordan decomposition of characters to
connected reductive groups whose center need not be connected.  The construction
uses a regular embedding into a connected reductive group with connected center.
We recall the relevant notation and auxiliary results, referring to
\cite{Lus88} for the proofs.

\begin{remark}
Let
\[
\pi:\bG\longrightarrow \bG_{\ad}
\]
be the adjoint quotient.  The finite group \(G_{\ad}\) acts on \(G\) by
conjugation through representatives in \(\bG_{\ad}\), and hence
\(G_{\ad}/\pi(G)\) acts on \(\Irr(G)\).  This action extends linearly to
virtual characters.  It preserves each Lusztig series \(\cE(G,s)\), since it
preserves the corresponding Deligne--Lusztig virtual characters
\(R_{\bT^*}^{\bG}(s)\).  Thus \(G_{\ad}/\pi(G)\) acts naturally on
\(\cE(G,s)\).
\end{remark}

\begin{definition}[Regular embedding]
Let \(\bG\) be a connected reductive group defined over \(\fq\).  A morphism
\[
i:\bG\longrightarrow \bG'
\]
defined over \(\fq\) is called a \emph{regular embedding} if:
\begin{itemize}
    \item \(\bG'\) is a connected reductive group over \(\fq\) with connected
    center;

    \item \(i\) identifies \(\bG\) with a closed subgroup of \(\bG'\);

    \item \(i(\bG)\) and \(\bG'\) have the same derived subgroup.
\end{itemize}
\end{definition}

Let \(i:\bG\hookrightarrow \bG'\) be a regular embedding.  We write
\(G'=\bG'^{F}\), and identify \(G\) with its image in \(G'\).  By
\cite[Theorem~1.7.12]{Book:GeckandMalle}, the dual morphism is a surjective
homomorphism
\[
i^*:\bG'^*\longrightarrow \bG^*
\]
defined over \(\fq\).  Its kernel
\[
\bK\coloneqq \ker(i^*)
\]
is an \(F^*\)-stable torus contained in \(Z(\bG'^*)\).  Put
\[
K\coloneqq \bK^{F^*}.
\]

There is a canonical isomorphism
\[
K \cong \Hom(G'/G,\Qlcl^\times).
\]
Thus \(K\) acts on \(\Irr(G')\) by tensoring with the corresponding linear
characters of \(G'/G\).  Under this action, if \(k\in K\), then tensoring by
the character corresponding to \(k\) sends
\[
\cE(G',s') \quad\text{onto}\quad \cE(G',ks').
\]

For \(s'\in G'^*=\bG'^{*F^*}\), define
\[
K_{s'}\coloneqq
\{\,k\in K\mid ks' \text{ is }G'^*\text{-conjugate to }s'\,\}.
\]
Equivalently, \(K_{s'}\) is the stabilizer of the Lusztig series
\(\cE(G',s')\) under the above action of \(K\).

\begin{lemma}\label{lem:H-component-group}
Let \(s'\in G'^*\), set
\[
s=i^*(s')\in G^*,
\qquad
\bH=C_{\bG^*}(s),
\qquad
H=\bH^{F^*},
\qquad
H^\circ=(\bH^\circ)^{F^*}.
\]
Then there is a canonical isomorphism
\[
H/H^\circ \cong K_{s'}.
\]
Consequently, \(H/H^\circ\) acts on \(\cE(G',s')\).
\end{lemma}

\begin{proof}
For \(x\in H\), choose a lift \(\dot{x}\in \bG'^*\) with
\(i^*(\dot{x})=x\).  Since \(x\) centralizes \(s=i^*(s')\), the element
\[
s'^{-1}\dot{x}s'\dot{x}^{-1}
\]
belongs to \(\bK\).  This construction induces a well-defined homomorphism
\[
H/H^\circ\longrightarrow K,
\]
whose image is precisely \(K_{s'}\).  The resulting map is the required
isomorphism.
\end{proof}

\begin{proposition}[{cf. \cite[Proposition~2.3.15]{Book:GeckandMalle}}]
\label{prop:H-action-on-unipotent-characters}
Let \(i:\bG\hookrightarrow \bG'\) be a regular embedding, and let
\(s'\in G'^*\).  Put
\[
s=i^*(s'),
\qquad
\bH=C_{\bG^*}(s),
\qquad
\bH'=C_{\bG'^*}(s').
\]
Then \(i^*\) induces a surjective morphism
\[
\bH'\twoheadrightarrow \bH^\circ
\]
with kernel \(\bK\).  Moreover, pullback along this morphism gives a canonical
bijection
\[
\Uch(H^\circ)\xrightarrow{\sim}\Uch(H'),
\qquad
\eta\longmapsto \eta\circ i^*|_{H'}.
\]
Via the identification \(H/H^\circ\cong K_{s'}\), this bijection transports the
natural \(H/H^\circ\)-action on \(\Uch(H^\circ)\) to the corresponding action
on \(\Uch(H')\).
\end{proposition}

The compatibility of the connected-center Jordan decomposition with central
twists implies the following equivariance statement.

\begin{lemma}[{\cite[Proposition~8.1]{Lus88}}]
\label{lem:regular-embedding-JD-equivariance}
The Jordan decomposition bijection for \(G'\),
\[
J_{s'}^{G'}:\cE(G',s')\xrightarrow{\sim}\Uch(H'),
\]
is compatible with the \(H/H^\circ\)-action, where \(H/H^\circ\) is identified
with \(K_{s'}\) as in Lemma~\ref{lem:H-component-group}.
\end{lemma}

\begin{remark}\label{rmk:regular-embedding-adjoint-action}
Let \(i:\bG\hookrightarrow \bG'\) be a regular embedding.  By
\cite[Remark~1.7.6]{Book:GeckandMalle}, \(i\) induces isomorphisms
\[
\bG_{\ad}\cong \bG'_{\ad},
\qquad
G_{\ad}\cong G'_{\ad}.
\]
Hence the conjugation action of \(G'\) on \(G\) induces a natural surjective
homomorphism
\[
G'/G \twoheadrightarrow G_{\ad}/\pi(G).
\]
In particular, the action of \(G'/G\) on \(\Irr(G)\) factors through the
action of \(G_{\ad}/\pi(G)\).

On the other hand, the canonical isomorphism
\[
K\cong \Hom(G'/G,\Qlcl^\times)
\]
identifies
\[
G'/G \cong \Hom(K,\Qlcl^\times).
\]
Restricting characters from \(K\) to \(K_{s'}\), and using
Lemma~\ref{lem:H-component-group}, gives a natural surjective homomorphism
\[
\alpha:G'/G
\twoheadrightarrow
\Hom(K_{s'},\Qlcl^\times)
\cong
\Hom(H/H^\circ,\Qlcl^\times).
\]
The preceding paragraph shows that \(\alpha\) factors through
\(G_{\ad}/\pi(G)\).
\end{remark}

\begin{theorem}[Multiplicity-freeness {\cite[\S 10]{Lus88}}]
Let \(i:\bG\hookrightarrow \bG'\) be a regular embedding.  Then, for every
\(\rho'\in\Irr(G')\), the restriction \(\rho'|_G\) is multiplicity-free.
\end{theorem}

We shall also use the following elementary Clifford-theoretic fact, in the
form used by Lusztig.

\begin{remark}[{\cite[\textsection 9]{Lus88}}]
\label{rmk:multi-regular-embedding-action}
Let \(N\triangleleft M\) be finite groups, and suppose that \(M/N\) is
abelian.  The conjugation action of \(M\) on \(N\) induces an action of
\(M/N\) on \(\Irr(N)\).  The dual group
\[
(M/N)^\vee\coloneqq \Hom(M/N,\Qlcl^\times)
\]
acts on \(\Irr(M)\) by tensoring with linear characters.

Assume that, for every \(\rho'\in\Irr(M)\), the restriction \(\rho'|_N\) is
multiplicity-free.  Then there is a unique bijection between orbit sets
\[
\Irr(N)/(M/N)
\longleftrightarrow
\Irr(M)/(M/N)^\vee
\]
with the following property.  If \(\mathcal O\subseteq\Irr(N)\) corresponds to
\(\mathcal O'\subseteq\Irr(M)\), and if
\(\rho_\circ\in\mathcal O\), \(\rho'_\circ\in\mathcal O'\), then
\[
\rho'_\circ|_N=\sum_{\rho\in\mathcal O}\rho,
\qquad
\Ind_N^M(\rho_\circ)=\sum_{\rho'\in\mathcal O'}\rho'.
\]
Moreover, the stabilizer of \(\rho_\circ\) in \(M/N\) and the stabilizer of
\(\rho'_\circ\) in \((M/N)^\vee\) are orthogonal under the natural pairing
\[
M/N\times (M/N)^\vee\longrightarrow \Qlcl^\times.
\]
\end{remark}

With this notation in place, we can now state Lusztig's disconnected-center
version of Jordan decomposition \cite[Proposition~5.1]{Lus88}.
\begin{theorem}[Jordan decomposition for disconnected centralizers]
\label{thm:Lusztig-orbit-JD-disconnected-centralizers}
Let \(s\in G^*\) be semisimple, and put
\[
\bH\coloneqq C_{\bG^*}(s),
\qquad
H\coloneqq \bH^{F^*},
\qquad
H^\circ\coloneqq (\bH^\circ)^{F^*}.
\]
Fix a regular embedding \(i:\bG\hookrightarrow\bG'\), and choose
\(s'\in G'^*\) such that \(i^*(s')=s\).  Then Lusztig constructs a
surjective map
\[
J_s=J_s^{(G,G')}:
\cE(G,s)\twoheadrightarrow \Uch(H^\circ)/(H/H^\circ)
\]
with the following properties.

\begin{enumerate}
    \item The fibers of \(J_s\) are precisely the
    \(G_{\ad}/\pi(G)\)-orbits on \(\cE(G,s)\).

    \item Let \(\mathcal O\subseteq \Uch(H^\circ)\) be an
    \(H/H^\circ\)-orbit, and let
    \[
    \Gamma=\Stab_{H/H^\circ}(u)
    \]
    be the stabilizer of any element \(u\in\mathcal O\).  Then
    \[
    |J_s^{-1}(\mathcal O)|=|\Gamma|.
    \]

    \item If \(\rho\in \cE(G,s)\) and \(J_s(\rho)=\mathcal O\), then, for
    every \(F^*\)-stable maximal torus \(\bT^*\subseteq\bG^*\) containing
    \(s\), one has
    \[
    \big\langle R_{\bT^*}^{\bG}(s),\rho\big\rangle_G
    =
    \epsilon_{\bG}\epsilon_{\bH^\circ}
    \sum_{u\in\mathcal O}
    \big\langle
    R_{\bT^*}^{\bH^\circ}(1_{T^*}),u
    \big\rangle_{H^\circ}.
    \]
\end{enumerate}
\end{theorem}

\begin{remark}
We recall how the map \(J_s\) in Theorem~\ref{thm:Lusztig-orbit-JD-disconnected-centralizers} is
constructed.  Let
\[
K=\ker(i^*:\bG'^*\to\bG^*)
\]
and choose \(s'\in G'^*\) with \(i^*(s')=s\).  By Clifford theory and the
multiplicity-freeness of restriction from \(G'\) to \(G\), there is a natural
bijection of orbit sets
\[
\cE(G,s)/(G_{\ad}/\pi(G))
\;\cong\;
\left(\bigcup_{k\in K}\cE(G',s'k)\right)/K.
\]
Since \(K\) acts by central twists, the right-hand side may be identified with
\[
\cE(G',s')/K_{s'},
\]
where
\[
K_{s'}=\{\,k\in K\mid s'k \text{ is }G'^*\text{-conjugate to }s'\,\}.
\]
By Lemma~\ref{lem:H-component-group}, we have
\[
K_{s'}\cong H/H^\circ.
\]
The connected-center Jordan decomposition for \(G'\) gives
\[
J_{s'}^{G'}:\cE(G',s')\xrightarrow{\sim}\Uch(H'),
\qquad
H'=C_{\bG'^*}(s')^{F^*},
\]
and Lemma~\ref{lem:regular-embedding-JD-equivariance} says that this bijection is
\(K_{s'}\)-equivariant.  Hence it descends to a bijection
\[
\cE(G',s')/K_{s'}
\;\xrightarrow{\sim}\;
\Uch(H')/K_{s'}.
\]
Finally, Proposition~\ref{prop:H-action-on-unipotent-characters} identifies \(\Uch(H')\) with
\(\Uch(H^\circ)\), compatibly with the action of
\(K_{s'}\cong H/H^\circ\).  Composing these maps gives
\[
\cE(G,s)/(G_{\ad}/\pi(G))
\;\xrightarrow{\sim}\;
\Uch(H^\circ)/(H/H^\circ),
\]
and \(J_s\) is the resulting orbit-valued map on \(\cE(G,s)\).
\end{remark}

\begin{theorem}[{\cite[Theorem~3.2.22]{Book:GeckandMalle}}]
\label{fact:cuspidal}
Let \(s\in G^*\) be semisimple, and put
\[
\bH=C_{\bG^*}(s),
\qquad
H^\circ=(\bH^\circ)^{F^*}.
\]
Let \(\rho\in\cE(G,s)\), and let
\[
J_s(\rho)=\mathcal O\subseteq \Uch(H^\circ)
\]
be its image under Theorem~\ref{thm:Lusztig-orbit-JD-disconnected-centralizers}.  Choose
\(u_\rho\in\mathcal O\).  Then \(\rho\) is cuspidal if and only if the
following two conditions hold:
\begin{enumerate}
    \item \(u_\rho\) is cuspidal;

    \item the maximal \(\fq\)-split subtorus of \(Z(\bH^\circ)^\circ\) is
    contained in \(Z(\bG^*)^\circ\).  Equivalently,
    \[
    \operatorname{rank}_{\fq}\bigl(Z(\bG^*)^\circ\bigr)
    =
    \operatorname{rank}_{\fq}\bigl(Z(\bH^\circ)^\circ\bigr).
    \]
\end{enumerate}
The first condition is independent of the choice of \(u_\rho\in\mathcal O\),
since cuspidality is preserved under the action of \(H/H^\circ\).
\end{theorem}

The following elementary consequence will be used later.  In fact, the
conclusion follows directly from the description of Levi subgroups as
centralizers of split central tori; the cuspidality hypothesis is not needed.

\begin{corollary}
Let \(\bL^*\) be an \(F^*\)-stable Levi factor of an \(F^*\)-stable parabolic
subgroup of \(\bG^*\), and suppose that \(s\in L^*\).  Then
\[
C_{\bL^*}(s)^\circ
\]
is an \(F^*\)-stable Levi subgroup of
\[
C_{\bG^*}(s)^\circ .
\]
\end{corollary}

\begin{proof}
Let \(A_{\bL^*}\) be the maximal \(\fq\)-split torus in
\(Z(\bL^*)^\circ\).  Since \(\bL^*\) is an \(F^*\)-stable Levi subgroup, we
may write
\[
\bL^*=C_{\bG^*}(A_{\bL^*}).
\]
Because \(s\in L^*\), the torus \(A_{\bL^*}\) centralizes \(s\).  Hence
\(A_{\bL^*}\subseteq C_{\bG^*}(s)^\circ\).  Therefore
\[
C_{\bL^*}(s)^\circ
=
C_{C_{\bG^*}(s)^\circ}(A_{\bL^*}).
\]
The centralizer of a torus in a connected reductive group is a Levi subgroup.
Thus \(C_{\bL^*}(s)^\circ\) is a Levi subgroup of
\(C_{\bG^*}(s)^\circ\).  Since both \(s\) and \(\bL^*\) are \(F^*\)-stable, this
Levi subgroup is \(F^*\)-stable.
\end{proof}

\begin{lemma}[{\cite[Lemma~11.2.1]{Book:DigneandMichel}}]
\label{lem:connected-center-levi-centralizer}
Let \(\bG\) be a connected reductive group over \(\fq\).
\begin{enumerate}
    \item If \(Z(\bG)\) is connected, then the center of every Levi subgroup
    of \(\bG\) is connected.

    \item If \(Z(\bG)\) is connected, then the centralizer of every semisimple
    element in \(\bG^*\) is connected.
\end{enumerate}
\end{lemma}
\medskip
\section{Harish--Chandra Series, Endomorphism Algebras, and Relative Weyl Groups}
\label{sec:HC-series-endomorphism-algebras}

In this section we recall the standard material on Harish--Chandra induction,
Harish--Chandra series, and the associated relative Weyl groups.  We also fix
the notation used later for the Howlett--Lehrer parametrization of
Harish--Chandra series.  The final result of the section is a uniqueness
criterion for certain Harish--Chandra series characters; it will be used later
in the construction of the Jordan decomposition.

We begin with the basic setup.  Let \(\bG\) be a connected reductive algebraic
group over \(\fq\), with Frobenius morphism \(F:\bG\to\bG\), and put
\(G=\bG^F\).  Let \(\bP\) be an \(F\)-stable parabolic subgroup of \(\bG\) with
an \(F\)-stable Levi decomposition
\[
\bP=\bL\bU,
\]
where \(\bL\) is a Levi subgroup and \(\bU\) is the unipotent radical.  We
write
\[
R_{\bL}^{\bG}
\qquad\text{and}\qquad
{}^*R_{\bL}^{\bG}
\]
for Harish--Chandra induction and Harish--Chandra restriction, respectively.

\begin{definition}[Harish--Chandra series]
Let \(\bL\) be an \(F\)-stable Levi subgroup of an \(F\)-stable parabolic
subgroup of \(\bG\), and let \(\tau\in\Irr(L)\) be cuspidal, where
\(L=\bL^F\).  The Harish--Chandra series associated with the cuspidal pair
\((L,\tau)\) is
\[
\Irr(G,(L,\tau))
\coloneqq
\left\{
\rho\in\Irr(G)
\ \middle|\
\big\langle {}^*R_{\bL}^{\bG}(\rho),\tau\big\rangle_L\neq 0
\right\}.
\]
Equivalently, \(\Irr(G,(L,\tau))\) is the set of irreducible constituents of
\(R_{\bL}^{\bG}(\tau)\).
\end{definition}

The Harish--Chandra series form a partition of \(\Irr(G)\); see
\cite[Corollary~3.1.17]{Book:GeckandMalle}.  More precisely, if
\((L,\tau)\) and \((M,\lambda)\) are cuspidal pairs that are not
\(G\)-conjugate, then
\[
\Irr(G,(L,\tau))\cap \Irr(G,(M,\lambda))=\emptyset,
\]
and every irreducible character of \(G\) belongs to a unique Harish--Chandra
series up to \(G\)-conjugacy of the cuspidal pair.

The following standard result relates Harish--Chandra series to Lusztig
series; see \cite[Corollary~3.3.21]{Book:DigneandMichel}.

\begin{proposition}
\label{prop:Lusztig-series-union-HC}
Let \(s\in G^*\) be semisimple.  Then the Lusztig series \(\cE(G,s)\) is a
union of Harish--Chandra series:
\[
\cE(G,s)
=
\bigsqcup_{(L,\tau)\in \Sigma_G(s)}
\Irr(G,(L,\tau)).
\]
Here \(\Sigma_G(s)\) denotes the set of \(G\)-conjugacy classes of cuspidal
pairs \((L,\tau)\) such that, for a dual Levi subgroup
\(\bL^*\subseteq\bG^*\) with \(s\in L^*\), one has
\[
\tau\in \cE(L,s).
\]
\end{proposition}

We next recall the relative Weyl group attached to a cuspidal pair.  Let
\(A_{\bL}\) be the maximal \(\fq\)-split torus in \(Z(\bL)\).  Then
\[
\bL=C_{\bG}(A_{\bL}).
\]
Define
\[
W_{\bG}(A_{\bL})
\coloneqq
N_G(A_{\bL})/L.
\]
This group acts on \(L\), and hence on \(\Irr(L)\), through representatives in
\(N_G(A_{\bL})\).  For a cuspidal character \(\tau\in\Irr(L)\), set
\[
W_{\tau}(\bG)
\coloneqq
\{\,w\in W_{\bG}(A_{\bL})\mid {}^w\tau\simeq \tau\,\}.
\]
Equivalently, \(W_{\tau}(\bG)\) is the stabilizer of the character
\(\chi_\tau\) of \(\tau\).  When the ambient group \(\bG\) is clear from
context, we write simply \(W_\tau\).

The structure of Harish--Chandra series is governed by Howlett--Lehrer theory.
In its most useful form for us, this theory describes the endomorphism algebra
of \(R_{\bL}^{\bG}(\tau)\) in terms of an Iwahori--Hecke algebra attached to
the relative Weyl group \(W_\tau\).  The possible twisting cocycle that appears
in the general construction is trivial in the finite reductive group setting;
this follows from Lusztig's work in the connected-center case and from Geck's
reduction to that case.  We refer to \cite{HowlettLehrer} and
\cite[\S 3.2]{Book:GeckandMalle} for details.

\begin{theorem}[Howlett--Lehrer]
\label{thm:Howlett-Lehrer}
Let \(\bL\) be the Levi factor of an \(F\)-stable parabolic subgroup of
\(\bG\), and let \(\tau\in\Irr(L)\) be cuspidal.  Then
\[
\End_G\!\bigl(R_{\bL}^{\bG}(\tau)\bigr)^{\mathrm{op}}
\]
is naturally described by the Howlett--Lehrer Hecke algebra attached to
\(W_\tau(\bG)\).  In particular, by the Tits deformation theorem, the simple
modules of this endomorphism algebra are naturally parametrized by
\(\Irr(W_\tau(\bG))\).  Consequently, there is a bijection
\[
I_{L,\tau}^{G}:\Irr(W_\tau(\bG))
\xrightarrow{\sim}
\Irr(G,(L,\tau)),
\qquad
\phi\longmapsto \rho_\phi .
\]
\end{theorem}

We shall use this parametrization together with its compatibility with
Harish--Chandra induction.  The following form is the one needed later; see
\cite[Theorem~3.2.7]{Book:GeckandMalle}.

\begin{theorem}[Howlett--Lehrer comparison theorem]
\label{thm:Howlett-Lehrer-comparison}
Let \((L,\tau)\) be a cuspidal pair of \(G\).  For every \(F\)-stable Levi
subgroup \(\bM\) of an \(F\)-stable parabolic subgroup of \(\bG\), with
\[
\bL\leq \bM\leq \bG,
\]
there is a bijection
\[
I_{L,\tau}^{M}:
\Irr(W_{\tau}(\bM))
\xrightarrow{\sim}
\Irr(M,(L,\tau)),
\qquad
\phi\longmapsto \rho_\phi ,
\]
and these bijections are compatible with induction in the following sense:
\[
\begin{tikzcd}
\Irr(W_{\tau}(\bG))
\arrow[r,"I_{L,\tau}^{G}"]
& \mathbb{Z}\Irr(G,(L,\tau)) \\
\Irr(W_{\tau}(\bM))
\arrow[u,"\Ind_{W_{\tau}(\bM)}^{W_{\tau}(\bG)}"]
\arrow[r,"I_{L,\tau}^{M}"]
& \mathbb{Z}\Irr(M,(L,\tau))
\arrow[u,"R_{\bM}^{\bG}"']
\end{tikzcd}
\]
Here the horizontal maps are extended linearly to Grothendieck groups on the
right-hand side.
\end{theorem}

Equivalently, if \(\phi\in\Irr(W_\tau(\bM))\), then
\[
R_{\bM}^{\bG}(\rho_\phi)
=
\sum_{\psi\in\Irr(W_\tau(\bG))}
\big\langle
\Ind_{W_{\tau}(\bM)}^{W_{\tau}(\bG)}(\phi),\psi
\big\rangle
\rho_\psi .
\]
Thus the decomposition of Harish--Chandra induced characters is controlled by
ordinary character induction for the corresponding relative Weyl groups.

We now record the promised uniqueness criterion.  The following result and its
proof were provided to us by Jay Taylor.

\begin{proposition}[Jay Taylor]\label{prop:unique-HC-char}
Let $(L,\tau)$ be a cuspidal pair in $G$.
Suppose $\phi,\phi' \in \Irr(W_{\tau}(\bG))$ satisfy the following:

\begin{enumerate}[label={\normalfont(\arabic*)}]
	\item ${}^*R_{\bM}^{\bG}(\rho_{\phi}) = {}^*R_{\bM}^{\bG}(\rho_{\phi'})$ for any proper $F$-stable Levi subgroup $\bL \leqslant \bM < \bG$, of an $F$-stable parabolic subgroup of $\bG$,
	\item $\rho_{\phi}(1) = \rho_{\phi'}(1)$.
\end{enumerate}

Then $\rho_{\phi} = \rho_{\phi'}$ except possibly when $W_{\tau}(\bG)$ is irreducible of type $\mathsf{E}_7$ or $\mathsf{E}_8$ and $\{\phi,\phi'\}$ is an exceptional family of characters of dimension $512$ or $4096$ as in \cite[6.3.6]{geck-pfeiffer:2000:characters-of-finite-coxeter-groups}.
\end{proposition}

\begin{proof}
Put
\[
\cW:=W_{\tau}(\bG).
\]
If $\cW$ is trivial, then there is nothing to prove. We therefore assume that
$\cW\neq 1$.

We first translate the hypothesis on Harish--Chandra restrictions into a
statement about restrictions of characters of $\cW$.  Let $\bM$ be an
$F$-stable Levi subgroup with $\bL\leq \bM\leq \bG$, and put
\[
\cW_{\bM}:=W_{\tau}(\bM).
\]
We extend the Howlett--Lehrer bijection
\[
I_{L,\tau}^{M}:\Irr(\cW_{\bM})\longrightarrow \Irr(M,(L,\tau))
\]
linearly to virtual characters.  By the comparison theorem and Frobenius
reciprocity, the projection of ${}^*R_{\bM}^{\bG}(\rho_{\psi})$ to the
Harish--Chandra series $\Irr(M,(L,\tau))$ is
\begin{equation}\label{eq:HC-restriction-W}
I_{L,\tau}^{M}\!\left(\Res_{\cW_{\bM}}^{\cW}\psi\right),
\qquad \psi\in \Irr(\cW).
\end{equation}
Indeed, for $\eta\in\Irr(\cW_{\bM})$, the comparison theorem gives
\[
R_{\bM}^{\bG}\bigl(I_{L,\tau}^{M}(\eta)\bigr)
=
I_{L,\tau}^{G}\!\left(\Ind_{\cW_{\bM}}^{\cW}\eta\right),
\]
and hence
\[
\left\langle {}^*R_{\bM}^{\bG}(\rho_\psi),
I_{L,\tau}^{M}(\eta)\right\rangle_M
=
\left\langle \psi,\Ind_{\cW_{\bM}}^{\cW}\eta\right\rangle_{\cW}
=
\left\langle \Res_{\cW_{\bM}}^{\cW}\psi,\eta\right\rangle_{\cW_{\bM}}.
\]
Thus hypothesis~\textup{(1)} implies
\begin{equation}\label{eq:proper-parabolic-restrictions}
\Res_{\cW'}^{\cW}\phi
=
\Res_{\cW'}^{\cW}\phi'
\end{equation}
for every proper parabolic subgroup $\cW'<\cW$.  In particular, taking
$\bM=\bL$, so that $\cW_{\bL}=1$, gives
\begin{equation}\label{eq:degree-of-W-chars}
\phi(1)=\phi'(1).
\end{equation}

We shall also use the degree formula for the constituents of
$R_{\bL}^{\bG}(\tau)$.  Let $c_{\psi}$ denote the Schur element attached to
the irreducible character $\psi$ of the Howlett--Lehrer Hecke algebra.  By
Lusztig's degree formula, see \cite[Cor.~8.7]{Lus84} and
\cite[\S8.1.8]{geck-pfeiffer:2000:characters-of-finite-coxeter-groups},
there is a nonzero constant $C_{\bG,\bL,\tau}$, independent of $\psi$, such
that
\begin{equation}\label{eq:degree-schur-element}
\rho_{\psi}(1)
=
C_{\bG,\bL,\tau}\frac{\psi(1)}{c_{\psi}}.
\end{equation}
Combining hypothesis~\textup{(2)} with \eqref{eq:degree-of-W-chars} and
\eqref{eq:degree-schur-element}, we obtain
\begin{equation}\label{eq:equal-schur-elements}
c_{\phi}=c_{\phi'}.
\end{equation}

We now reduce to the case where $\cW$ is irreducible.  Write
\[
\cW=\cW_1\times\cdots\times \cW_r
\]
as a product of its irreducible Coxeter components, and write
\[
\phi=\phi_1\boxtimes\cdots\boxtimes \phi_r,
\qquad
\phi'=\phi'_1\boxtimes\cdots\boxtimes \phi'_r.
\]
If $r>1$, then each factor $\cW_i$ is a proper parabolic subgroup of $\cW$,
arising from an intermediate Levi subgroup in the Howlett--Lehrer
correspondence.  Hence \eqref{eq:proper-parabolic-restrictions} gives
\[
\frac{\phi(1)}{\phi_i(1)}\,\phi_i
=
\Res_{\cW_i}^{\cW}\phi
=
\Res_{\cW_i}^{\cW}\phi'
=
\frac{\phi'(1)}{\phi'_i(1)}\,\phi'_i.
\]
Since $\phi_i$ and $\phi'_i$ are irreducible characters of $\cW_i$, this
implies $\phi_i=\phi'_i$.  This holds for every $i$, and therefore
$\phi=\phi'$.  We may therefore assume from now on that $\cW$ is irreducible.

Suppose first that $\cW$ is of type $\mathsf A_n$, with $n\geq 1$.  If neither
$\phi$ nor $\phi'$ is the trivial character $\mathbf 1_{\cW}$, then
\eqref{eq:proper-parabolic-restrictions} implies $\phi=\phi'$ by
\cite[Cor.~5.4.8]{geck-pfeiffer:2000:characters-of-finite-coxeter-groups}.
We may therefore assume that one of the two characters is $\mathbf 1_{\cW}$.
By \eqref{eq:degree-of-W-chars}, both characters have degree one, so
\[
\{\phi,\phi'\}\subseteq \{\mathbf 1_{\cW},\varepsilon\},
\]
where $\varepsilon$ is the sign character.  If $n\geq 2$, then restriction to
a rank-one parabolic subgroup distinguishes $\mathbf 1_{\cW}$ from
$\varepsilon$, and hence $\phi=\phi'$.  If $n=1$, then the corresponding
Hecke algebra has parameter $q^k$ for some $k\geq 1$, and the two Schur
elements are
\[
c_{\mathbf 1_{\cW}}=\Phi_2(q^k),
\qquad
c_{\varepsilon}=q^{-k}\Phi_2(q^k).
\]
These are distinct.  Therefore \eqref{eq:equal-schur-elements} again forces
$\phi=\phi'$.

If $\cW$ is of type $\mathsf B_n$, $\mathsf C_n$, or $\mathsf D_n$, in the
usual irreducible ranges, then \eqref{eq:proper-parabolic-restrictions}
already implies $\phi=\phi'$ by
\cite[Thm.~6.2.9]{geck-pfeiffer:2000:characters-of-finite-coxeter-groups}.
It remains to consider the exceptional crystallographic types.  Let $\cW$ be
of type $\mathsf G_2$.  By the classification of the possible failures of
separation by proper parabolic restrictions, the only possible distinct pair
satisfying \eqref{eq:proper-parabolic-restrictions} is
\[
\{\phi,\phi'\}=\{\phi_{2,1},\phi_{2,2}\},
\]
where $\phi_{2,1}$ and $\phi_{2,2}$ are the two irreducible characters of
degree $2$, distinguished by their $b$-invariants.  According to
\cite[Table~II]{Lus84} and \cite[Thm.~8.6]{Lus84}, the possible parameters for
the corresponding Hecke algebra are of the form
\[
(q,q^{2k-1}),\qquad k\in\{1,2,5\}.
\]
By \cite[Thm.~8.3.4]{geck-pfeiffer:2000:characters-of-finite-coxeter-groups},
the Schur elements are
\[
c_{\phi_{2,b}}
=
2q^{-2k+1}\,
\Phi_3(q^{k+b-2})\,
\Phi_6(q^{k-b+1}),
\qquad b\in\{1,2\}.
\]
Thus it suffices to compare the two products
\[
\Phi_3(q^{k-1})\Phi_6(q^k)
\quad\text{and}\quad
\Phi_3(q^k)\Phi_6(q^{k-1}).
\]
For $k=1,2,5$ these factor as follows:
\[
\begin{array}{c|c|c}
k &
\Phi_3(q^{k-1})\Phi_6(q^k) &
\Phi_3(q^k)\Phi_6(q^{k-1}) \\
\hline
1 &
3\Phi_6(q) &
\Phi_3(q) \\
2 &
\Phi_3(q)\Phi_{12}(q) &
\Phi_3(q)\Phi_6(q)^2 \\
5 &
\Phi_3(q)\Phi_6(q)^2\Phi_{12}(q)\Phi_{30}(q) &
\Phi_3(q)\Phi_{15}(q)\Phi_{24}(q).
\end{array}
\]
For $q=2$ the two entries in each row are unequal by direct evaluation.  For
$q\neq 2$, Zsigmondy's theorem gives a primitive prime divisor of
$\Phi_6(q)$, $\Phi_{12}(q)$, and $\Phi_{30}(q)$ respectively in the three
rows; such a prime cannot divide any of the cyclotomic factors appearing on
the other side.  Hence
\[
c_{\phi_{2,1}}\neq c_{\phi_{2,2}},
\]
contradicting \eqref{eq:equal-schur-elements}.  Therefore $\phi=\phi'$ in
type $\mathsf G_2$ as well.

Finally suppose that $\cW$ is of type
\[
\mathsf F_4,\quad \mathsf E_6,\quad \mathsf E_7,\quad \text{or}\quad
\mathsf E_8.
\]
By \cite[6.3.6]{geck-pfeiffer:2000:characters-of-finite-coxeter-groups},
proper parabolic restrictions determine irreducible characters in these types
except for the exceptional pairs occurring in types $\mathsf E_7$ and
$\mathsf E_8$.  These exceptional pairs have degrees $512$ and $4096$,
respectively.  Thus, outside precisely the exceptional cases stated in the
proposition, \eqref{eq:proper-parabolic-restrictions} forces
$\phi=\phi'$, and hence $\rho_{\phi}=\rho_{\phi'}$.
\end{proof}

Let $\bT$ be an $F$-stable maximal torus contained in $\bL$, and let $\theta : T \to \overline{\mathbb{Q}}_{\ell}^{\times}$ be a character. Let $\bL^{*}$ be an $F^{*}$-stable Levi subgroup of $\bG^{*}$ corresponding to the Levi subgroup $\bL$ of $\bG$. Note that $\bL^{*}$ is determined up to $G^{*}$-conjugacy.
Let $(\bT^{*}, s)$ be a pair corresponding to $(\bT, \theta)$ as in Proposition~\ref{prop:geometric-conjugacy}, with $\bT^{*} \subset \bL^{*}$.

\begin{lemma}\label{lem:dual-stabilizers-geometric-pairs}
Let $[(\bT, \theta)]$ denote the $L$-conjugacy class of $(\bT, \theta)$, and let $[(\bT^{*}, s)]$ denote the $L^{*}$-conjugacy class of $(\bT^{*}, s)$. Then the isomorphism
\[
\delta : W_{\bG}(\bT_{0}) \longrightarrow W_{\bG^{*}}(\bT_{0}^{*}),
\]
induces an isomorphism
\[
\Stab_{W_{\bG}(A_{\bL})} \left( [(\bT, \theta)] \right) \cong \Stab_{W_{\bG^{*}}(A_{\bL^{*}})} \left( [(\bT^{*}, s)] \right).
\]
\end{lemma}
\begin{proof}
By \cite[Proposition~2.1]{HowlettLehrer}, we have
\[
W_{\bG}(A_{\bL}) \cong \frac{N_{W_{\bG}(\bT_{0})^{F}}(W_{\bL}(\bT_{0})^{F})}{W_{\bL}(\bT_{0})^{F}}.
\]
Thus, $W_{\bG}(A_{\bL}) \cong W_{\bG^{*}}(A_{\bL^{*}})$ under the isomorphism $\delta$, from which the result follows.
\end{proof}

\begin{lemma}\label{lem:centralizer-stabilizer-geometric-pair}
Let $s \in T^{*}$ be a semisimple element of $L^{*}$ and let $\tau\in \cE(L,s)$ be a cuspidal representation.
Let $\bH = C_{\bG^{*}}(s)$ be the centraliser of $s$ in $\bG^{*}$, and let $\bH_{\bL} = C_{\bL^{*}}(s)$ be the centraliser of $s$ in $\bL^{*}$.
Then
\[
\Stab_{W_{\bG^{*}}(A_{\bL^{*}})} \left( [(\bT^{*}, s)] \right) \cong \Stab_{W_{\bH}(A_{\bL^{*}})} \left( [(\bT^{*}, 1)] \right),
\]
where
\[
W_{\bH}(A_{\bL^{*}}) = \frac{N_{H}(A_{\bL^{*}})}{H_{\bL}}.
\]
\end{lemma}

\begin{proof}
Since $A_{\bL^{*}}$ is the maximal $\mathbb{F}_{q}$-split subtorus of $Z(\bL^{*})^{\circ}$, Theorem~\ref{fact:cuspidal} implies that $A_{\bL^{*}}$ is also the maximal $\mathbb{F}_{q}$-split subtorus of $Z((\bH_L)^{\circ})$. Moreover, we have $C_{\bG^{*}}(A_{\bL^{*}}) = \bL^{*}$ and $C_{\bH}(A_{\bL^{*}}) = \bH_{\bL}$. The result follows from these observations and the definition of the relevant stabilisers.
\end{proof}


\section{Jordan decomposition, relative Weyl groups, and endomorphism algebras}

In this section we compare the action of relative Weyl groups on cuspidal
characters with the corresponding action on the Jordan-decomposition side.
We first treat the case where the Levi subgroup has connected centre, and
then reduce the general case to this one by means of a regular embedding.
The final consequence is that the endomorphism algebra of Harish--Chandra
induction depends only on the relevant fibre of the disconnected Jordan
decomposition.

\subsection{The connected-centre case}

Let $\bG$ be a connected reductive group defined over $\fq$, and let
$\bL$ be an $F$--stable Levi factor of an $F$--stable parabolic subgroup of
$\bG$.  Throughout this subsection we assume that $Z(\bL)$ is connected.
Let $A_{\bL}$ denote the maximal $\fq$--split torus in $Z(\bL)$.

Let
\[
\delta:W_{\bG}(A_{\bL})\xrightarrow{\sim}W_{\bG^*}(A_{\bL^*})
\]
be the duality isomorphism.  For $w\in W_{\bG}(A_{\bL})$, write
$w^*=\delta(w)$.  We choose representatives inducing $F$--equivariant
automorphisms
\[
\sigma_w:\bL\longrightarrow\bL,
\qquad
\sigma_{w^*}:\bL^*\longrightarrow\bL^*,
\]
dual to one another.  The induced action on characters of $L$ is given by
\[
\rho\longmapsto \rho\circ\sigma_w^{-1}.
\]

Let $s\in L^*$ be semisimple and put
\[
s^w:=\sigma_{w^*}(s),\qquad
\bH_{L,s}:=C_{\bL^*}(s),\qquad
\bH_{L,s^w}:=C_{\bL^*}(s^w).
\]
Then $\sigma_{w^*}$ identifies $\bH_{L,s}$ with $\bH_{L,s^w}$.  By the
functoriality of Deligne--Lusztig induction under dual automorphisms, the map
\[
f_w^s:\mathcal E(L,s)\longrightarrow \mathcal E(L,s^w),
\qquad
\rho\longmapsto \rho\circ\sigma_w^{-1},
\]
is a bijection.  Similarly, $\sigma_{w^*}$ induces a bijection
\[
f_{w^*}^s:\Uch(H_{L,s})\longrightarrow \Uch(H_{L,s^w}),
\qquad
u\longmapsto u\circ\sigma_{w^*}^{-1}.
\]

Let
\[
J_s^L:\mathcal E(L,s)\longrightarrow \Uch(H_{L,s})
\]
denote the Digne--Michel Jordan decomposition for $\bL$, as normalized in
Theorem~\ref{thm:JD-connected-center}.  Transporting the Jordan
decomposition at $s^w$ back to the Lusztig series of $s$ gives another
bijection
\[
\widetilde J_{s,w}^L
:=
(f_{w^*}^s)^{-1}\circ J_{s^w}^L\circ f_w^s
:
\mathcal E(L,s)\longrightarrow \Uch(H_{L,s}).
\]
The following proposition says that the Digne--Michel normalization is
invariant under this transport.
\begin{proposition}
\label{prop:compatibility-JDs-for-Levi-with-connected-center}
Let $\bG$ be a connected reductive group defined over $\fq$, and let
$\bL$ be an $F$--stable Levi factor of an $F$--stable parabolic subgroup of
$\bG$ such that $Z(\bL)$ is connected.  Let
$w\in W_{\bG}(A_{\bL})$, and let
$w^*\in W_{\bG^*}(A_{\bL^*})$ be the corresponding element under the duality
isomorphism.  Choose representatives inducing $F$--equivariant automorphisms
\[
\sigma_w:\bL\longrightarrow \bL,\qquad
\sigma_{w^*}:\bL^*\longrightarrow \bL^* .
\]
For a semisimple element $s\in L^*$, put
\[
s^w:=\sigma_{w^*}(s),\qquad
\bH_s:=C_{\bL^*}(s),\qquad
\bH_{s^w}:=C_{\bL^*}(s^w).
\]
Then $\sigma_{w^*}$ identifies $\bH_s$ with $\bH_{s^w}$, and for every
$\tau\in \mathcal{E}(L,s)$ one has
\[
J_{s^w}^{L}\bigl(\tau\circ\sigma_w^{-1}\bigr)
=
J_s^L(\tau)\circ\sigma_{w^*}^{-1}.
\]
Equivalently, if
\[
f_w:\mathcal{E}(L,s)\longrightarrow \mathcal{E}(L,s^w),
\qquad
\tau\longmapsto \tau\circ\sigma_w^{-1},
\]
and
\[
f_{w^*}:\Uch(H_s)\longrightarrow \Uch(H_{s^w}),
\qquad
u\longmapsto u\circ\sigma_{w^*}^{-1},
\]
then
\[
f_{w^*}\circ J_s^L
=
J_{s^w}^L\circ f_w,
\]
or, equivalently,
\[
J_s^L
=
f_{w^*}^{-1}\circ J_{s^w}^{L}\circ f_w .
\]
\end{proposition}

\begin{proof}
Changing the chosen representatives of $w$ and $w^*$ changes
$\sigma_w$ and $\sigma_{w^*}$ only by inner automorphisms of $\bL$ and
$\bL^*$, respectively.  Hence the induced maps on irreducible characters and
on unipotent characters are independent of these choices, after the usual
identification of characters under inner conjugacy.

The automorphisms $\sigma_w$ and $\sigma_{w^*}$ are dual to one another.
Therefore Deligne--Lusztig induction is transported by these automorphisms:
for every $F^*$--stable maximal torus $\bT^*\leq \bH_s$,
\[
R_{\sigma_{w^*}(\bT^*)}^{\bL}(s^w)
=
R_{\bT^*}^{\bL}(s)\circ\sigma_w^{-1},
\]
and similarly
\[
R_{\sigma_{w^*}(\bT^*)}^{\bH_{s^w}}(1)
=
R_{\bT^*}^{\bH_s}(1)\circ\sigma_{w^*}^{-1}.
\]
It follows at once that $f_w$ maps $\mathcal{E}(L,s)$ bijectively onto
$\mathcal{E}(L,s^w)$, and that $f_{w^*}$ maps $\Uch(H_s)$ bijectively onto
$\Uch(H_{s^w})$.

Define
\[
\widetilde J_s^L
:=
f_{w^*}^{-1}\circ J_{s^w}^{L}\circ f_w
:
\mathcal{E}(L,s)\longrightarrow \Uch(H_s).
\]
We claim that $\widetilde J_s^L$ satisfies the same Digne--Michel
normalizing conditions as $J_s^L$.  For the first condition, let
$\rho\in\mathcal{E}(L,s)$ and let $\bT^*\leq \bH_s$ be an $F^*$--stable
maximal torus.  Using the preceding functoriality and then the defining
property of $J_{s^w}^L$, we obtain
\begin{align*}
\bigl\langle R_{\bT^*}^{\bL}(s),\rho\bigr\rangle_L
&=
\bigl\langle
R_{\sigma_{w^*}(\bT^*)}^{\bL}(s^w),
\rho\circ\sigma_w^{-1}
\bigr\rangle_L                                      \\
&=
\epsilon_{\bL}\epsilon_{\bH_{s^w}}
\bigl\langle
R_{\sigma_{w^*}(\bT^*)}^{\bH_{s^w}}(1),
J_{s^w}^L(\rho\circ\sigma_w^{-1})
\bigr\rangle_{H_{s^w}}                              \\
&=
\epsilon_{\bL}\epsilon_{\bH_s}
\bigl\langle
R_{\bT^*}^{\bH_s}(1),
\widetilde J_s^L(\rho)
\bigr\rangle_{H_s}.
\end{align*}
Here $\epsilon_{\bH_s}=\epsilon_{\bH_{s^w}}$, since
$\sigma_{w^*}$ identifies $\bH_s$ with $\bH_{s^w}$.

The remaining Digne--Michel conditions are preserved by the same transport
of structure.  The condition for $s=1$ is preserved because the relevant
Deligne--Lusztig varieties, Frobenius actions, and Iwahori--Hecke algebras
are identified by $\sigma_w$ and $\sigma_{w^*}$.  The compatibility with
central characters follows from the fact that, if $z\in Z(L^*)$, then the
linear character $\widehat z\circ\sigma_w^{-1}$ corresponds to
$\widehat{\sigma_{w^*}(z)}$.  Compatibility with Harish--Chandra induction
is transported by the identities
\[
f_w\circ R_{\bM}^{\bL}
=
R_{\sigma_w(\bM)}^{\bL}\circ f_w
\]
and the corresponding identity on the dual side.  The exceptional
$\mathsf E_8$ conditions, the compatibility with central quotients, and the
direct-product condition are likewise invariant under the same
$F$--equivariant isomorphisms.

Thus the family obtained from the Digne--Michel Jordan decomposition by
transport through $(\sigma_w,\sigma_{w^*})$ again satisfies the
Digne--Michel characterizing properties.  By the uniqueness assertion in
Theorem~\ref{thm:JD-connected-center}, it must coincide with the
original family.  Hence
\[
\widetilde J_s^L=J_s^L,
\]
which is precisely
\[
J_s^L
=
f_{w^*}^{-1}\circ J_{s^w}^{L}\circ f_w .
\]
Equivalently,
\[
J_{s^w}^{L}\bigl(\tau\circ\sigma_w^{-1}\bigr)
=
J_s^L(\tau)\circ\sigma_{w^*}^{-1}
\]
for all $\tau\in\mathcal{E}(L,s)$.
\end{proof}

\begin{proposition}\label{prop:weyl-isom-connected-levi-center}
Let $\bG$ be a connected reductive group over $\mathbb F_q$, and let
$\bL$ be an $F$--stable Levi factor of an $F$--stable parabolic subgroup of
$\bG$ such that $Z(\bL)$ is connected.  Let $s\in L^*$ be semisimple, and let
\[
\tau\in \mathcal E(L,s)
\]
be a cuspidal irreducible character of $L$.  Put
\[
\bH:=C_{\bG^*}(s),\qquad
\bH_{\bL}:=C_{\bL^*}(s),
\]
and set
\[
u:=J_s^L(\tau)\in \Uch(H_{\bL}).
\]
Define
\[
W_\tau(\bG)
:=
\left\{
w\in W_{\bG}(A_{\bL})
\;\middle|\;
\tau\circ \sigma_w^{-1}=\tau
\right\},
\]
and
\[
W_u(\bH)
:=
\left\{
w^*\in W_{\bH}(A_{\bL^*})
\;\middle|\;
u\circ \sigma_{w^*}^{-1}=u
\right\},
\]
where
\[
W_{\bH}(A_{\bL^*})
=
N_H(A_{\bL^*})/H_{\bL}.
\]
Then the duality isomorphism
\[
\delta:W_{\bG}(A_{\bL})\xrightarrow{\sim} W_{\bG^*}(A_{\bL^*})
\]
restricts, after the standard identification of the stabilizer of the
$L^*$--class of $s$ with $W_{\bH}(A_{\bL^*})$, to a canonical isomorphism
\[
W_\tau(\bG)
\xrightarrow{\;\sim\;}
W_u(\bH).
\]
Equivalently,
\[
W_\tau(\bG)
\cong
W_{J_s^L(\tau)}(\bH).
\]
\end{proposition}

\begin{proof}
We first recall the identification of the target relative Weyl group.  The
inclusion
\[
N_H(A_{\bL^*})\subset N_{G^*}(A_{\bL^*})
\]
induces an injective homomorphism
\[
W_{\bH}(A_{\bL^*})
=
N_H(A_{\bL^*})/H_{\bL}
\longrightarrow
W_{\bG^*}(A_{\bL^*})
=
N_{G^*}(A_{\bL^*})/L^*.
\]
Its image is precisely the stabilizer of the $L^*$--conjugacy class of $s$.
Indeed, if $n^*\in N_{G^*}(A_{\bL^*})$ represents an element stabilizing the
$L^*$--class of $s$, then there is an element $\ell^*\in L^*$ such that
\[
\ell^* n^* s(n^*)^{-1}(\ell^*)^{-1}=s.
\]
Thus $\ell^* n^*\in N_H(A_{\bL^*})$, and its class modulo $H_{\bL}$ depends
only on the class of $n^*$ modulo $L^*$.

Let $w\in W_\tau(\bG)$.  Then
\[
\tau\circ\sigma_w^{-1}=\tau.
\]
If $w^*=\delta(w)$, then the usual compatibility of Lusztig series with dual
automorphisms gives
\[
\tau\circ\sigma_w^{-1}\in
\mathcal E\bigl(L,\sigma_{w^*}(s)\bigr).
\]
Since the Lusztig series of $L$ are disjoint and
$\tau\in\mathcal E(L,s)$, the elements $s$ and $\sigma_{w^*}(s)$ are
$L^*$--conjugate.  Hence $w^*$ lies in the stabilizer of the $L^*$--class of
$s$, and therefore determines an element of
$W_{\bH}(A_{\bL^*})$.  We denote this element by $\delta_s(w)$.

Choose a representative of $\delta_s(w)$ in $N_H(A_{\bL^*})$.  Then it
centralizes $s$, and hence its action preserves $\bH_{\bL}=C_{\bL^*}(s)$.
By Proposition~\ref{prop:compatibility-JDs-for-Levi-with-connected-center},
applied in the form
\[
J_s^L(\rho\circ\sigma_w^{-1})
=
J_s^L(\rho)\circ\sigma_{\delta_s(w)}^{-1},
\qquad
\rho\in\mathcal E(L,s),
\]
we obtain
\[
u\circ\sigma_{\delta_s(w)}^{-1}
=
J_s^L(\tau)\circ\sigma_{\delta_s(w)}^{-1}
=
J_s^L(\tau\circ\sigma_w^{-1})
=
J_s^L(\tau)
=
u.
\]
Thus $\delta_s(w)\in W_u(\bH)$.

Conversely, let $w^*\in W_u(\bH)$, and let
\[
w:=\delta^{-1}(w^*)
\]
under the inclusion
\[
W_{\bH}(A_{\bL^*})\hookrightarrow W_{\bG^*}(A_{\bL^*}).
\]
Since $w^*$ is represented by an element of $N_H(A_{\bL^*})$, it centralizes
$s$.  Hence $\sigma_{w^*}(s)=s$, and therefore
$\tau\circ\sigma_w^{-1}$ again lies in $\mathcal E(L,s)$.  Using the same
equivariance of the Jordan decomposition, we get
\[
J_s^L(\tau\circ\sigma_w^{-1})
=
J_s^L(\tau)\circ\sigma_{w^*}^{-1}
=
u\circ\sigma_{w^*}^{-1}
=
u
=
J_s^L(\tau).
\]
Since $J_s^L:\mathcal E(L,s)\to \Uch(H_{\bL})$ is a bijection, it follows that
\[
\tau\circ\sigma_w^{-1}=\tau.
\]
Thus $w\in W_\tau(\bG)$.

The two constructions are inverse to each other, and hence $\delta$ induces
the desired canonical isomorphism
\[
W_\tau(\bG)\cong W_u(\bH)
=
W_{J_s^L(\tau)}(\bH).
\]
\end{proof}

\begin{theorem}\label{thm:HC-endo-match-connected-centre}
Let $\bG$ be a connected reductive group over $\fq$ with a connected center.
Let $\bL$ be an $F$-stable Levi factor of an $F$-stable parabolic subgroup of $G$ and
let $\tau\in \cE(L,s)$ be a cuspidal irreducible representation of $L^F$.
Then there is a canonical isomorphism, $$\End_{G}(R_{\bL}^{\bG}(\tau))\cong\End_{H}(R_{\bH_L}^{\bH}(J_s^L(\tau))),$$
where $J_s^L$ is the unique Jordan decomposition as in Theorem \ref{thm:JD-connected-center}.
\end{theorem}
\begin{proof}
The Theorem follows from Proposition \ref{prop:weyl-isom-connected-levi-center} and Theorem \ref{thm:Howlett-Lehrer}.
\end{proof}
\subsection{General case}

We now remove the hypothesis that the centre of the Levi subgroup is connected.
Let $\bL$ be an $F$--stable Levi subgroup of $\bG$.  Choose an $F$--stable
torus $\bS$ containing $Z(\bL)$, and form
\[
\bG'=(\bG\times \bS)/\{(z,z^{-1})\mid z\in Z(\bG)\}.
\]
Let $i:\bG\hookrightarrow \bG'$ be the regular embedding induced by
$g\mapsto (g,1)$, and let $\bS'$ be the image of $\{1\}\times \bS$ in
$\bG'$.  Then $Z(\bG')=\bS'$ is connected.  Put
\[
\bL'=\bL\cdot \bS'.
\]
Then $\bL'$ is an $F$--stable Levi subgroup of $\bG'$, and the restriction
\[
i_L:\bL\hookrightarrow \bL'
\]
is a regular embedding.  We write, as usual,
\[
L=\bL^F,\qquad L'=\bL'^F.
\]

Let
\[
i_L^*:\bL'^*\twoheadrightarrow \bL^*
\]
be the dual epimorphism, and let
\[
\bK=\ker(i_L^*).
\]
For $s\in L^*$ choose $s'\in L'^*$ such that $i_L^*(s')=s$.  Put
\[
\bH_L=C_{\bL^*}(s),\qquad
\bH_{L'}=C_{\bL'^*}(s').
\]
By Lusztig's construction for a regular embedding, we have a surjective map
\[
J_s^L:\mathcal E(L,s)\longrightarrow
\Uch(H_L^\circ)/(H_L/H_L^\circ).
\]
Its fibres are the diagonal-automorphism orbits on $\mathcal E(L,s)$; with
the above regular embedding, these may be described as the $L'/L$--orbits.
For $\tau\in\mathcal E(L,s)$, we denote its $L'/L$--orbit by
\[
[\tau]_{L'}.
\]

We shall use the following equivariance property of the disconnected Jordan
map.

\begin{lemma}\label{lem:disconnected-JD-Weyl-equivariance}
Let $w\in W_{\bG}(A_{\bL})$, and let
\[
w^*=\delta(w)\in W_{\bG^*}(A_{\bL^*})
\]
be its dual element.  Suppose that $w^*$ stabilizes the $L^*$--conjugacy
class of $s$.  Under the standard identification of this stabilizer with
$W_{\bH}(A_{\bL^*})$, where $\bH=C_{\bG^*}(s)$, let the corresponding element
be denoted by $w_s^*$.  Then, for every $\rho\in\mathcal E(L,s)$,
\[
J_s^L(\rho\circ\sigma_w^{-1})
=
w_s^*\cdot J_s^L(\rho)
\]
as elements of
\[
\Uch(H_L^\circ)/(H_L/H_L^\circ).
\]
\end{lemma}

\begin{proof}
The assertion is a formal consequence of Lusztig's construction of
$J_s^L$ by means of the regular embedding $i_L:\bL\hookrightarrow \bL'$,
together with the equivariance of the connected-centre Jordan decomposition.

Indeed, the element $w$ also acts on $\bL'$, since $\bL'=\bL\cdot \bS'$ and
$\bS'$ is central in $\bG'$.  Let $\widetilde w^*$ be the corresponding
dual element acting on $\bL'^*$.  Since $i_L^*(s')=s$ and $w_s^*$ centralizes
$s$, the element $\widetilde w^*$ need not fix $s'$ itself; rather,
\[
\widetilde w^*(s')=s'k_w
\]
for a uniquely determined element $k_w\in K$.  Thus $\widetilde w^*$ permutes
the union
\[
\bigcup_{k\in K}\mathcal E(L',s'k),
\]
and this is precisely the ambiguity accounted for in Lusztig's construction
of the disconnected Jordan map.

Each step in the construction of $J_s^L$ is natural with respect to the
automorphisms induced by $w$ and $\widetilde w^*$: the passage to
$L'/L$--orbits, the correspondence between $L'/L$--orbits in
$\mathcal E(L,s)$ and $K$--orbits in the union
$\bigcup_{k\in K}\mathcal E(L',s'k)$, the passage from $K$ to $K_{s'}$, and
the identification
\[
K_{s'}\cong H_L/H_L^\circ.
\]
The only non-formal part is the connected-centre Jordan decomposition
\[
J_{s'}^{L'}:\mathcal E(L',s')\longrightarrow \Uch(H_{L'}),
\]
and its equivariance is exactly Proposition
\ref{prop:compatibility-JDs-for-Levi-with-connected-center}, applied to the
Levi subgroup $\bL'$ of $\bG'$.

Therefore the disconnected map $J_s^L$ is equivariant for the induced Weyl
actions, giving
\[
J_s^L(\rho\circ\sigma_w^{-1})
=
w_s^*\cdot J_s^L(\rho).
\]
\end{proof}

\begin{theorem}\label{thm:weyl-isom-general-levi}
Let $\bG$ be a connected reductive group over $\mathbb F_q$, and let
$\bL$ be an $F$--stable Levi factor of an $F$--stable parabolic subgroup of
$\bG$.  Let $s\in L^*$ be semisimple, and let
\[
\tau\in\mathcal E(L,s)
\]
be cuspidal.  Put
\[
\bH=C_{\bG^*}(s),\qquad
\bH_L=C_{\bL^*}(s),
\]
and let
\[
\bar u:=J_s^L(\tau)\in \Uch(H_L^\circ)/(H_L/H_L^\circ).
\]
Define
\[
\Stab_{W_{\bG}(A_{\bL})}([\tau]_{L'})
=
\left\{
w\in W_{\bG}(A_{\bL})
\;\middle|\;
[\tau\circ\sigma_w^{-1}]_{L'}=[\tau]_{L'}
\right\},
\]
and
\[
W_{\bar u}(\bH)
=
\left\{
w^*\in W_{\bH}(A_{\bL^*})
\;\middle|\;
w^*\cdot\bar u=\bar u
\right\}.
\]
Then the duality isomorphism
\[
\delta:W_{\bG}(A_{\bL})\xrightarrow{\sim}W_{\bG^*}(A_{\bL^*})
\]
restricts to a canonical isomorphism
\[
\Stab_{W_{\bG}(A_{\bL})}([\tau]_{L'})
\;\xrightarrow{\;\sim\;}\;
W_{\bar u}(\bH).
\]
Equivalently,
\[
\Stab_{W_{\bG}(A_{\bL})}([\tau]_{L'})
\cong
W_{J_s^L(\tau)}(\bH),
\]
where the right-hand side is the stabilizer of the
$H_L/H_L^\circ$--orbit \(J_s^L(\tau)\).
\end{theorem}

\begin{proof}
We first recall the relevant identification on the dual side.  The natural
map
\[
W_{\bH}(A_{\bL^*})
=
N_H(A_{\bL^*})/H_L
\longrightarrow
W_{\bG^*}(A_{\bL^*})
\]
is injective, and its image is the stabilizer of the $L^*$--conjugacy class
of $s$.  Indeed, if $n^*\in N_{G^*}(A_{\bL^*})$ represents an element
stabilizing the $L^*$--class of $s$, then for some $\ell^*\in L^*$ one has
\[
\ell^* n^* s(n^*)^{-1}(\ell^*)^{-1}=s.
\]
Thus $\ell^*n^*\in N_H(A_{\bL^*})$, and its class modulo $H_L$ depends only
on the class of $n^*$ modulo $L^*$.

Let
\[
w\in \Stab_{W_{\bG}(A_{\bL})}([\tau]_{L'}).
\]
Then
\[
[\tau\circ\sigma_w^{-1}]_{L'}=[\tau]_{L'}.
\]
In particular, $\tau\circ\sigma_w^{-1}$ belongs to the same disconnected
Lusztig series $\mathcal E(L,s)$.  On the other hand, the usual compatibility
of Lusztig series with dual automorphisms gives
\[
\tau\circ\sigma_w^{-1}
\in
\mathcal E(L,\delta(w)(s)).
\]
Since Lusztig series are disjoint, $\delta(w)(s)$ is $L^*$--conjugate to
$s$.  Hence $\delta(w)$ lies in the stabilizer of the $L^*$--class of $s$,
and therefore determines an element
\[
\delta_s(w)\in W_{\bH}(A_{\bL^*}).
\]

By Lemma~\ref{lem:disconnected-JD-Weyl-equivariance},
\[
J_s^L(\tau\circ\sigma_w^{-1})
=
\delta_s(w)\cdot J_s^L(\tau).
\]
But $J_s^L$ is constant on $L'/L$--orbits, and
\[
[\tau\circ\sigma_w^{-1}]_{L'}=[\tau]_{L'}.
\]
Therefore
\[
\delta_s(w)\cdot J_s^L(\tau)=J_s^L(\tau).
\]
Thus
\[
\delta_s(w)\in W_{\bar u}(\bH).
\]

Conversely, let
\[
w^*\in W_{\bar u}(\bH),
\]
and let $w\in W_{\bG}(A_{\bL})$ be the inverse image of $w^*$ under the
duality isomorphism, using the inclusion
\[
W_{\bH}(A_{\bL^*})
\hookrightarrow
W_{\bG^*}(A_{\bL^*}).
\]
Since $w^*$ is represented by an element of $N_H(A_{\bL^*})$, it stabilizes
$s$.  Hence
\[
\tau\circ\sigma_w^{-1}\in\mathcal E(L,s).
\]
Again by Lemma~\ref{lem:disconnected-JD-Weyl-equivariance},
\[
J_s^L(\tau\circ\sigma_w^{-1})
=
w^*\cdot J_s^L(\tau).
\]
Since $w^*$ fixes $\bar u=J_s^L(\tau)$, this gives
\[
J_s^L(\tau\circ\sigma_w^{-1})=J_s^L(\tau).
\]
The fibres of $J_s^L$ are precisely the diagonal-automorphism orbits, which
in the present regular embedding are the $L'/L$--orbits.  Therefore
\[
[\tau\circ\sigma_w^{-1}]_{L'}=[\tau]_{L'}.
\]
Thus
\[
w\in \Stab_{W_{\bG}(A_{\bL})}([\tau]_{L'}).
\]

The two constructions are inverse to one another.  Hence duality induces the
desired canonical isomorphism
\[
\Stab_{W_{\bG}(A_{\bL})}([\tau]_{L'})
\cong
W_{J_s^L(\tau)}(\bH).
\]
\end{proof}
The preceding theorem also shows that the endomorphism algebra of the
Harish--Chandra induction depends only on the fibre of the disconnected Jordan
decomposition, up to non-canonical algebra isomorphism.

\begin{corollary}\label{cor:end-algebras-same-JD-fiber}
Let $\bG$ be a connected reductive group over $\mathbb F_q$, and let
$\bL$ be an $F$--stable Levi factor of an $F$--stable parabolic subgroup of
$\bG$.  Let $s\in L^*$ be semisimple, and let
\[
\tau,\tau'\in \mathcal E(L,s)
\]
be irreducible cuspidal characters.  Suppose that
\[
J_s^L(\tau)=J_s^L(\tau')
\]
in
\[
\Uch(H_L^\circ)/(H_L/H_L^\circ),
\qquad
H_L=C_{\bL^*}(s).
\]
Then
\[
\End_G\bigl(R_{\bL}^{\bG}(\tau)\bigr)
\cong
\End_G\bigl(R_{\bL}^{\bG}(\tau')\bigr).
\]
\end{corollary}

\begin{proof}
Choose a regular embedding
\[
i:\bG\hookrightarrow \bG'
\]
with connected centre, and put $\bL'=\bL Z(\bG')$.  By Lusztig's construction of
the disconnected Jordan decomposition, the fibres of
\[
J_s^L:\mathcal E(L,s)\longrightarrow
\Uch(H_L^\circ)/(H_L/H_L^\circ)
\]
are precisely the diagonal-automorphism orbits, equivalently the $L'/L$--orbits
on $\mathcal E(L,s)$.  Hence the equality
\[
J_s^L(\tau)=J_s^L(\tau')
\]
implies that there exists $x\in L'$ such that
\[
\tau'\simeq \tau\circ \sigma_x^{-1},
\]
where $\sigma_x$ denotes the automorphism of $L$ induced by conjugation by $x$.

Let $\bP$ be an $F$--stable parabolic subgroup of $\bG$ with Levi factor $\bL$.
Since $\bL'$ is a Levi factor of the corresponding parabolic subgroup of $\bG'$
and $\bG$ is normalized by $\bG'$, the automorphism $\sigma_x$ preserves the pair
$(\bG,\bL)$ and transports Harish--Chandra induction.  Thus
\[
R_{\bL}^{\bG}(\tau')
\simeq
R_{\bL}^{\bG}(\tau\circ\sigma_x^{-1})
\simeq
R_{\bL}^{\bG}(\tau)\circ\sigma_x^{-1}.
\]
Twisting by the automorphism $\sigma_x$ is an equivalence of the category of
$G$--modules, and it preserves endomorphism algebras.  Therefore
\[
\End_G\bigl(R_{\bL}^{\bG}(\tau')\bigr)
\cong
\End_G\bigl(R_{\bL}^{\bG}(\tau)\circ\sigma_x^{-1}\bigr)
\cong
\End_G\bigl(R_{\bL}^{\bG}(\tau)\bigr).
\]
This proves the claim.
\end{proof}

\section{Harish--Chandra compatibility in the connected-centre case}
\label{sec:HC-compatible-JD-connected-centre}

In this section we record the Harish--Chandra compatibility of the
Digne--Michel Jordan decomposition in the connected-centre case.

Let $\bG$ be a connected reductive group over $\fq$ with connected centre, and
let $s\in G^*$ be semisimple.  We fix, once and for all, the unique
Jordan decomposition of Digne--Michel,
\[
J_s^G:\cE(G,s)\longrightarrow \Uch(H),
\qquad
H=C_{\bG^*}(s)^{F^*},
\]
as in Theorem~\ref{thm:JD-connected-center}.  If $\bL$ is an
$F$-stable Levi factor of an $F$-stable parabolic subgroup of $\bG$, then
$Z(\bL)$ is again connected.  Hence the same theorem applies to $\bL$ and gives
\[
J_s^L:\cE(L,s)\longrightarrow \Uch(H_L),
\qquad
H_L=C_{\bL^*}(s)^{F^*}.
\]

Let $\tau\in \cE(L,s)$ be cuspidal and put
\[
u_\tau:=J_s^L(\tau).
\]
By Theorem~\ref{thm:HC-endo-match-connected-centre}, there is a canonical
isomorphism of relative endomorphism algebras
\[
\End_G\!\bigl(R_{\bL}^{\bG}(\tau)\bigr)
\cong
\End_H\!\bigl(R_{\bH_L}^{\bH}(u_\tau)\bigr).
\]
Consequently, this isomorphism induces a canonical bijection between the
corresponding Harish--Chandra series,
\[
\J_{L,\tau}^{G}:
\Irr(G,(L,\tau))
\xrightarrow{\sim}
\Irr(H,(H_L,u_\tau)).
\]
Equivalently,
\[
\J_{L,\tau}^{G}
\]
is the bijection obtained by transporting the simple modules of the relative
endomorphism algebra on the $G$-side to those of the corresponding relative
endomorphism algebra on the centralizer side.
Using Proposition~\ref{prop:Lusztig-series-union-HC}, the Lusztig series
$\cE(G,s)$ is the disjoint union of the Harish--Chandra series
$\Irr(G,(L,\tau))$ with $\tau\in \cE(L,s)$ cuspidal, taken up to
$G$-conjugacy of cuspidal pairs.  We therefore define
\[
\J_{G,s}:\cE(G,s)\longrightarrow \Uch(H)
\]
by requiring
\[
\J_{G,s}|_{\Irr(G,(L,\tau))}=\J_{L,\tau}^{G}.
\]
From now on, $u_{\tau}=J_s^L(\tau)$ where $J_s^L$ is the unique Jordan decomposition as in Theorem \ref{thm:JD-connected-center} and $\J_{G,s}$ is constructed as above using this collection of bijections at the cupidal levels.
\smallskip
\begin{lemma}[The Harish--Chandra series map is a Jordan decomposition]
\label{lem:HC-series-map-is-JD}
The map
\[
\J_{G,s}:\cE(G,s)\longrightarrow \Uch(H)
\]
constructed above is a Jordan decomposition of $\cE(G,s)$.

More precisely, for every $\rho\in\cE(G,s)$ one has
\[
\J_{G,s}(\rho)=J_s^G(\rho)
\]
unless the two characters lie in one of the exceptional pairs occurring in
Proposition~\ref{prop:unique-HC-char}.  In those exceptional cases the two
possible characters have the same degree and the same multiplicities in the
relevant Deligne--Lusztig virtual characters.  Hence the possible interchange
inside such a pair does not affect the defining character-theoretic properties
of Jordan decomposition.
\end{lemma}

\begin{proof}
We argue by induction on the semisimple rank of $\bG$.  The assertion is clear
for tori.  By the usual reduction to simple adjoint factors, it is enough to
consider the case where $\bG/Z(\bG)$ is simple.

Let
\[
\rho\in\Irr(G,(L,\tau))\subseteq\cE(G,s),
\]
and put
\[
u_\tau=J_s^L(\tau),\qquad
\eta=J_s^G(\rho),\qquad
\eta'=\J_{G,s}(\rho).
\]
By the Harish--Chandra compatibility of the Digne--Michel Jordan decomposition
\cite[Theorem~4.7.5]{Book:GeckandMalle}, if
\[
\rho\in \Irr(G,(L,\tau))\subseteq \cE(G,s),
\]
then
\[
J_s^G(\rho)\in \Irr(H,(H_L,J_s^L(\tau))).
\]
Moreover, for every $F$-stable Levi subgroup $\bM$ with
$\bL\leq \bM\leq \bG$, one has
\[
{}^*R_{\bH_M}^{\bH}\bigl(J_s^G(\rho)\bigr)
=
J_s^M\!\left({}^*R_{\bM}^{\bG}(\rho)\right).
\]
By the induction hypothesis applied to $\bM$, the right-hand side agrees with
\[
\J_{M,s}\!\left({}^*R_{\bM}^{\bG}(\rho)\right),
\]
up to the exceptional ambiguity described in
Proposition~\ref{prop:unique-HC-char}.  On the other hand, the construction of
$\J_{G,s}$ and the comparison theorem give
\[
\J_{M,s}\!\left({}^*R_{\bM}^{\bG}(\rho)\right)
=
{}^*R_{\bH_M}^{\bH}\bigl(\J_{G,s}(\rho)\bigr).
\]
By the induction hypothesis applied to $\bM$, the maps $J_s^M$ and $\J_{M,s}$
agree on the constituents of ${}^*R_{\bM}^{\bG}(\rho)$, except possibly for the
exceptional pairs of Proposition~\ref{prop:unique-HC-char}.  In the exceptional
cases, the two possible constituents have the same proper Harish--Chandra
restrictions and the same Deligne--Lusztig multiplicities, so the ambiguity is
invisible to the character-theoretic tests used below.  Thus $J_s^G(\rho)$ and
$\J_{G,s}(\rho)$ have the same Harish--Chandra restrictions to all proper
Levi subgroups of $\bH$, up to precisely the exceptional ambiguity already
listed in Proposition~\ref{prop:unique-HC-char}.

It remains to compare degrees.  The degree formula for Jordan decomposition \cite[Theorem 3.2.18]{Book:GeckandMalle} and \cite[Corollary 2.6.6]{Book:GeckandMalle},
together with the degree formula for the Howlett--Lehrer parametrization and
the parameter-preserving endomorphism algebra comparison in
Theorem~\ref{thm:HC-endo-match-connected-centre}, gives
\[
J_s^G(\rho)(1)=\J_{G,s}(\rho)(1).
\]
Applying Proposition~\ref{prop:unique-HC-char} to the connected reductive
group $\bH$ and the cuspidal pair $(H_L,u_\tau)$, we obtain
\[
J_s^G(\rho)=\J_{G,s}(\rho)
\]
unless the two characters lie in one of the exceptional pairs described in
that proposition.

In the remaining cases, the two possible characters have the same degree
\(512\) in type $\mathsf E_7$ and \(4096\) in type $\mathsf E_8$, and the two
members of each pair have the same multiplicities in the relevant
Deligne--Lusztig virtual characters.  Therefore replacing one by the other
does not change the defining scalar-product identities for Jordan
decomposition.  Hence $\J_{G,s}$ is a Jordan decomposition.
\end{proof}

By Lemma~\ref{lem:HC-series-map-is-JD}, the map
\[
\J_{G,s}:\cE(G,s)\longrightarrow \Uch(H)
\]
constructed above from the Harish--Chandra series bijections is a Jordan
decomposition.  We shall use this Jordan decomposition in the following
compatibility statement.
\begin{proposition}\label{prop:HC-induction-connected-centre}
Let $\bG$ be a connected reductive group over $\fq$ with a connected centre.
Let $s\in G^{*}$ be any semisimple element and $\bL^*\leq \bG^*$  be an  $F^*$-stable  Levi factor of an $F^*$-stable parabolic subgroup of $\bG^*$ containing $s$ with dual $\bL$.
Let $\J_{G,s}:\cE(G,s)\rightarrow\Uch(H)$ be the Jordan decomposition as above. Then the following diagram commutes:
\[
\begin{tikzcd}
\mathbb{Z}\cE(G,s)\arrow[r,"\J_{G,s}"]
&\mathbb{Z}\Uch(H)\\
\mathbb{Z}\cE(M,s)\arrow[r,"\J_{M,s}"]\arrow[u,"R_{\bM}^{\bG}"]
& \mathbb{Z}\Uch(H_M)\arrow[u,"R_{\bH_M}^{\bH}"],
\end{tikzcd}
\]
where $\bL^*\leq \bM^*\leq \bG^*$ are $F^*$-stable Levi subgroups of $F^*$-stable parabolic subgroups of $\bG^*$.
\end{proposition}
\begin{proof}
The proposition follows from Theorem \ref{thm:Howlett-Lehrer-comparison} and the definition of bijections $\J_{G,s}$ and $\J_{M,s}$.
\end{proof}

\section{Preferred extensions of cuspidal unipotent characters}
\label{sec:preferred-extensions-cuspidal-unipotents}

\subsection{Preferred extensions in the quasisimple case}
Let $\mathbf H^\circ$ be a connected reductive group over $\overline{\mathbb F}_p$
equipped with a Frobenius endomorphism $F$, and put $H^\circ:=\mathbf H^{\circ F}$.
Let $\mathbf H$ be an $F$--stable (possibly disconnected) algebraic group with identity component
$\mathbf H^\circ$ such that $\pi_0(\mathbf H):=\mathbf H/\mathbf H^\circ$ is abelian, and set
$H:=\mathbf H^F$.
Let $u^\circ\in\Irr(H^\circ)$ be a cuspidal unipotent character, and write
\[
I \ :=\ I_H(u^\circ)\ =\ \{\,h\in H:\ {}^hu^\circ\simeq u^\circ\,\}.
\]

We shall use the word \emph{canonical} in this section in the following
pinned-normalized sense.  For a fixed pinning, the component action on the based
root datum is fixed.  If an invariant cuspidal unipotent character admits
several extensions to its inertia group, those extensions differ by the usual
characters of the component quotient.  Lusztig's preferred-extension
normalization, recalled below, selects one of them without any further auxiliary
choice.  Thus the extensions constructed here are canonical relative to the
fixed pinning and to Lusztig's preferred-extension convention; this is the sense
in which they are used later in the Jordan decomposition.

The semidirect-product notation in this section is used only when a pinned
semidirect component model, in the sense of
Definition~\ref{def:pinned-semidirect-component-model}, has been specified.  In
the centralizer situations needed later, this model is supplied by
Lemma~\ref{lem:DM-semisimple-centralizer-splits}; no assertion is being made
that an arbitrary disconnected reductive group splits over its identity
component.

\begin{lemma}[Existence of preferred extensions on cyclic subgroups]
\label{lem:cyclic-preferred-extension}
Let $\mathbf H^\circ$ be a quasisimple adjoint algebraic group with Frobenius $F$ and put
$H^\circ=\mathbf H^{\circ F}$. Let $\mathbf H$ be an $F$--stable (possibly disconnected) algebraic
group with identity component $\mathbf H^\circ$ and set $H=\mathbf H^F$.
Let $u^\circ\in\Irr(H^\circ)$ be a cuspidal unipotent character, and let
\[
I \ :=\ I_H(u^\circ)=\{\,h\in H:\ {}^h u^\circ\simeq u^\circ\,\}.
\]
Fix a cyclic subgroup $C\le I/H^\circ$, and let $H_C\le H$ be its inverse image.
Assume that $H_C$ has a pinned semidirect component model
\[
        H_C=H^\circ\rtimes A_C,
        \qquad A_C\cong C,
\]
in the sense of Definition~\ref{def:pinned-semidirect-component-model}.

Then there exists an extension $u_C\in\Irr(H_C)$ of $u^\circ$ to $H_C$ which is
canonical once one fixes Lusztig's preferred extension convention for Weyl-group representations.
\end{lemma}

\begin{proof}
Choose an element $\sigma\in A_C$ whose image generates $C$.  By the pinned
semidirect component model,
\[
        H_C=H^\circ\rtimes\langle\sigma\rangle .
\]
This is the only point in the proof where a semidirect-product decomposition is
used; for semisimple centralizers it is supplied by
Lemma~\ref{lem:DM-semisimple-centralizer-splits}.  Since $C\le I/H^\circ$, the
character $u^\circ$ is $H_C$--stable, hence extends to $H_C$ (cyclic-quotient
extension criterion, e.g.\ \cite[Cor.\ 11.22]{Isaacs}); different extensions correspond to multiplying the
intertwiner $T_\sigma$ by a $|C|$-th root of unity.

\textbf{Lusztig's preferred extension on the Weyl-group side.}
Let $W$ be the Weyl group of $\mathbf H^\circ$.
The automorphism $\sigma$ induces an automorphism (still denoted) $\sigma$ of the based root datum,
hence of $W$, and we form the semidirect product $C_cW:=W\rtimes\langle\sigma\rangle$ where
$c=|\sigma|$.
Lusztig shows that if an irreducible $W$--module $E$ is extendable to $C_cW$, then it has $c$
extensions, and he \emph{singles out one} such extension and calls it the \emph{preferred extension},
defined case-by-case and functorially with respect to products.  \cite[\S17.2]{Lusztig1986character}
More generally, for an abelian group acting on $W$, Lusztig constructs canonical intertwiners
$\sigma:E\to E$ that realize the preferred extensions and build an extension of $E$ to the relevant
semidirect product. \cite[\S17.3]{Lusztig1986character}

\smallskip\noindent
\emph{Transfer from the Weyl-group side.}
Let $c=|\sigma|$.
Fix an irreducible representation $(\rho,V)$ of $H^\circ$ affording $u^\circ$.
Since $u^\circ$ is $H_C$--stable, there exists an intertwiner
$T_\sigma:V\to V$ with
\[
T_\sigma\,\rho(h)\,T_\sigma^{-1} \;=\; \rho(\sigma(h))
\qquad(h\in H^\circ),
\]
unique up to scalar by Schur's lemma. Choosing $T_\sigma$ with $T_\sigma^c=1$
is equivalent to choosing an extension of $u^\circ$ to $H_C$, and different
extensions correspond to multiplying $T_\sigma$ by a $c$-th root of unity.

On the geometric side, $u^\circ$ lies in a (generalized) Springer block $I$
and corresponds to a $\sigma$--stable Weyl-group character $\varphi$
(of a suitable relative Weyl group $W$, in particular $\varphi$ is extendable
to $W\rtimes\langle\sigma\rangle$).  By \cite[\S2]{DLM03} (following Lusztig's
construction of characteristic functions from $F$--stable pairs, cf.
\cite[24.1--24.2]{Lus84}), an extension $\tilde\varphi$ of $\varphi$ to
$W\rtimes\langle\sigma\rangle$ canonically determines the corresponding
$\sigma$--equivariant structure on the local system in the block, hence fixes
the scalar by which $\sigma$ acts on the resulting function/representation.
Equivalently, it fixes the scalar ambiguity in the intertwiner $T_\sigma$
above, and therefore picks out a unique extension of $u^\circ$ to $H_C$.

We now take $\tilde\varphi$ to be Lusztig's preferred extension
\cite[\S17.2--\S17.3]{Lusztig1986character}.  The resulting intertwiner $T_\sigma$ (hence the
corresponding extension) is then canonical once the preferred-extension
convention is fixed; denote this extension by $u_C\in\Irr(H_C)$.

\end{proof}
\begin{lemma}[Gluing preferred extensions in the non-cyclic abelian case]
\label{lem:gluing-preferred-extension-non-cyclic}
Let \(u^\circ\in \Irr(H^\circ)\) be \(I\)-invariant, and put
\[
\Omega \ :=\ I/H^\circ .
\]
Assume that \(\Omega\) is non-cyclic abelian and that a decomposition
\[
\Omega \ \cong\ \Gamma \times \Phi
\]
has been fixed with \(\Gamma\) and \(\Phi\) cyclic. Let
\[
I_\Gamma,\ I_\Phi \ \leq\ I
\]
be the inverse images of \(\Gamma\) and \(\Phi\) under the quotient map
\(I\twoheadrightarrow \Omega\).

Assume that:

\begin{enumerate}
\item there is an extension \(u_\Gamma\in \Irr(I_\Gamma)\) of \(u^\circ\) which is
\(\Phi\)-invariant (equivalently, \(I\)-invariant, since \(I/I_\Gamma\simeq \Phi\));

\item there is an extension \(u_\Phi\in \Irr(I_\Phi)\) of \(u^\circ\).
\end{enumerate}

Then there exists a unique extension \(\widetilde u\in \Irr(I)\) of \(u^\circ\) such that
\[
\widetilde u|_{I_\Gamma}=u_\Gamma,
\qquad
\widetilde u|_{I_\Phi}=u_\Phi.
\]
In particular, if \(u_\Gamma\) and \(u_\Phi\) are canonically determined, then so is
\(\widetilde u\).
\end{lemma}

\begin{proof}
Since \(I_\Gamma\triangleleft I\) and \(I/I_\Gamma\simeq \Phi\) is cyclic, the
\(\Phi\)-invariance of \(u_\Gamma\) implies, by the cyclic-quotient extension criterion
(e.g.\ \cite[Cor.~11.22]{Isaacs}), that \(u_\Gamma\) extends to some
\(\widetilde u_0\in \Irr(I)\).

We claim that \(\widetilde u_0|_{I_\Phi}\) is irreducible. Indeed, let
\(\rho:I\to \GL(V)\) afford \(\widetilde u_0\). Since \(\widetilde u_0\) extends
\(u^\circ\), the restriction \(\rho|_{H^\circ}\) is irreducible. Any \(I_\Phi\)-stable
subspace of \(V\) is in particular \(H^\circ\)-stable, hence must be \(0\) or \(V\).
Thus \(\rho|_{I_\Phi}\) is irreducible, so \(\widetilde u_0|_{I_\Phi}\in \Irr(I_\Phi)\),
and it extends \(u^\circ\).

Now \(I_\Phi/H^\circ\simeq \Phi\) is abelian, and both
\(\widetilde u_0|_{I_\Phi}\) and \(u_\Phi\) lie above \(u^\circ\). By Clifford theory,
there is therefore a unique linear character \(\lambda_\Phi\in \Irr(\Phi)\) such that
\[
\widetilde u_0|_{I_\Phi}
\;=\;
u_\Phi \otimes \widehat{\lambda_\Phi},
\]
where \(\widehat{\lambda_\Phi}\) denotes the inflation of \(\lambda_\Phi\) to \(I_\Phi\).

Let \(\lambda\in \Irr(\Omega)\) be the inflation of \(\lambda_\Phi\) along the projection
\(\Omega=\Gamma\times\Phi\twoheadrightarrow \Phi\), and let \(\widehat\lambda\) be its
inflation to \(I\). Set
\[
\widetilde u \ :=\ \widetilde u_0\otimes \widehat\lambda^{-1}.
\]
Because \(\widehat\lambda\) is trivial on \(I_\Gamma\), we have
\[
\widetilde u|_{I_\Gamma}
=
\widetilde u_0|_{I_\Gamma}
=
u_\Gamma.
\]
Because \(\widehat\lambda|_{I_\Phi}=\widehat{\lambda_\Phi}\), we obtain
\[
\widetilde u|_{I_\Phi}
=
\bigl(\widetilde u_0|_{I_\Phi}\bigr)\otimes \widehat{\lambda_\Phi}^{-1}
=
u_\Phi.
\]
Thus \(\widetilde u\) is the required extension of \(u^\circ\).

For uniqueness, let \(\widetilde u_1,\widetilde u_2\in \Irr(I)\) be two extensions of
\(u^\circ\) with the same restrictions to \(I_\Gamma\) and \(I_\Phi\). Since
\(\Omega=I/H^\circ\) is abelian and both characters lie above \(u^\circ\), Clifford theory
gives
\[
\widetilde u_2 \ =\ \widetilde u_1\otimes \widehat\mu
\]
for some \(\mu\in \Irr(\Omega)\). Restricting to \(I_\Gamma\) shows that
\(\mu|_\Gamma=1\), and restricting to \(I_\Phi\) shows that \(\mu|_\Phi=1\).
Since \(\Omega=\Gamma\times \Phi\), it follows that \(\mu=1\), hence
\(\widetilde u_1=\widetilde u_2\).

Therefore the extension \(\widetilde u\) is unique. The final assertion on canonicity is
immediate.
\end{proof}

\begin{lemma}[\(\Phi\)-stable preferred extension along the graph factor]
\label{lem:preferred-extension-cuspidal-abelian}
Assume that \(\bH^\circ\) is simple adjoint. Assume that
\(\Omega:= H/H^\circ\) is non-cyclic abelian with decomposition
\[
\Omega \ \cong\ \Gamma \times \Phi
\]
where \(\Gamma\) and \(\Phi\) are cyclic.  Assume moreover that this decomposition is realized by a
pinned semidirect component model
\[
        H=H^\circ\rtimes (\Gamma\times\Phi).
\]
Set $K:=H^\circ\rtimes\Gamma\lhd H$.  Then there exists a canonically preferred extension $\tilde u_\Gamma\in\Irr(K)$ of $u^\circ$ which is $\Phi$--invariant.
\end{lemma}

\begin{proof}
For a quasisimple adjoint $\mathbf H^\circ$, the only Lie types in which the relevant
outer action can contribute a non-cyclic abelian inertia quotient are the three families
appearing in Malle's reduction in \cite[Prop.~2.2]{Malle08}, namely:
\[
A_{n-1}\ (n\ge3),\qquad D_n\ (n\ge4),\qquad E_6.
\]
In type $A_{n-1}$ (equivalently $L_n(q)$ with $n\ge3$) there are \emph{no} cuspidal unipotent
characters, so there is nothing to prove.

Thus we may assume $\mathbf H^\circ$ is of type $D_n$ or $E_6$.
By the pinned semidirect component model in the statement, the decomposition
\(\Omega\cong \Gamma\times\Phi\) is represented inside $H$ and we may write
\[
H \ =\ H^\circ \rtimes (\Gamma\times\Phi),
\]
where $\Gamma$ is generated by a (quasi-central) graph automorphism of order $r\in\{2,3\}$,
and $\Phi$ is cyclic and commutes with $\Gamma$ (``field part'' in the sense of Malle).
Set $K:=H^\circ\rtimes\Gamma\lhd H$.
\smallskip\noindent
The argument is identical to the one used by Malle in the proof of \cite[Prop.~2.3]{Malle08}
in the present cuspidal-unipotent situation: in each of the cases
$D_4$ (i.e.\ $O_8^+(q)$), $E_6(q)$, and $D_n(q)$ with $n=s^2$ ($s$ even),
he produces a $\Phi$--stable virtual character $R$ of $K$ (a suitable generalized
Deligne--Lusztig character for the disconnected group)
whose decomposition into irreducibles separates the different extensions of $u^\circ$
to $K$ by \emph{pairwise distinct multiplicities}. More precisely:
\begin{itemize}
\item In type $D_4$, one takes $R=R^{G\sigma}_{T\sigma}(1)$ in the disconnected group
$G\rtimes\langle\sigma\rangle$; its decomposition contains the $r$ extensions of the cuspidal
unipotent character with \emph{different multiplicities}, and $T^F$ is stabilized by field
automorphisms, hence $R$ is $\Phi$--stable \cite[Prop.~2.3]{Malle08}.
\item In type $E_6$, the proof of \cite[Prop.~2.3]{Malle08} similarly uses an explicit
decomposition of a suitable generalized Deligne--Lusztig character in $G\rtimes\langle\gamma\rangle$
showing that the extensions occur with different multiplicities.
\item In type $D_n$ ($n=s^2$ with $s$ even), Malle takes a $\gamma$-- and $F$--stable Levi $L$,
extends a suitable unipotent character $\rho_1$ to $\widetilde L=L\rtimes\langle\gamma\rangle$,
and applies \cite[Cor.~2.4]{DigneMichel94} to obtain a $\Phi$--stable virtual character
$R=R^{\widetilde G}_{\widetilde L}(\widetilde\rho_1)$ in which the two extensions of $u^\circ$
to $K$ occur with distinct multiplicities \cite[Prop.~2.3]{Malle08}.
\end{itemize}

Now write the resulting $\Phi$--stable virtual character as
\[
R \;=\; \sum_{\eta\in\Irr(K)} m_\eta\,\eta.
\]
Since $R$ is $\Phi$--stable, the multiplicity $m_\eta$ is constant on $\Phi$--orbits.
Since the irreducible constituents $\eta$ lying above $u^\circ$ occur with \emph{pairwise distinct}
multiplicities (as stated in \cite[Prop.~2.3]{Malle08}), there is a unique such constituent with
maximal $m_\eta$; denote it by $\tilde u_\Gamma$.  Its $\Phi$--orbit must be trivial (otherwise an
orbit would contain two constituents with the same multiplicity), hence $\tilde u_\Gamma$ is
$\Phi$--invariant.  This is the canonical choice used below.


\end{proof}

\begin{lemma}[Cuspidal unipotent invariance on quasi-simple factors]\label{lem:cuspidal-unipotent-invariant}
Let \(\mathbf X\) be a connected quasi-simple algebraic group over
\(\overline{\mathbb F}_p\), equipped with a Frobenius endomorphism \(F\), and
put \(X=\mathbf X^F\). If $\chi\in\Uch(X)$ is cuspidal unipotent, then $\chi$ is invariant under all automorphisms of $X$.
In particular, for any subgroup $A\le \Aut(X)$ one has $I_A(\chi)=A$.
\end{lemma}

\begin{proof}
As recalled by Malle, in the classical types Lusztig’s parametrisation of unipotent characters by
symbols implies that a classical group has \emph{at most one} cuspidal unipotent character (see
Table~1 in \cite[\S2.1]{Malle2017}); consequently this cuspidal unipotent character is fixed
by \emph{all} automorphisms of $G$.
For groups of exceptional type, see the proof in \cite[Thm.~ 3.2]{Malle2017})  and \cite[Thm.~2.5]{Malle08}.
Together these references yield that any cuspidal unipotent character of $X=\mathbf X^F$ is
$\Aut(X)$--invariant.

\end{proof}

We recall a standard fact from the literature in the form of a Lemma.
\begin{lemma}\label{lem:unipotent-trivial-central-character}
Let $\mathbf G$ be a connected reductive group defined over $\mathbb F_q$ with Frobenius $F$.
If $\chi\in\cE(\mathbf G^F,(1))$ is a unipotent character (in particular, a cuspidal unipotent
character), then $Z(\mathbf G)^F\subseteq\ker(\chi)$. Equivalently, the central character of $\chi$
is trivial.
\end{lemma}

\begin{proof}
Let $\pi:\mathbf G\to \mathbf G_{\ad}:=\mathbf G/Z(\mathbf G)$ be the adjoint quotient.  Then $\pi$
is defined over $\mathbb F_q$, has central kernel $Z(\mathbf G)$, and is surjective, hence
$\pi(\mathbf G)$ contains the derived group $(\mathbf G_{\ad})'$.
Therefore \cite[Prop.~13.20]{Book:DigneandMichel} applies to $\pi$ and shows that every unipotent character of
$\mathbf G^F$ is of the form $\bar\chi\circ \pi^F$ with $\bar\chi$ a unipotent character of
$\mathbf G_{\ad}^F$.
In particular, $\chi$ is inflated from $\mathbf G_{\ad}^F$, so $\ker(\pi^F)=Z(\mathbf G)^F$
acts trivially; hence $Z(\mathbf G)^F\subseteq\ker(\chi)$, i.e.\ the central character is trivial.
\end{proof}
We now combine the preceding lemmas in the form needed later. For the construction of the
canonical Jordan decomposition, it is enough to know that a cuspidal unipotent character admits a
canonically determined extension across an abelian component group.

\begin{proposition}[Canonical extension in the abelian component case]
\label{prop:preferred-extension-abelian-component}
Let $\mathbf H^\circ$ be a quasisimple algebraic group over $\overline{\mathbb F}_p$
equipped with a Frobenius endomorphism $F$, and set $H^\circ:=\mathbf H^{\circ F}$.
Let $\mathbf H$ be an $F$--stable (possibly disconnected) algebraic group with identity component
$\mathbf H^\circ$, and put $H:=\mathbf H^F$. Assume that the component group
\[
\pi_0(\mathbf H)=\mathbf H/\mathbf H^\circ
\]
is abelian, so that $H/H^\circ$ is abelian as well.  Assume moreover that $H$ has a pinned semidirect
component model in the sense of Definition~\ref{def:pinned-semidirect-component-model}.  By
Remark~\ref{rmk:DM-supplies-pinned-component-models}, this extra hypothesis is automatic for the
semisimple centralizers to which the proposition is applied below. Let
\[
u^\circ\in\Irr(H^\circ)
\]
be a cuspidal unipotent character. Then there exists a canonically determined extension
\[
\widetilde u\in\Irr(H)
\]
of $u^\circ$ to $H$.

We shall refer to this canonically determined extension as the \emph{preferred extension} of
$u^\circ$ to $H$.
\end{proposition}
\begin{proof}
Let
\[
\pi:\mathbf H \longrightarrow \overline{\mathbf H}:=\mathbf H/Z(\mathbf H^\circ)
\]
be the quotient by the center of \(\mathbf H^\circ\). Then
\[
\overline{\mathbf H}^{\,\circ}=\mathbf H^\circ/Z(\mathbf H^\circ)
\]
is adjoint, and
\[
\overline{\mathbf H}/\overline{\mathbf H}^{\,\circ}\;\cong\;\mathbf H/\mathbf H^\circ,
\]
so the component group is unchanged.  The chosen pinned semidirect component model for $H$
maps to a pinned semidirect component model for $\overline H$ with the same component quotient.
By Lemma~\ref{lem:unipotent-trivial-central-character},
the cuspidal unipotent character \(u^\circ\) is trivial on \(Z(\mathbf H^\circ)^F\), hence is the
inflation of a cuspidal unipotent character \(\overline u^\circ\in\Irr(\overline H^\circ)\), where
\(\overline H^\circ=\overline{\mathbf H}^{\,\circ F}\). Any extension of \(\overline u^\circ\) to
\(\overline H:=\overline{\mathbf H}^{\,F}\) inflates to an extension of \(u^\circ\) to \(H\).
Therefore, for the existence and canonicity statement, we may replace \(\mathbf H\) by
\(\overline{\mathbf H}\) and hence assume from the outset that \(\mathbf H^\circ\) is simple adjoint.

Set
\[
\Omega \ :=\ H/H^\circ .
\]
By Lemma~\ref{lem:cuspidal-unipotent-invariant}, the character \(u^\circ\) is fixed by all
automorphisms of \(H^\circ\), hence
\[
I_H(u^\circ)=H.
\]

If \(\Omega\) is cyclic, then applying
Lemma~\ref{lem:cyclic-preferred-extension} with \(C=\Omega\) yields a canonical
extension
\[
\widetilde u\in\Irr(H)
\]
of \(u^\circ\). This settles the cyclic case.

Assume now that \(\Omega\) is non-cyclic abelian. By Malle's reduction
(\cite[Prop.~2.2]{Malle08}), the only simple adjoint types for which such a non-cyclic abelian
component quotient can occur in the present setting are
\[
A_{n-1}\ (n\ge 3),\qquad D_n\ (n\ge 4),\qquad E_6.
\]
In type \(A_{n-1}\) there are no cuspidal unipotent characters, so this case does not arise.

Thus \(\mathbf H^\circ\) is of type \(D_n\) or \(E_6\). Again by the same reduction, after fixing
commuting representatives of the component action, we may choose a decomposition
\[
\Omega \ \cong\ \Gamma\times\Phi
\]
with \(\Gamma\) and \(\Phi\) cyclic, where \(\Gamma\) is the graph part and \(\Phi\) the cyclic
field part. Let
\[
I_\Gamma,\ I_\Phi \le H
\]
denote the inverse images of \(\Gamma\) and \(\Phi\) under the quotient map \(H\twoheadrightarrow\Omega\).

By Lemma~\ref{lem:preferred-extension-cuspidal-abelian}, there exists a canonically preferred
extension
\[
u_\Gamma\in\Irr(I_\Gamma)
\]
of \(u^\circ\) which is \(\Phi\)-invariant. By
Lemma~\ref{lem:cyclic-preferred-extension}, applied to the cyclic subgroup
\(\Phi\le \Omega\), there exists a canonical extension
\[
u_\Phi\in\Irr(I_\Phi)
\]
of \(u^\circ\).

Now Lemma~\ref{lem:gluing-preferred-extension-non-cyclic} applies and yields a unique extension
\[
\widetilde u\in\Irr(H)
\]
of \(u^\circ\) such that
\[
\widetilde u|_{I_\Gamma}=u_\Gamma,
\qquad
\widetilde u|_{I_\Phi}=u_\Phi.
\]
Since \(u_\Gamma\) and \(u_\Phi\) are canonical, the character \(\widetilde u\) is canonical as
well.

Hence in all cases \(u^\circ\) admits a preferred extension to \(H\).
\end{proof}


\subsection{Pinned representatives for component actions}
Fix a pinning $\cP$ of $G$ and hence a pinning of $G^*$ and of $H^\circ=C_{G^*}(s)^\circ$.
Write $\cP_{H^\circ}$ for this pinning.  Let $\Aut(H^\circ,\cP_{H^\circ})$ denote automorphisms of
the algebraic group $H^\circ$ preserving the pinning; it is canonically identified with the
automorphism group of the based root datum of $H^\circ$.

\begin{lemma}[Pinned representatives of the component action]\label{lem:pinned-reps-Omega}
Let $H=C_{G^*}(s)$ and $\Omega:=H/H^\circ$.  Conjugation by $H$ induces a homomorphism
$\Omega\to\Out(H^\circ)$.  The pinning $\cP_{H^\circ}$ provides a canonical section
$\Out(H^\circ)\to \Aut(H^\circ,\cP_{H^\circ})$.  Composing, we obtain a canonical homomorphism
\[
\Omega\ \longrightarrow\ \Aut(H^\circ,\cP_{H^\circ}),\qquad \omega\ \longmapsto\ \sigma_\omega,
\]
whose image consists of pinned automorphisms of $H^\circ$ representing the outer action class of $\omega$.
\end{lemma}

\begin{proof}
Self evident.
\end{proof}

\subsection{Permutation factors and wreath-product extensions}

\begin{lemma}[Canonical wreath-product extension]\label{lem:wreath-extension}
Let $K$ be a finite group and $\theta\in\Irr(K)$.  For $m\ge 1$, put $K^m:=K\times\cdots\times K$ and
$\theta^{\boxtimes m}:=\theta\boxtimes\cdots\boxtimes\theta\in\Irr(K^m)$.  Let $\mathfrak S_m$ act on $K^m$
by permuting factors and form the semidirect product $K\wr \mathfrak S_m:=K^m\rtimes \mathfrak S_m$.

Then $\theta^{\boxtimes m}$ has a \emph{canonical} extension $\widetilde\theta^{(m)}\in\Irr(K\wr\mathfrak S_m)$
characterized as follows: if $(\rho,V)$ affords $\theta$, then $K^m$ acts on $V^{\otimes m}$ by the tensor
product representation and $\mathfrak S_m$ acts by permuting tensor factors:
\[
\sigma\cdot(v_1\otimes\cdots\otimes v_m):=v_{\sigma^{-1}(1)}\otimes\cdots\otimes v_{\sigma^{-1}(m)}.
\]
This yields a representation of $K\wr\mathfrak S_m$ on $V^{\otimes m}$ whose character is $\widetilde\theta^{(m)}$.
Moreover,
\[
\Res^{K\wr\mathfrak S_m}_{K^m}\widetilde\theta^{(m)}=\theta^{\boxtimes m}
\qquad\text{and hence}\qquad
\bigl\langle \Res^{K\wr\mathfrak S_m}_{K^m}\widetilde\theta^{(m)},\ \theta^{\boxtimes m}\bigr\rangle_{K^m}=1.
\]
\end{lemma}

\begin{proof}
Let $(\rho,V)$ afford $\theta$.  With the actions described in the statement,
$V^{\otimes m}$ becomes a representation of $K\wr \mathfrak S_m$, because the
permutation action of $\mathfrak S_m$ on tensor factors satisfies the semidirect
product relation
\[
\sigma (k_1,\dots,k_m)\sigma^{-1}=(k_{\sigma^{-1}(1)},\dots,k_{\sigma^{-1}(m)}).
\]
Its restriction to $K^m$ is exactly the external tensor product
$\theta^{\boxtimes m}$.

Since $\theta\in\Irr(K)$, the character $\theta^{\boxtimes m}$ is irreducible on
$K^m$.  As $K^m\triangleleft K\wr \mathfrak S_m$, any nonzero
$(K\wr \mathfrak S_m)$-submodule of $V^{\otimes m}$ is in particular a nonzero
$K^m$-submodule, hence must be all of $V^{\otimes m}$.  Therefore the above
representation is irreducible, and its character
$\widetilde\theta^{(m)}\in\Irr(K\wr \mathfrak S_m)$ extends
$\theta^{\boxtimes m}$.

The equality
\[
\Res^{K\wr \mathfrak S_m}_{K^m}\widetilde\theta^{(m)}=\theta^{\boxtimes m}
\]
is built into the construction, so the inner product statement follows
immediately.  Finally, the construction is basis-free and functorial in $(\rho,V)$,
so the resulting character depends only on $\theta$; hence the extension is canonical.
\end{proof}

\subsection{The general reductive case}
\begin{proposition}[Preferred characters arising out of a cuspidal unipotent of the identity component]
\label{prop:preferred-extension-abelian-component-general}
Let $\mathbf H^\circ$ be a connected reductive  group over $\overline{\F}_p$ equipped with a
Frobenius endomorphism $F$, and put $H^\circ:=\mathbf H^{\circ F}$.
Let $\mathbf H$ be an $F$--stable (possibly disconnected) algebraic group with identity component
$\mathbf H^\circ$ and set $H:=\mathbf H^F$.
Assume that $\pi_0(\mathbf H):=\mathbf H/\mathbf H^\circ$ is abelian (hence $H/H^\circ$ is abelian),
and assume that $H$ has a pinned semidirect component model in the sense of
Definition~\ref{def:pinned-semidirect-component-model}.  This holds for $F$-stable
semisimple centralizers by Remark~\ref{rmk:DM-supplies-pinned-component-models}.
Let $u^\circ\in\Irr(H^\circ)$ be a cuspidal unipotent character and let
\[
I \ :=\ I_H(u^\circ)\ =\ \{\,h\in H:\ {}^h u^\circ\simeq u^\circ\,\}.
\]

Then $u^\circ$ admits a preferred extension to $I$.
In particular, if $u^\circ$ is $H$--invariant (equivalently $I=H$), then $u^\circ$ admits a preferred
extension to $H$.
\end{proposition}

\begin{proof}
The pinned semidirect component model for $H$ restricts to $I$, by intersecting the chosen
complement with the inverse image of $I/H^\circ$.  Hence, replacing $H$ by $I$ and using this
restricted model, we may and do assume that $u^\circ$ is $H$--invariant, so that
$\Omega:=H/H^\circ$ is abelian and $u^\circ$ is $\Omega$--stable.

\medskip
Since $u^\circ$ is unipotent cuspidal, it factors through the adjoint quotient, so we may assume
$\mathbf H^\circ$ is semisimple of adjoint type.  Decompose
\[
H^\circ \ \cong\ \prod_{a\in\mathcal A} X_a
\]
as a direct product of finite groups of Lie type $X_a=\mathbf X_a^F$ coming from connected almost
simple adjoint algebraic groups $\mathbf X_a$ , possibly with repetitions.
Then every unipotent character of $H^\circ$ is an external tensor product, and cuspidality forces
each factor to be cuspidal; thus we may write
\[
u^\circ \ =\ \boxtimes_{a\in\mathcal A}\chi_a,
\qquad \chi_a\in\Uch(X_a)\ \text{cuspidal unipotent.}
\]
For each isomorphism type $X$ occurring among the $X_a$, and each cuspidal unipotent character
$\chi\in\Uch(X)$, let $m_{X,\chi}$ be the multiplicity of the pair $(X,\chi)$ among the factors
$\{(X_a,\chi_a)\}_{a\in\mathcal A}$.  Reordering factors gives a canonical refinement
\begin{equation}\label{eq:block-decomp}
H^\circ \ \cong\ \prod_{(X,\chi)} X^{m_{X,\chi}},
\qquad
u^\circ \ =\ \boxtimes_{(X,\chi)} \chi^{\boxtimes m_{X,\chi}},
\end{equation}
where $(X,\chi)$ runs over the finite set of pairs consisting of an absolute quasi-simple factor $X$ and a
cuspidal unipotent character $\chi$ of $X$ which appears in $u^\circ$.

\medskip

Conjugation by $H$ induces an action of $\Omega$ on the set of almost simple direct factors of
$H^\circ$.  Since $\Omega$ stabilizes $u^\circ$, it cannot permute a factor carrying $(X,\chi)$ into
a factor carrying $(X,\chi')$ with $\chi'\neq\chi$, because that would send $u^\circ$ to a distinct
tensor product character.  Hence the action of $\Omega$ preserves each block $X^{m_{X,\chi}}$ in
\eqref{eq:block-decomp}, and even preserves the partition of the $m_{X,\chi}$ copies by the label
$\chi$.  Consequently, for each pair $(X,\chi)$ the action of $\Omega$ on the block
$X^{m_{X,\chi}}$ factors through an abelian subgroup
\[
\Omega_{X,\chi}\ \le\ \Aut(X)\wr \mathfrak S_{m_{X,\chi}},
\]
whose permutation image lies in the Young subgroup permuting only the $m_{X,\chi}$ identical
$\chi$--factors.

Moreover, by Lemma~\ref{lem:cuspidal-unipotent-invariant}, $\chi$ is invariant under \emph{all}
automorphisms of $X$.  In particular, the character $\chi^{\boxtimes m_{X,\chi}}$ is stabilized by
$\Omega_{X,\chi}$.

\medskip
Fix a pair $(X,\chi)$ and put $m:=m_{X,\chi}$.
Let $A_{X,\chi}$ be the (abelian) projection of $\Omega_{X,\chi}$ to $\Aut(X)$ (i.e.\ the ``outer''
part acting on a single factor), and let $P_{X,\chi}$ be the (abelian) permutation image in
$\mathfrak S_m$.
Then the extension \(\widetilde{\chi}\) of \(\chi\) across the automorphism part $A_{X,\chi}$ follows from Proposition \ref{prop:preferred-extension-abelian-component} , i.e.,  the disconnected group $\mathbf X\rtimes A_{X,\chi}$ (whose component group is
abelian) yields a preferred extension
\[
\widetilde\chi\ \in\ \Irr(X\rtimes A_{X,\chi})
\qquad\text{with}\qquad
\Res^{X\rtimes A_{X,\chi}}_{X}\widetilde\chi=\chi.
\]

\smallskip

Consider the wreath product
\[
(X\rtimes A_{X,\chi})\wr \mathfrak S_m \ :=\ (X\rtimes A_{X,\chi})^m \rtimes \mathfrak S_m.
\]
By Lemma~\ref{lem:wreath-extension} (applied with $K=X\rtimes A_{X,\chi}$ and $\theta=\widetilde\chi$),
the character $\widetilde\chi^{\boxtimes m}\in\Irr\bigl((X\rtimes A_{X,\chi})^m\bigr)$ has a
canonical extension to $(X\rtimes A_{X,\chi})\wr\mathfrak S_m$.
Restrict this canonical extension to the subgroup
\[
X^m \rtimes \Omega_{X,\chi}\ \le\ (X\rtimes A_{X,\chi})\wr\mathfrak S_m
\]
(realizing $\Omega_{X,\chi}$ via its given action on the $m$ copies).  Denote the resulting
character by $\widetilde u_{X,\chi}\in\Irr(X^m\rtimes \Omega_{X,\chi})$.
\medskip
\noindent\textit{Irreducibility after restriction.}
Note that $X^m\triangleleft X^m\rtimes\Omega_{X,\chi}$ and $\Res^{\,X^m\rtimes\Omega_{X,\chi}}_{X^m}\tilde u_{X,\chi}=\chi^{\boxtimes m}\in\Irr(X^m)$. Thus any nonzero $(X^m\rtimes\Omega_{X,\chi})$--submodule is a nonzero $X^m$--submodule, so $\tilde u_{X,\chi}$ is irreducible.
By construction,
\[
\Res^{X^m\rtimes \Omega_{X,\chi}}_{X^m}\widetilde u_{X,\chi}=\chi^{\boxtimes m},
\]
and $\widetilde u_{X,\chi}$ is canonical once the preferred-extension and wreath-extension
conventions are fixed.

\medskip

Now, form the external tensor product over all pairs $(X,\chi)$ appearing in \eqref{eq:block-decomp}:
\[
\widetilde u \ :=\ \boxtimes_{(X,\chi)} \widetilde u_{X,\chi}.
\]
Then $\widetilde u$ is the desired canonical extension of $u^\circ$ to $H$.

\end{proof}
\begin{theorem}[Pinned canonical extension of cuspidal unipotents]
\label{thm:pinned-canonical-extension}
Let \(G\) be a connected reductive group over \(\F_q\), let \(s\in G^*\) be semisimple, and put \(H:=C_{G^*}(s)\) with identity component \(H^\circ\). Fix a pinning \(\cP\) of \(G\), and let \(\cP_H\) be the induced pinning of \(H^\circ\). Let \(u_\rho^\circ\in \Uch(H^\circ)\) be a cuspidal unipotent character, and set
\[
I:=\Stab_H(u_\rho^\circ).
\]
Then the pinning \(\cP\), together with Lusztig's preferred-extension
normalization fixed above, determines a distinguished extension
\[
        \dot u_\rho\in \Irr(I_H(u^\circ))
\]
of \(u^\circ\).  We will call this extension the \(\cP\)-preferred extension.
Consequently,
\[
u_\rho:=\Ind_I^H(\dot u_\rho)
\]
is a canonically determined irreducible character of \(H\) lying above the \(H\)-orbit of
\(u_\rho^\circ\).
\end{theorem}

\begin{proof}
The algebraic group \(\bH=C_{\bG^*}(s)\) is an \(F^*\)-stable semisimple centralizer.
Hence Lemma~\ref{lem:DM-semisimple-centralizer-splits} and
Remark~\ref{rmk:DM-supplies-pinned-component-models} give pinned semidirect component models for
\(H\) and for the inertia subgroup \(I\).  By Lemma~\ref{lem:H-component-group}, the quotient
\(H/H^\circ\) is abelian. Hence
\[
I/H^\circ \le H/H^\circ
\]
is abelian as well. Applying Proposition~\ref{prop:preferred-extension-abelian-component-general}
to \(H\) and the character \(u_\rho^\circ\), we obtain a canonical preferred
extension
\[
\dot u_\rho\in \Irr(I)
\]
of \(u_\rho^\circ\).

Now set
\[
u_\rho:=\Ind_I^H(\dot u_\rho).
\]
By Clifford theory, \(u_\rho\) is irreducible, and its restriction to \(H^\circ\) is the sum of the
\(H\)-conjugates of \(u_\rho^\circ\). Thus \(u_\rho\) lies above the \(H\)-orbit of \(u_\rho^\circ\).
Since \(\dot u_\rho\) is canonical, so is \(u_\rho\).
\end{proof}
\medskip
\section{The cuspidal part of the pinned Jordan decomposition}
\label{sec:cuspidal-pinned-JD}
The preceding section supplied the extension-theoretic input needed on the
unipotent side: after fixing the pinning, cuspidal unipotent characters admit
canonical preferred extensions to the relevant stabilisers.  We now use this
input to isolate the cuspidal part of the pinned Jordan decomposition.

\subsection{Canonical Jordan decomposition of cuspidal characters}
\begin{theorem}[Pinned Malle matching for simple groups]
\label{thm:pinned-Malle-matching-simple}
Let \(\mathbf G\) be simple and simply connected, equipped with a Frobenius
endomorphism \(F\), and put \(G=\mathbf G^F\).  Fix an \(F\)-stable pinning
\(\cP\) of \(\mathbf G\), and let \(\psi_{\cP}\) be the corresponding
non-degenerate character of the unipotent radical of the fixed Borel subgroup.
Let
\[
        A:=G_{\ad}/G
\]
denote the group of diagonal automorphisms of \(G\).  Let
\(s\in G^{*F}\), and let
\[
        \mathcal S_s\subset \cE(G,s)
\]
be the \(A\)-orbit of semisimple characters in \(\cE(G,s)\).  Denote by
\[
        \rho_{s,\psi_{\cP}}\in \mathcal S_s
\]
the Whittaker-normalized semisimple character determined by \(\psi_{\cP}\).

Let
\[
        F_O\subset \cE(G,s)_{\cusp}
\]
be an \(A\)-orbit of cuspidal characters.  Then the pinning \(\cP\) determines
a distinguished element
\[
        \rho_{O,\cP}\in F_O
\]
such that
\[
        \Aut(G)_{\rho_{O,\cP}}
        =
        \Aut(G)_{\rho_{s,\psi_{\cP}}}.
\]
Consequently there is a unique \(A\)-equivariant bijection
\[
        m_{\cP,O}:F_O\xrightarrow{\sim}\mathcal S_s
\]
satisfying
\[
        m_{\cP,O}(\rho_{O,\cP})=\rho_{s,\psi_{\cP}}.
\]
Moreover, for every \(\rho\in F_O\), one has
\[
        \Aut(G)_\rho
        =
        \Aut(G)_{m_{\cP,O}(\rho)}.
\]
We call \(m_{\cP,O}\) the \(\cP\)-preferred Malle matching on the orbit \(F_O\).
\end{theorem}
\begin{proof}
Write the action of \(a\in A\) on characters as
\[
        \eta^a:=\eta\circ \ad(a).
\]
Malle's theorem gives, for each cuspidal character
\(\rho\in \cE(G,s)\), a semisimple character
\(\chi\in \cE(G,s)\) satisfying
\[
        \Aut(G)_\chi=\Aut(G)_\rho.
\]
In particular, the \(A\)-orbit \(F_O\) is isomorphic, as an \(A\)-set, to the
\(A\)-orbit \(\mathcal S_s\) of semisimple characters.  The point is to choose
the isomorphism in a way normalized by the pinning.

We first treat type \(A\).  In this case the diagonal automorphism group may
have order larger than two, so the stabilizer condition alone does not choose
a distinguished point.  Instead one uses the Whittaker normalization.  By the
type \(A\) equivariant character correspondence of Cabanes--Späth, as used in
Malle's proof, each diagonal orbit of cuspidal characters in \(\cE(G,s)\)
contains a unique \(\psi_{\cP}\)-generic constituent.  We define
\[
        \rho_{O,\cP}
\]
to be this constituent.  The semisimple orbit \(\mathcal S_s\) similarly
contains the unique \(\psi_{\cP}\)-generic semisimple character
\(\rho_{s,\psi_{\cP}}\).  The type \(A\) equivariance theorem gives
\[
        \Aut(G)_{\rho_{O,\cP}}
        =
        \Aut(G)_{\rho_{s,\psi_{\cP}}}.
\]

We now assume that \(\mathbf G\) is not of type \(A\).  Let
\[
        S_\rho:=\Aut(G)_\rho,\qquad A_\rho:=S_\rho\cap A,
\]
and let \(B_\rho\) be the image of \(S_\rho\) in the graph-field part of
\(\Out(G)\).  If \(\chi\in\mathcal S_s\) is one semisimple character satisfying
\[
        \Aut(G)_\chi=S_\rho,
\]
then the other diagonal translates of \(\chi\) with the same stabilizer
\(S_\rho\) are parametrized by
\[
        (A/A_\rho)^{B_\rho}.
\]
For the non-type \(A\) groups occurring in Malle's reduction, this group is
either trivial or of order two.

If
\[
        (A/A_\rho)^{B_\rho}=1,
\]
then there is no residual choice: the semisimple character in
\(\mathcal S_s\) whose full automorphism stabilizer equals \(S_\rho\) is unique.
Equivalently, there is a unique element of \(F_O\) whose stabilizer equals
\(\Aut(G)_{\rho_{s,\psi_{\cP}}}\).  We define this element to be
\(\rho_{O,\cP}\).

It remains to discuss the residual twofold cases.  These are precisely the
cases in Malle's proof where the stabilizer condition leaves two diagonal
translates undistinguished.  Malle's proof treats them by explicit Brauer-tree
chains, by Deligne--Lusztig induction chains with multiplicity-one
constituents, and by the dihedral actions occurring in type \(D\).  The pinning
removes the remaining sign ambiguity as follows.

In the Brauer-tree cases, Malle connects the cuspidal character to the
semisimple character by a labelled chain of non-exceptional vertices on
Brauer trees.  We root this chain at the pinned semisimple endpoint
\(\rho_{s,\psi_{\cP}}\).  Since the labels of hooks, cohooks, and the relevant
simple factors are fixed by the pinned root datum, this rooted chain has a
unique cuspidal endpoint in \(F_O\).  We define this endpoint to be
\(\rho_{O,\cP}\).  Malle's Brauer-tree argument then gives equality of the
full automorphism stabilizers.

In the Deligne--Lusztig induction cases, choose the Levi subgroup \(L\) to be
the standard Levi determined by the pinned root datum.  The character on \(L\)
appearing in Malle's induction chain is chosen with the corresponding pinned
Whittaker normalization.  The relevant constituents of
\[
        R_L^G(\psi)
\]
occur with multiplicity one.  Hence, among the two possible diagonal translates
in the residual orbit, the pinned induction datum selects a unique constituent.
This constituent is defined to be \(\rho_{O,\cP}\), and Malle's
Deligne--Lusztig argument gives
\[
        \Aut(G)_{\rho_{O,\cP}}
        =
        \Aut(G)_{\rho_{s,\psi_{\cP}}}.
\]

Finally, in the type \(D\) cases where the residual ambiguity is spinorial, the
non-trivial ambiguity is the interchange of the two half-spin labels, or its
graph-field analogue.  The pinning labels the two terminal vertices of the
\(D_n\)-diagram.  We choose the constituent corresponding to the preferred
terminal vertex determined by \(\cP\).  Malle's case analysis shows that the
cuspidal and semisimple diagonal orbits carry the same dihedral action of the
relevant graph-field automorphisms.  Therefore this pinned choice again has
the same full automorphism stabilizer as the pinned semisimple character.

Thus in every case we have constructed a distinguished element
\[
        \rho_{O,\cP}\in F_O
\]
with
\[
        \Aut(G)_{\rho_{O,\cP}}
        =
        \Aut(G)_{\rho_{s,\psi_{\cP}}}.
\]
In particular,
\[
        \Stab_A(\rho_{O,\cP})
        =
        \Stab_A(\rho_{s,\psi_{\cP}}).
\]
There is therefore a unique \(A\)-equivariant bijection
\[
        m_{\cP,O}:F_O\longrightarrow \mathcal S_s
\]
sending \(\rho_{O,\cP}\) to \(\rho_{s,\psi_{\cP}}\), namely
\[
        m_{\cP,O}\bigl(\rho_{O,\cP}^a\bigr)
        :=
        \rho_{s,\psi_{\cP}}^a
        \qquad(a\in A).
\]
This formula is well-defined because the two basepoints have the same
\(A\)-stabilizer.  It is visibly \(A\)-equivariant and is the only
\(A\)-equivariant map with the prescribed value at \(\rho_{O,\cP}\).

For any \(a\in A\), conjugating the equality
\[
        \Aut(G)_{\rho_{O,\cP}}
        =
        \Aut(G)_{\rho_{s,\psi_{\cP}}}
\]
by \(a\) gives
\[
        \Aut(G)_{\rho_{O,\cP}^a}
        =
        \Aut(G)_{\rho_{s,\psi_{\cP}}^a}.
\]
Hence
\[
        \Aut(G)_\rho
        =
        \Aut(G)_{m_{\cP,O}(\rho)}
\]
for all \(\rho\in F_O\), as claimed.
\end{proof}
\begin{lemma}\label{lemma:bijection-cuspidal}
Fix a pinning \(\cP\) of \((\bG,F)\). Then there exists a canonically defined bijection
\[
\J_{s,\cP}:\cE(G,s)_{\cusp}\xrightarrow{\ \sim\ }\Uch(H)_{\cusp}
\]
satisfying
\[
\langle R_{\bT^{*}}^{\bG}(s), \rho \rangle_G = \epsilon_{\bG} \epsilon_{\bH} \langle R_{\bT^{*}}^{\bH}(1), \J_s(\rho) \rangle_H.
\]
We call \(\J_{s,\cP}\) the preferred bijection attached to \(\cP\).
\end{lemma}
\begin{proof}
Let
\[
\bar J_s:\cE(G,s)_{\cusp}\longrightarrow
\Uch(H^\circ)_{\cusp}/\!\sim(H/H^\circ)
\]
be Lusztig's disconnected Jordan decomposition restricted to the cuspidal part. This is independent of the choice of $J_s$
For an \(H/H^\circ\)-orbit
\[
O\subset \Uch(H^\circ)_{\cusp},
\]
put
\[
F_O:=\bar J_s^{-1}(O)\subset \cE(G,s)_{\cusp}.
\]

Choose \(u_O^\circ\in O\), and set
\[
I_O:=\Stab_H(u_O^\circ),\qquad \Omega_O:=I_O/H^\circ .
\]
By Theorem~\ref{thm:pinned-canonical-extension}, \(u_O^\circ\) admits a canonical preferred
extension
\[
\dot u_O\in \Irr(I_O).
\]
For \(\xi\in \Omega_O^\vee=\Irr(\Omega_O)\), define
\[
u_{O,\xi}:=\Ind_{I_O}^H(\dot u_O\otimes \xi)\in \Irr(H).
\]
By Clifford theory, each \(u_{O,\xi}\) is irreducible, and the map
\[
\Omega_O^\vee\longrightarrow \Uch(H)_O,\qquad
\xi\longmapsto u_{O,\xi},
\]
is a bijection onto
\[
\Uch(H)_O:=
\left\{
u\in \Uch(H)_{\cusp}:\Res^H_{H^\circ}u=\sum_{\lambda\in O}\lambda
\right\}.
\]

This parametrization is independent of the chosen representative \(u_O^\circ\in O\).
Indeed, if \(u_1^\circ={}^h u_O^\circ\) for some \(h\in H\), then
\[
I_1=hI_Oh^{-1}.
\]
Since \(H/H^\circ\) is abelian, conjugation by \(h\) induces the identity on
\[
I_O/H^\circ=\Omega_O.
\]
Moreover, by canonicity of the preferred extension, the preferred extension of \(u_1^\circ\) is
\[
{}^h\!\dot u_O.
\]
Hence for every \(\xi\in \Omega_O^\vee\),
\[
\Ind_{I_1}^H\bigl({}^h\!\dot u_O\otimes \xi\bigr)
=
\Ind_{I_O}^H(\dot u_O\otimes \xi),
\]
so the characters \(u_{O,\xi}\) depend only on the orbit \(O\). In particular,
\begin{equation}\label{eq:HO-torsor-rewrite}
|\Uch(H)_O|=|\Omega_O|.
\end{equation}

Now put
\[
A:=G_{\ad}/G.
\]
If \(\rho,\rho'\in F_O\), then \(\rho'=\rho\circ \ad(a)\) for some \(a\in A\), and since \(A\) is
abelian we have
\[
\Stab_A(\rho')=\Stab_A(\rho).
\]
Thus the stabilizer \(A_\rho\) depends only on the orbit \(O\); we denote it by
\[
A_O:=A_\rho\qquad (\rho\in F_O).
\]

By the construction of the disconnected Jordan decomposition together with the orthogonality
statement in Remark~\ref{rmk:multi-regular-embedding-action}, the subgroup \(A_O\) is the orthogonal complement of
\(\Omega_O\) under the pairing induced by
\[
\alpha:A\longrightarrow (H/H^\circ)^\vee.
\]
Hence \(\alpha\) induces an isomorphism
\begin{equation}\label{eq:alpha-fiber-iso-rewrite}
\bar\alpha_O:\ A/A_O \xrightarrow{\ \sim\ } \Omega_O^\vee.
\end{equation}
Since \(F_O\) is a single \(A\)-orbit and \(|F_O|=|\Omega_O|\) by
Theorem~\ref{thm:Lusztig-orbit-JD-disconnected-centralizers}, we obtain
\begin{equation}\label{eq:FO-size-rewrite}
|F_O|=|A/A_O|=|\Omega_O|.
\end{equation}

For every \(\rho\in \cE(G,s)_{\cusp}\), with \(O=\bar J_s(\rho)\), we now attach a class
\[
c_\rho\in A/A_O
\]
canonically determined by \(\rho\) and the chosen pinning \(\cP\).

Case (1): If \(\bG\) is simple and simply connected, let \(c_\rho\) be the class \(g_\rho A_\rho\) obtained
as follows: let
\[
        \chi_\rho:=m_{\cP,O}(\rho)\in \mathcal S_s
\]
be the \(\cP\)-preferred semisimple character attached to \(\rho\) by
Theorem~\ref{thm:pinned-Malle-matching-simple}. Let \(\rho_{s,\psi}\) be the semisimple character attached to the pinned Whittaker datum, and choose
\(g_\rho\in G_{\ad}\) such that
\[
\chi_\rho\circ \ad(g_\rho)=\rho_{s,\psi}.
\]
Then \(c_\rho:=g_\rho A_\rho\).

Case (2): If \(\bG\) is semisimple and simply connected, write
\[
\bG=\prod_i \bG_i,\qquad \rho=\boxtimes_i \rho_i,
\]
and define \(c_\rho\) as the product of the classes attached to the factors by the preceding case.

Case~(3): If \(\bG_{\der}\) is simply connected, choose a semisimple character
\[
\chi_0\in \cE(G_0,s_0)
\]
corresponding to \(\rho_0\). Let
\[
\chi_\rho\in \cE(G,s)
\]
be the unique semisimple character lying above \(\chi_0\), whose existence and
uniqueness are given by \cite[Lem.~4.1(b)]{Malle2017}.  The reduction argument in
\cite[proof of Thm.~1, immediately after Lem.~4.1]{Malle2017} shows that the
stabilizer-matching statement for \(G_0=[G,G]\) passes to \(G\): indeed,
\(G_0\) is characteristic in \(G\), and both \(\rho\) and \(\chi_\rho\) are uniquely
determined, inside the fixed Lusztig series \(\cE(G,s)\), by their constituents
\(\rho_0\) and \(\chi_0\) on \(G_0\).  Since \(\rho_0\) and \(\chi_0\) have the same
stabilizer under diagonal automorphisms by the simply connected derived case, it
follows that
\[
\Stab_A(\chi_\rho)=\Stab_A(\rho)=A_\rho .
\]
Hence
\[
\Stab_A(\chi_\rho)=\Stab_A(\rho)=A_\rho.
\]
Choose \(g_\rho\in G_{\ad}\) such that
\[
\chi_\rho\circ \ad(g_\rho)=\rho_{s,\psi},
\]
and set
\[
c_\rho:=g_\rho A_\rho.
\]
Case (4): Finally, for general \(G\), choose any \(z\)-extension
\[
1\longrightarrow Z\longrightarrow G'\xrightarrow{\,f\,} G\longrightarrow 1
\]
defined over \(\F_q\), with \(Z\) a central torus and \(G'_{\der}\) simply connected. Choose a lift
\(s'\in G'^*\) of \(s\), and let \(\rho'\) be the inflation of \(\rho\) to \(G'^F\). Applying the
preceding case to \((G',\rho')\), we obtain a class
\[
c_{\rho'}\in G'_{\ad}/G'.
\]
Via the canonical identification
\[
G'_{\ad}/G'\xrightarrow{\sim} G_{\ad}/G=A
\]
of Remark~\ref{rmk:regular-embedding-adjoint-action}, we transport \(c_{\rho'}\) to a class \(c_\rho\in A/A_O\).
By Lemma~\ref{lem:zext-independence} below, this class is independent of the chosen \(z\)-extension.
Thus \(c_\rho\) is canonically determined in all cases.

The construction is \(A\)-equivariant: for \(a\in A\),
\begin{equation}\label{eq:crho-equiv}
c_{\rho\circ \ad(a)}=a^{-1}c_\rho.
\end{equation}
In Cases~(1)--(3) this is immediate from the definitions, and in Case~(4) it
follows after inflation to a \(z\)-extension and descent via Lemma~\ref{lem:zext-independence}.

We may therefore define
\[
\J_{s,\cP}(\rho):=u_{O,\bar\alpha_O(c_\rho)},
\qquad O=\bar J_s(\rho).
\]
Everything entering this definition is canonical once \(\cP\) is fixed, so
\(\J_{s,\cP}\) is canonically defined.

It remains to prove that \(\J_{s,\cP}\) is a bijection. Fix an orbit \(O\), and choose
\(\rho_O\in F_O\). Every \(\rho\in F_O\) may be written uniquely in the form
\[
\rho=\rho_O\circ \ad(a)
\]
with \(aA_O\in A/A_O\). By \eqref{eq:crho-equiv},
\[
c_\rho=a^{-1}c_{\rho_O}.
\]
Hence
\[
\J_{s,\cP}(\rho)=u_{O,\bar\alpha_O(a^{-1}c_{\rho_O})}.
\]
Translation by the fixed class \(c_{\rho_O}\) is a bijection of \(A/A_O\), and
\eqref{eq:alpha-fiber-iso-rewrite} is an isomorphism. Therefore the map
\[
F_O\longrightarrow \Uch(H)_O,\qquad
\rho\longmapsto \J_{s,\cP}(\rho),
\]
is a bijection. Taking the disjoint union over all \(H/H^\circ\)-orbits \(O\) yields a bijection
\[
\J_{s,\cP}:\cE(G,s)_{\cusp}\xrightarrow{\ \sim\ }\Uch(H)_{\cusp}.
\]

Finally, let \(\rho\in F_O\). By construction,
\[
\Res^H_{H^\circ}\J_{s,\cP}(\rho)=\sum_{\lambda\in O}\lambda.
\]
Therefore Theorem~\ref{thm:Lusztig-orbit-JD-disconnected-centralizers}(3) gives
\[
\langle R_{\bT^*}^{\bG}(s),\rho\rangle_G
=
\epsilon_{\bG}\epsilon_{\bH^\circ}
\sum_{\lambda\in O}
\langle R_{\bT^*}^{\bH^\circ}(1),\lambda\rangle_{H^\circ}
=
\epsilon_{\bG}\epsilon_{\bH^\circ}
\Bigl\langle
R_{\bT^*}^{\bH^\circ}(1),\Res^H_{H^\circ}\J_{s,\cP}(\rho)
\Bigr\rangle_{H^\circ}.
\]
By the definition of \(R_{\bT^*}^{\bH}(1)\) for the disconnected group \(H\), the last expression is
\[
\epsilon_{\bG}\epsilon_{\bH}\,
\langle R_{\bT^*}^{\bH}(1),\J_{s,\cP}(\rho)\rangle_H.
\]
This is the required formula.
\end{proof}
\begin{lemma}[Independence of the \(z\)-extension]
\label{lem:zext-independence}
Let
\[
1\longrightarrow Z_i \longrightarrow G_i' \xrightarrow{\,f_i\,} G \longrightarrow 1
\qquad (i=1,2)
\]
be two \(F\)-stable \(z\)-extensions of \(G\), with \(Z_i\) a connected central torus and
\((G_i')_{\der}\) simply connected. Fix a pinning \(\cP\) of \((\bG,F)\), and equip each
\(G_i'\) with the induced pinning. Let \(s_i'\in (G_i'^*)^F\) be the image of \(s\) under the
dual regular embedding \(G^* \hookrightarrow G_i'^*\), and let
\[
\rho_i' := \rho\circ f_i^F \in \cE(G_i',s_i')_{\cusp}
\]
be the inflation of \(\rho\in \cE(G,s)_{\cusp}\).

Apply Case~\textup{(3)} to \((G_i',\rho_i')\), and let
\[
c_i(\rho)\in G_{\ad}/G
\]
be the resulting class, via the natural identification of the corresponding diagonal-automorphism
quotients with \(G_{\ad}/G\). Then
\[
c_1(\rho)=c_2(\rho).
\]
Consequently the character
\[
\J_s(\rho)=u_\rho\otimes \alpha(c_i(\rho))
\]
is independent of the chosen \(z\)-extension.
\end{lemma}

\begin{proof}
Consider the fibre product
\[
\widetilde G:= G_1' \times_G G_2'
=
\{(g_1,g_2)\in G_1'\times G_2' \mid f_1(g_1)=f_2(g_2)\}.
\]
Then \(\widetilde G\) is a connected reductive \(F\)-stable group, and the projections
\[
p_i:\widetilde G \twoheadrightarrow G_i'
\]
have kernels
\[
\ker(p_1)=\{(1,z_2):z_2\in Z_2\}\cong Z_2,
\qquad
\ker(p_2)=\{(z_1,1):z_1\in Z_1\}\cong Z_1,
\]
which are connected central tori. Thus \(\widetilde G\) is again a \(z\)-extension of each
\(G_i'\), and also of \(G\).

Moreover, each \(p_i\) induces an isomorphism on derived groups:
\[
p_i:\widetilde G_{\der}\xrightarrow{\sim} (G_i')_{\der}.
\]
Indeed, surjectivity follows from the surjectivity of \(p_i\), while
\(\ker(p_i)\cap \widetilde G_{\der}\) is a connected subgroup of the torus \(\ker(p_i)\) contained
in the semisimple group \(\widetilde G_{\der}\), hence is trivial. In particular,
\(\widetilde G_{\der}\) is simply connected.

It is therefore enough to prove the following functoriality statement: if
\[
f:\widetilde G \twoheadrightarrow G'
\]
is a surjective \(F\)-morphism of connected reductive groups whose kernel is a connected central
torus, and if both \(\widetilde G_{\der}\) and \(G'_{\der}\) are simply connected, then the class
constructed in Case~\textup{(3)} for the inflated character on \(\widetilde G\) is the same as the
class constructed for \(G'\).

So let \(f:\widetilde G\twoheadrightarrow G'\) be such a morphism, let \(s'\in (G'^*)^F\) be the
image of \(s\), let \(\widetilde s\in (\widetilde G^*)^F\) be the image of \(s'\) under the dual
regular embedding \(G'^*\hookrightarrow \widetilde G^*\), and let
\[
\rho'\in \cE(G',s')_{\cusp},
\qquad
\widetilde\rho:=\rho'\circ f^F\in \cE(\widetilde G,\widetilde s)_{\cusp}.
\]
Write
\[
G_0':=G'_{\der},\qquad \widetilde G_0:=\widetilde G_{\der}.
\]
As observed above, the induced morphism
\[
f_0:\widetilde G_0\xrightarrow{\sim} G_0'
\]
is an isomorphism.

Since Case~\textup{(3)} is already canonical once the pinning is fixed, we may compute both
classes using compatible auxiliary choices. Choose an irreducible constituent
\[
\rho_0\in \Irr(G_0'^F\mid \rho')
\]
and transport it via \(f_0\) to
\[
\widetilde\rho_0:=\rho_0\circ f_0 \in \Irr(\widetilde G_0^F\mid \widetilde\rho).
\]
Applying Case~\textup{(2)} to \(G_0'\) and \(\widetilde G_0\), we obtain corresponding semisimple
characters
\[
\chi_0\in \cE(G_0',s_0') ,\qquad
\widetilde\chi_0:=\chi_0\circ f_0 \in \cE(\widetilde G_0,\widetilde s_0),
\]
with the same stabilizers as \(\rho_0\) and \(\widetilde\rho_0\), respectively.

Now apply \cite[Lem.~4.1(b)]{Malle2017}. There is a unique semisimple character
\[
\rho_s' \in \cE(G',s')
\]
lying above \(\chi_0\), and a unique semisimple character
\[
\widetilde\rho_s \in \cE(\widetilde G,\widetilde s)
\]
lying above \(\widetilde\chi_0\). Since \(f^F\) is surjective and \(\widetilde\chi_0\) is a constituent
of the restriction of \(\infl_f(\rho_s')\) to \(\widetilde G_0^F\), the uniqueness in
\cite[Lem.~4.1(b)]{Malle2017} forces
\[
\widetilde\rho_s=\infl_f(\rho_s').
\]

Let \(\rho'_{s,\psi}\) and \(\widetilde\rho_{s,\psi}\) denote the semisimple characters attached to the
pinned Whittaker data on \(G'\) and \(\widetilde G\), respectively. Because the pinning on
\(\widetilde G\) is the pullback of that on \(G'\), the corresponding maximal unipotent subgroups
and nondegenerate characters are compatible under \(f\). Hence the Gelfand--Graev characters are
related by inflation, and therefore so are their Alvis--Curtis duals:
\[
\widetilde\rho_{s,\psi}=\infl_f(\rho'_{s,\psi}).
\]

Let \(g\in G'_{\ad}\). Via the canonical identification \(\widetilde G_{\ad}\cong G'_{\ad}\), we may
also regard \(g\) as an element of \(\widetilde G_{\ad}\). Then
\[
\rho_s'\circ \ad(g)=\rho'_{s,\psi}
\]
if and only if
\[
\widetilde\rho_s\circ \ad(g)=\widetilde\rho_{s,\psi},
\]
because inflation commutes with conjugation and \(f^F:\widetilde G^F\twoheadrightarrow G'^F\) is
surjective. Thus the same element \(g\) is admissible in the Case~\textup{(3)} construction for
both \(G'\) and \(\widetilde G\), and so the resulting classes coincide.

Applying this functoriality to the two projections
\[
p_i:\widetilde G\to G_i'
\qquad (i=1,2)
\]
gives
\[
c_1(\rho)=c_{\widetilde G}(\widetilde\rho)=c_2(\rho).
\]
This proves the lemma. The final assertion follows immediately, since \(u_\rho\) depends only on the
\(H/H^\circ\)-orbit attached to \(\rho\).
\end{proof}

\section{Preliminaries about disconnected groups}
This section recalls some basic facts about disconnected reductive groups over \(\mathbb{F}_q\). The contents are taken from \cite{Bonnafe1999Produit}.
\medskip
\subsection{Disconnected reductive groups over $\mathbb{F}_q$}
\label{subsec:disc-red}

\paragraph{General set-up.}
Let $\mathbf G$ be a (not necessarily connected) reductive algebraic group defined over $\F_q$
with Frobenius endomorphism $F:\mathbf G\to \mathbf G$.
Write $\mathbf G^\circ$ for the identity component and set
\[
G:=\mathbf G^F,\qquad G^\circ := (\mathbf G^\circ)^F,\qquad
\Omega := G/G^\circ \cong (\mathbf G/\mathbf G^\circ)^F .
\]

\subsubsection{Parabolics, Levis, quasi-Borels, quasi-tori, Weyl groups}

\begin{definition}[Parabolics and Levis in $\mathbf G$ {\cite[Def.~6.1.1--6.1.2]{Bonnafe1999Produit}}]\label{def:disconn-parabolic}
A \emph{parabolic subgroup} of $\mathbf G$ is a closed subgroup $\mathbf P\le \mathbf G$
containing a Borel subgroup of $\mathbf G^\circ$.
Let $\mathbf U:=R_u(\mathbf P^\circ)$ be the unipotent radical of $\mathbf P^\circ$.
Choose a Levi subgroup $\mathbf L_0\le \mathbf P^\circ$ and put
\[
\mathbf L:=N_{\mathbf P}(\mathbf L_0).
\]
Then $\mathbf L^\circ=\mathbf L_0$ and $\mathbf P=\mathbf L\ltimes \mathbf U$.
We call $\mathbf L$ a \emph{Levi subgroup} of $\mathbf P$.
A subgroup $\mathbf L\le \mathbf G$ is called \emph{regular} if it is a Levi of some parabolic of $\mathbf G$.
\end{definition}

\begin{definition}[Quasi-Borels and quasi-maximal tori {\cite[Def.~6.1.3]{Bonnafe1999Produit}}]
A \emph{quasi-Borel} of $\mathbf G$ is the normalizer $N_{\mathbf G}(\mathbf B^\circ)$
of a Borel subgroup $\mathbf B^\circ\le \mathbf G^\circ$.
A \emph{quasi-maximal torus} of $\mathbf G$ is a Levi subgroup of a quasi-Borel.
\end{definition}

\begin{definition}[Weyl group of $\mathbf G$ relative to a torus {\cite[Def.~6.1.4]{Bonnafe1999Produit}}]
If $\mathbf T^\circ$ is a maximal torus of $\mathbf G^\circ$, set
\[
W_{\mathbf G}(\mathbf T^\circ):=N_{\mathbf G}(\mathbf T^\circ)/\mathbf T^\circ .
\]
If $\mathbf B^\circ\le \mathbf G^\circ$ is a Borel containing $\mathbf T^\circ$ and we put
$\mathbf T:=N_{\mathbf G}(\mathbf T^\circ,\mathbf B^\circ)$,
then $\mathbf T$ is a quasi-maximal torus of $\mathbf G$ with $(\mathbf T)^\circ=\mathbf T^\circ$ and
\[
W_{\mathbf G^\circ}(\mathbf T^\circ)
=
N_{\mathbf G^\circ}(\mathbf T^\circ)/\mathbf T^\circ
\trianglelefteq
W_{\mathbf G}(\mathbf T^\circ)
=
N_{\mathbf G}(\mathbf T^\circ)/\mathbf T^\circ .
\]
Moreover,  multiplication induces an isomorphism
\[
W_{\mathbf G^\circ}(\mathbf T^\circ)\rtimes
(\mathbf T/\mathbf T^\circ)
\xrightarrow{\;\sim\;}
W_{\mathbf G}(\mathbf T^\circ),
\]
where \(\mathbf T/\mathbf T^\circ\) acts on
\(W_{\mathbf G^\circ}(\mathbf T^\circ)\) by conjugation.
\end{definition}

\subsubsection{Pinning and pinned representatives for the component action}

\paragraph{Pinning of $\mathbf G^\circ$.}
Fix an $F$-stable Borel $\mathbf B^\circ\le \mathbf G^\circ$ and an $F$-stable maximal torus
$\mathbf T^\circ\le \mathbf B^\circ$.
Let $\Phi$ be the root system of $\mathbf G^\circ$ relative to $\mathbf T^\circ$ and
$\Delta\subset\Phi$ the simple roots determined by $\mathbf B^\circ$.
A (standard) \emph{pinning} of $\mathbf G^\circ$ is the data
\[
\mathcal P=(\mathbf G^\circ,\mathbf B^\circ,\mathbf T^\circ,(x_\alpha)_{\alpha\in\Delta}),
\]
where $x_\alpha:\mathbb G_a\to \mathbf U_\alpha$ are root group isomorphisms for the simple roots.
We say $\mathcal P$ is \emph{$F$-stable} if $F$ permutes $\Delta$ and
$F(x_\alpha(t))=x_{F(\alpha)}(t^q)$ (after identifying $F$ on $\mathbb G_a$ with $t\mapsto t^q$).

\begin{lemma}[Bonnafé {\cite[Lem.~6.2.1]{Bonnafe1999Produit}}]
Let $\mathbf T:=N_{\mathbf G}(\mathbf T^\circ,\mathbf B^\circ)$.
There exists a closed $F$-stable subgroup $\mathbf A\le \mathbf T$ such that
\begin{enumerate}[label=\textup{(\alph*)}]
\item $\mathbf G=\mathbf G^\circ\cdot \mathbf A$;
\item $\mathbf A\cap \mathbf G^\circ\subset Z(\mathbf G^\circ)$;
\item every $a\in \mathbf A$ induces a \emph{quasi-central} automorphism of $\mathbf G^\circ$
(in the sense of Digne--Michel).
\end{enumerate}
Moreover one may take
\[
\mathbf A=\bigl\{a\in \mathbf T \ \big|\
a\,x_\alpha(t)\,a^{-1}=x_{a(\alpha)}(t)\ \text{for all }\alpha\in\Delta,\ t\in\mathbb G_a\bigr\},
\]
so that conjugation by $\mathbf A$ preserves the chosen pinning $\mathcal P$.
\end{lemma}

\begin{definition}[Pinned representatives of the algebraic component action]\label{def:pinned-Omega-action}
Conjugation by $\mathbf G$ induces a homomorphism
\[
\Omega=\mathbf G/\mathbf G^\circ \longrightarrow \Out(\mathbf G^\circ).
\]
The pinning $\mathcal P$ identifies $\Aut(\mathbf G^\circ,\mathcal P)$ with the automorphism group of the
\emph{based} root datum of $(\mathbf G^\circ,\mathbf T^\circ,\Delta)$ and yields the usual
\emph{pinned lift} section $\Out(\mathbf G^\circ)\to \Aut(\mathbf G^\circ,\mathcal P)$.
Composing, we obtain a canonical homomorphism
\[
\Omega \longrightarrow \Aut(\mathbf G^\circ,\mathcal P),\qquad \omega\longmapsto \sigma_\omega,
\]
whose image consists of pinning-preserving automorphisms representing the algebraic outer action class of $\omega$.
This definition records the algebraic outer action only.
\end{definition}

\subsubsection{Lusztig induction/restriction for disconnected groups}

\paragraph{Deligne--Lusztig varieties and Lusztig functors.}
Let $\mathbf P$ be a parabolic subgroup of $\mathbf G$ with unipotent radical $\mathbf U:=R_u(\mathbf P^\circ)$,
and let $\mathbf L$ be an $F$-stable Levi subgroup of $\mathbf P$.
Set
\[
Y_{\mathbf U}:=\{g\in \mathbf G \mid g^{-1}F(g)\in \mathbf U\}.
\]
Then $G$ acts on $Y_{\mathbf U}$ by left translation and $L:=\mathbf L^F$ acts by right translation.
The virtual $(G,L)$--bimodule $H_c^\ast(Y_{\mathbf U})$ yields Lusztig induction and restriction functors
\[
R_{\mathbf L\subset \mathbf P}^{\mathbf G} \quad\text{and}\quad {}^\ast R_{\mathbf L\subset \mathbf P}^{\mathbf G}
\]
between Grothendieck groups (or between spaces of central class functions), as in
{\cite[\S6.3]{Bonnafe1999Produit}} (cf.\ {\cite[Def.~2.2]{DigneMichel94}}.
If $\mathbf L\supset \mathbf G^\circ$ then $\mathbf U=1$ and
\[
R_{\mathbf L\subset \mathbf L}^{\mathbf G}=\Ind_L^G,\qquad
{}^\ast R_{\mathbf L\subset \mathbf L}^{\mathbf G}=\Res_L^G
\]
{\cite[Prop.~6.3.2]{Bonnafe1999Produit}}.
These functors satisfy the usual transitivity for chains of parabolics
{\cite[Prop.~6.3.3]{Bonnafe1999Produit}}.

\medskip
We now state a Lemma of Digne–Michel \cite[Cor. 2.4]{DigneMichel94} to be used in \S \ref{sec:disc-full-JD}.
\begin{lemma}[Restriction to the identity component for Lusztig functors]\label{lem:DM94-restriction-to-identity-component}
Let $G$ be a (possibly disconnected) reductive $\F_q$--group, $P$ an $F$--stable ``parabolic'' of $G$
(in the sense of \cite[Def.\ 1.4]{DigneMichel94}) with Levi $L$ and unipotent radical $U=R_u(P^\circ)\subset G^\circ$.
Let $P^\circ=P\cap G^\circ$ and $L^\circ=L\cap G^\circ$.
Then for any class function (or virtual character) $\varphi$ of $L^F$ one has, as class functions on $G^{\circ F}$,
\[
\Res^{G^F}_{G^{\circ F}}\bigl(R^G_{L\subset P}(\varphi)\bigr)
\;=\;
\sum_{a\in G^F/(L^F\cdot G^{\circ F})} {}^{a}\!\Bigl(R^{G^\circ}_{L^\circ\subset P^\circ}\bigl(\Res^{L^F}_{L^{\circ F}}\varphi\bigr)\Bigr),
\]
and for any class function $\psi$ of $G^F$,
\[
\Res^{L^F}_{L^{\circ F}}\bigl({}^*R^G_{L\subset P}(\psi)\bigr)
\;=\;
{}^*R^{G^\circ}_{L^\circ\subset P^\circ}\bigl(\Res^{G^F}_{G^{\circ F}}\psi\bigr).
\]
\end{lemma}

\begin{corollary}\label{cor:HC-induction-restricts-contains-connected}
With the notation of Lemma~\ref{lem:DM94-restriction-to-identity-component},
$\Res^{G^F}_{G^{\circ F}} R^G_{L\subset P}(\tau)$ contains $R^{G^\circ}_{L^\circ\subset P^\circ}(\tau^\circ)$
as a direct summand for every $\tau^\circ\prec\Res^{L^F}_{L^{\circ F}}\tau$.
\end{corollary}

\subsubsection{Unipotent characters for disconnected groups}\label{subsubsec:unipotent-characters-for-disconnected}

\paragraph{Unipotent characters.}
Following {\cite[\S6.4]{Bonnafe1999Produit}}, an irreducible character $\gamma\in \Irr(G)$ is called
\emph{unipotent} if it occurs in some $R_{\mathbf T^\circ\subset \mathbf B^\circ}^{\mathbf G}(1)$,
where $\mathbf T^\circ\le \mathbf G^\circ$ is an $F$-stable maximal torus and $\mathbf B^\circ\le \mathbf G^\circ$
a Borel containing it.
Equivalently, $\gamma$ is unipotent iff it occurs in $\Ind_{G^\circ}^G(\gamma^\circ)$ for some unipotent
$\gamma^\circ\in \Irr(G^\circ)$ {\cite[Lem.~6.4.2]{Bonnafe1999Produit}}.

\begin{definition}[A semidirect-product model for $G^*$]\label{def:dual-semidirect-pinning}
Let $\bG$ be a (possibly disconnected) reductive algebraic group over
$\overline{\F}_p$ equipped with a Frobenius endomorphism $F:\bG\to\bG$.
Write $\bG^\circ$ for the identity component and
\[
\Omega \ :=\ \pi_0(\bG)\ \cong\ \bG/\bG^\circ .
\]
Fix an $F$-stable pinning
\[
\mathcal{P}=(\bB,\bT,\{x_\alpha\}_{\alpha\in\Delta})
\qquad\text{of }\bG^\circ.
\]
Let $\bG^{\circ*}$ be the Langlands dual complex reductive group of
$\bG^\circ$, equipped with the \emph{dual pinning}
$\mathcal{P}^*=(\bB^*,\bT^*,\{x_{\alpha^*}\}_{\alpha^*\in\Delta^*})$
and the dual Frobenius automorphism $F^*$.

For each $\omega\in\Omega$, let
\[
        \sigma_\omega\in\Aut(\bG^\circ,\mathcal P)
\]
be the pinned automorphism representing the algebraic outer action of $\omega$,
as in Definition~\ref{def:pinned-Omega-action}.  This semidirect-product
model records the pinned algebraic component action on the dual side.

By duality of pinned automorphisms, $\sigma_\omega$ corresponds uniquely to a
pinned automorphism
\[
\sigma_\omega^* \ \in\ \Aut(\bG^{\circ*},\mathcal{P}^*),
\]
and $\omega\mapsto \sigma_\omega^*$ defines a group homomorphism
\[
\iota:\Omega\longrightarrow \Aut(\bG^{\circ*}),\qquad
\omega\longmapsto \sigma_\omega^*.
\]

\smallskip
\noindent
\textbf{Definition of $\,\bG^*$:}
Define the \emph{(possibly disconnected) dual group} of $\bG$ to be the
semidirect product
\[
\bG^* \ :=\ \bG^{\circ*}\rtimes_{\iota}\Omega,
\]
i.e. as a variety $\bG^*=\bG^{\circ*}\times\Omega$ with multiplication
\[
(g_1,\omega_1)\cdot(g_2,\omega_2)
\ :=\
\bigl(g_1\,\iota(\omega_1)(g_2),\ \omega_1\omega_2\bigr),
\qquad
g_i\in\bG^{\circ*},\ \omega_i\in\Omega .
\]
Then $(\bG^*)^\circ=\bG^{\circ*}$ and $\pi_0(\bG^*)\cong\Omega$.

\smallskip
\noindent
\textbf{Frobenius on $\,\bG^*$:}
Let $F_\Omega:\Omega\to\Omega$ be the action induced by $F$ on components.
Assuming the pinning $\mathcal{P}$ is $F$-stable, define
\[
F^*(g,\omega)\ :=\ \bigl(F^*(g),\,F_\Omega(\omega)\bigr),
\qquad (g,\omega)\in \bG^{\circ*}\rtimes \Omega,
\]
so that $F^*$ is an algebraic automorphism of $\bG^*$ and restricts to the
usual dual Frobenius on $(\bG^*)^\circ=\bG^{\circ*}$.
\end{definition}
\medskip
\section{Disconnected Howlett--Lehrer theory}\label{sec:disconnected-Howlett-Lehrer-theory}
  \subsection{Clifford labels for component extensions}\label{subsec:clifford-labels}

Let $N\lhd M$ be finite with $M/N$ abelian and $\theta\in \Irr(N)$. Put
\[
I_M(\theta):=\{m\in M\mid {}^m\theta\simeq \theta\},\qquad \Omega_\theta:=I_M(\theta)/N.
\]

Clifford theory attaches to $\theta$ a canonical cohomology class
$[\alpha_\theta]\in H^2(\Omega_\theta,\C^\times)$.  We use the following
convention throughout the paper.  Choose a projective extension $\tilde\theta$
of $\theta$ to $I_M(\theta)$, and let
$c_\theta\in Z^2(\Omega_\theta,\C^\times)$ be its multiplier:
\[
\tilde\theta(i_1)\tilde\theta(i_2)
=
c_\theta(\bar i_1,\bar i_2)\tilde\theta(i_1i_2),
\qquad i_1,i_2\in I_M(\theta).
\]
The Clifford-label cocycle is defined to be
\[
\alpha_\theta:=c_\theta^{-1}.
\]
Thus $\alpha_\theta$ is the cocycle whose projective representations cancel the
multiplier of $\tilde\theta$.  Changing $\tilde\theta$ changes $c_\theta$, and
hence $\alpha_\theta$, by a coboundary; the cohomology class
$[\alpha_\theta]$ is canonical.  With this convention Clifford theory gives a
bijection
\[
\Irr(\C_{\alpha_\theta}[\Omega_\theta]) \xrightarrow{\ \sim\ } \Irr(M\mid \theta),
\]
depending on the chosen representative but only through the canonical class
$[\alpha_\theta]$.  Concretely,
\[
E\in \Irr(\C_{\alpha_\theta}[\Omega_\theta])
\quad\longmapsto\quad
\Ind_{I_M(\theta)}^{M}
\bigl(\tilde\theta\otimes \infl_{\Omega_\theta}^{I_M(\theta)}E\bigr)
\ \in \Irr(M\mid \theta).
\]
Here the tensor product is an honest representation of $I_M(\theta)$ because
$\tilde\theta$ has multiplier $c_\theta$ while $E$ has multiplier
$c_\theta^{-1}$.
\medskip
\subsection{A semisimple crossed-product description}
\label{subsec:semisimple-crossed-product}

The next lemma gives the corresponding semisimple crossed-product description at the
level of endomorphism algebras.

Let $\bX$ be a possibly disconnected reductive group over $\F_q$ with Frobenius
endomorphism $F$, and let $\bP$ be an $F$-stable parabolic subgroup of $\bX$ with
$F$-stable Levi factor $\bM$. Write
\[
X:=\bX^F,\qquad X^\circ:=(\bX^\circ)^F,\qquad
M:=\bM^F,\qquad M^\circ:=(\bM^\circ)^F,\qquad P^\circ:=(\bP^\circ)^F.
\]
Let $\theta^\circ\in \Irr(M^\circ)$ be cuspidal, and put
\[
\Omega_{\theta^\circ}:=I_M(\theta^\circ)/M^\circ.
\]
Choose a normalized $2$--cocycle
\[
c=c_{\theta^\circ}\in Z^2(\Omega_{\theta^\circ},\C^\times),
\]
a section
\[
\Omega_{\theta^\circ}\longrightarrow I_M(\theta^\circ),\qquad
\omega\longmapsto \dot\omega,\qquad \dot 1=1,
\]
and intertwiners
\[
\phi_\omega:\theta^\circ\xrightarrow{\ \sim\ }{}^{\dot\omega}\theta^\circ,
\qquad \phi_1=\id,
\]
realizing the factor set $c$.  In accordance with
\S\ref{subsec:clifford-labels}, the corresponding Clifford-label cocycle is
\[
\alpha=\alpha_{\theta^\circ}:=c^{-1}.
\]

Set
\[
M_{X^\circ}:=R^{X^\circ}_{M^\circ\subset P^\circ}(\theta^\circ),
\qquad
N_X:=R^X_{M\subset P}\!\bigl(\Ind_{M^\circ}^{M}\theta^\circ\bigr).
\]
Assume, as in the later applications, that the chosen parabolics are compatible so that
transitivity of Lusztig induction yields
\[
N_X
\;\cong\;
\Ind_{X^\circ}^{X}(M_{X^\circ}),
\]
and that the chosen $\Omega_{\theta^\circ}$--equivariant structure on $\theta^\circ$
induces an action of $\Omega_{\theta^\circ}$ on $M_{X^\circ}$ and identifies
$\Omega_{\theta^\circ}$ with the stabilizer of $M_{X^\circ}$ in $X/X^\circ$.

For each $\omega\in \Omega_{\theta^\circ}$, Harish--Chandra functoriality gives an
isomorphism
\[
\Phi_\omega:
M_{X^\circ}
=
R^{X^\circ}_{M^\circ\subset P^\circ}(\theta^\circ)
\xrightarrow{\ \sim\ }
R^{X^\circ}_{M^\circ\subset P^\circ}({}^{\dot\omega}\theta^\circ)
\cong
{}^{\dot\omega}M_{X^\circ}.
\]
These isomorphisms satisfy the same factor-set relation as the $\phi_\omega$.
They therefore define an action of $\Omega_{\theta^\circ}$ on
\[
E_{X^\circ}:=\End_{X^\circ}(M_{X^\circ})
\]
by
\[
\omega\cdot f
:=
\Phi_\omega^{-1}\circ {}^{\dot\omega}\!f\circ \Phi_\omega,
\qquad f\in E_{X^\circ}.
\]

\begin{lemma}\label{lem:semisimple-crossed-product}
With the above notation, there is a canonical algebra isomorphism
\[
\End_X(N_X)\xrightarrow{\sim} E_{X^\circ}\rtimes_c \Omega_{\theta^\circ},
\]
where the crossed product is formed with respect to the above action of
$\Omega_{\theta^\circ}$ on $E_{X^\circ}$ and the multiplier $c$ of the chosen
projective equivariant structure.
\end{lemma}

\begin{proof}
By the transitivity assumption,
\[
N_X\cong \Ind_{X^\circ}^{X}(M_{X^\circ}).
\]
Hence Frobenius reciprocity and Mackey theory for the normal subgroup
$X^\circ\lhd X$ give
\[
\End_X(N_X)
\cong
\bigoplus_{x\in X/X^\circ}
\Hom_{X^\circ}\!\bigl(M_{X^\circ},{}^xM_{X^\circ}\bigr).
\]
By the hypothesis that $\Omega_{\theta^\circ}$ identifies with the stabilizer of
$M_{X^\circ}$ in $X/X^\circ$, the summand indexed by $xX^\circ$ is nonzero if and only if
$xX^\circ\in \Omega_{\theta^\circ}$. Thus
\[
\End_X(N_X)
\cong
\bigoplus_{\omega\in \Omega_{\theta^\circ}}
\Hom_{X^\circ}\!\bigl(M_{X^\circ},{}^{\dot\omega}M_{X^\circ}\bigr).
\]

For each $\omega\in\Omega_{\theta^\circ}$, composition with $\Phi_\omega$ identifies
\[
E_{X^\circ}
=
\End_{X^\circ}(M_{X^\circ})
\xrightarrow{\sim}
\Hom_{X^\circ}\!\bigl(M_{X^\circ},{}^{\dot\omega}M_{X^\circ}\bigr),
\qquad
f\longmapsto \Phi_\omega\circ f.
\]
Under these identifications, $\End_X(N_X)$ becomes, as a vector space,
\[
\End_X(N_X)\cong \bigoplus_{\omega\in\Omega_{\theta^\circ}} E_{X^\circ}T_\omega,
\]
where $T_\omega$ denotes the basis element corresponding to $\Phi_\omega$.

Now let $f,g\in E_{X^\circ}$ and $\omega,\omega'\in\Omega_{\theta^\circ}$.
Using the definition of the action of $\Omega_{\theta^\circ}$ on $E_{X^\circ}$ and the
factor-set relation satisfied by the $\Phi_\omega$, one computes that
\[
(fT_\omega)(gT_{\omega'})
=
f\,(\omega\cdot g)\,c(\omega,\omega')\,T_{\omega\omega'}.
\]
This is exactly the defining multiplication in the crossed product
$E_{X^\circ}\rtimes_c \Omega_{\theta^\circ}$.
Therefore the assignment
\[
fT_\omega\longmapsto fT_\omega
\]
extends to an algebra isomorphism
\[
E_{X^\circ}\rtimes_c \Omega_{\theta^\circ}
\xrightarrow{\ \sim\ }
\End_X(N_X).
\]
This proves the lemma.
\end{proof}

The formulation is deliberately at the level of the full Harish--Chandra induced
module \(M_{X^\circ}\), rather than at the level of its irreducible constituents.
Thus no irreducibility assumption is imposed on
\[
M_{X^\circ}=R^{X^\circ}_{M^\circ\subset P^\circ}(\theta^\circ);
\]
the connected Howlett--Lehrer algebra is retained as the coefficient algebra in
the crossed product.
\medskip
\subsection{A corner lemma}
\label{subsec:corner-lemma}

Keep the notation of \S\ref{subsec:semisimple-crossed-product}.  In addition, put
\[
\cA_{\theta^\circ}:=\C_\alpha[\Omega_{\theta^\circ}],
\]
where \(\alpha=\alpha_{\theta^\circ}=c_{\theta^\circ}^{-1}\) is the
Clifford-label cocycle.  Thus \(\cA_{\theta^\circ}\) is the twisted group algebra
which parametrizes the irreducible representations of \(M\) lying above
\(\theta^\circ\).  Equivalently, it is the component algebra that appears after
passing from the \(c_{\theta^\circ}\)-twisted crossed product in
Lemma~\ref{lem:semisimple-crossed-product} to the opposite algebra.

Let \(A_{\bM}\) denote the maximal \(F\)-split torus in \(Z(\bM^\circ)\).  Thus
\[
\bM^\circ=C_{\bX^\circ}(A_{\bM}),
\]
and we write
\[
N_{X^\circ}(A_{\bM}):=N_{\bX^\circ}(A_{\bM})^F .
\]
Set, as before,
\[
N_X:=R^X_{M\subset P}\!\bigl(\Ind_{M^\circ}^{M}\theta^\circ\bigr),
\qquad
M_{X^\circ}:=R^{X^\circ}_{M^\circ\subset P^\circ}(\theta^\circ).
\]
Let
\[
W_{\theta^\circ}
:=
\Bigl\{
w\in N_{X^\circ}(A_{\bM})/M^\circ
\;\Bigm|\;
{}^{\,w}\theta^\circ\simeq \theta^\circ
\Bigr\}
\]
be the connected relative Weyl group attached to the cuspidal pair
\((M^\circ,\theta^\circ)\) inside \(X^\circ\).  By the connected
Howlett--Lehrer theorem, after choosing the usual normalized
Howlett--Lehrer intertwiners, we have an algebra isomorphism
\[
\End_{X^\circ}(M_{X^\circ})^{\op}\;\cong\;\C[W_{\theta^\circ}].
\]

For each
\[
E\in \Irr(\cA_{\theta^\circ}),
\]
let
\[
\sigma_E\in \Irr(M\mid \theta^\circ)
\]
be the irreducible representation corresponding to \(E\) under the
Clifford-theoretic bijection of \S\ref{subsec:clifford-labels}.  If
\(w\in W_{\theta^\circ}\), then \({}^{\,w}\sigma_E\) again lies above
\(\theta^\circ\), and hence there exists a unique
\[
w\cdot E\in \Irr(\cA_{\theta^\circ})
\]
such that
\[
{}^{\,w}\sigma_E \;\cong\; \sigma_{\,w\cdot E}.
\]
Thus \(W_{\theta^\circ}\) acts on \(\Irr(\cA_{\theta^\circ})\), and hence on the
set of primitive central idempotents of \(\cA_{\theta^\circ}\), by
\[
w\cdot e_E = e_{\,w\cdot E}.
\]
Put
\[
W_E:=\operatorname{Stab}_{W_{\theta^\circ}}(E)
=
\Bigl\{
w\in W_{\theta^\circ}\;\Bigm|\;{}^{\,w}\sigma_E\simeq \sigma_E
\Bigr\}.
\]
\smallskip
The following consequence of Lemma~\ref{lem:semisimple-crossed-product}
and the connected Howlett--Lehrer theorem will be used to pass from the
full induced object containing all Clifford extensions to the summand
attached to a fixed Clifford label.
\begin{lemma}\label{lem:corner-lemma}
With the above notation, there is a natural algebra isomorphism
\[
\End_X(N_X)^{\op}\xrightarrow{\sim}
\cA_{\theta^\circ}\rtimes W_{\theta^\circ},
\]
where the skew group algebra is formed with respect to the above action of
\(W_{\theta^\circ}\) on \(\cA_{\theta^\circ}\).

Let \(E\in \Irr(\cA_{\theta^\circ})\), and let
\(e_E\in Z(\cA_{\theta^\circ})\) be the corresponding primitive central
idempotent.  Then there exists a \(2\)-cocycle
\[
\beta_E\in Z^2(W_E,\C^\times)
\]
such that
\[
e_E\bigl(\cA_{\theta^\circ}\rtimes W_{\theta^\circ}\bigr)e_E
\;\cong\;
\End_\C(E^\vee)\otimes_\C \C_{\beta_E}[W_E].
\]
In particular, if the induced projective action of \(W_E\) on \(E^\vee\) lifts to
a genuine linear action, equivalently if \([\beta_E]=1\) in
\(H^2(W_E,\C^\times)\), then
\[
e_E\bigl(\cA_{\theta^\circ}\rtimes W_{\theta^\circ}\bigr)e_E
\;\cong\;
\End_\C(E^\vee)\otimes_\C \C[W_E].
\]
\end{lemma}

\begin{proof}
By Clifford theory and the convention \(\alpha=c_{\theta^\circ}^{-1}\),
\[
\Ind_{M^\circ}^{M}\theta^\circ
\;\cong\;
\bigoplus_{E\in \Irr(\cA_{\theta^\circ})}
\sigma_E\boxtimes E^\vee
\]
as an \(M\)-module with its natural right
\(\cA_{\theta^\circ}\)-action on the multiplicity factor.  The dual
multiplicity space \(E^\vee\) has multiplier \(c_{\theta^\circ}\), matching the
component intertwiners in Lemma~\ref{lem:semisimple-crossed-product}; this is
precisely why \(\cA_{\theta^\circ}\) appears after taking the opposite algebra.
Applying Harish--Chandra induction gives
\[
N_X
\;\cong\;
\bigoplus_{E\in \Irr(\cA_{\theta^\circ})}
R^X_{M\subset P}(\sigma_E)\boxtimes E^\vee
\]
as an \(X\)-module with the induced right \(\cA_{\theta^\circ}\)-action.

For each \(w\in W_{\theta^\circ}\), let
\[
B_w\in \End_{X^\circ}(M_{X^\circ})^{\op}
\]
denote the Howlett--Lehrer operator corresponding to \(w\).  Under the
identification
\[
\End_{X^\circ}(M_{X^\circ})^{\op}\cong \C[W_{\theta^\circ}],
\]
the operator \(B_w\) acts on the Clifford labels by
\[
B_w\, e_E \;=\; e_{\,w\cdot E}\, B_w.
\]
Equivalently, \(B_w\) sends the \(E\)-isotypic summand
\[
R^X_{M\subset P}(\sigma_E)\boxtimes E^\vee
\]
isomorphically onto the \((w\cdot E)\)-isotypic summand
\[
R^X_{M\subset P}(\sigma_{w\cdot E})\boxtimes (w\cdot E)^\vee.
\]

It follows that the subalgebra of \(\End_X(N_X)^{\op}\) generated by
\(\cA_{\theta^\circ}\) and the operators \(B_w\), \(w\in W_{\theta^\circ}\), is a
homomorphic image of the skew group algebra
\[
\cA_{\theta^\circ}\rtimes W_{\theta^\circ}.
\]
This homomorphism is an isomorphism.  Indeed, by
Lemma~\ref{lem:semisimple-crossed-product} and the connected
Howlett--Lehrer theorem,
\[
\dim_\C \End_X(N_X)
=
|\Omega_{\theta^\circ}|\cdot \dim_\C \End_{X^\circ}(M_{X^\circ})
=
|\Omega_{\theta^\circ}|\,|W_{\theta^\circ}|
=
\dim_\C\bigl(\cA_{\theta^\circ}\rtimes W_{\theta^\circ}\bigr).
\]
Thus the above surjective homomorphism between finite-dimensional algebras is
bijective, proving the first assertion.

We now compute the corner attached to a fixed
\(E\in \Irr(\cA_{\theta^\circ})\).  Inside the skew group algebra
\(\cA_{\theta^\circ}\rtimes W_{\theta^\circ}\), write \([w]\) for the standard
basis element corresponding to \(w\in W_{\theta^\circ}\).  The defining relation is
\[
[w]\, e_E = e_{\,w\cdot E}\,[w].
\]
Hence
\[
e_E [w] e_E=0
\qquad\text{if }w\notin W_E,
\]
and therefore
\[
e_E\bigl(\cA_{\theta^\circ}\rtimes W_{\theta^\circ}\bigr)e_E
=
\bigoplus_{w\in W_E} e_E\cA_{\theta^\circ}[w]e_E.
\]

The block algebra
\[
e_E\cA_{\theta^\circ}e_E
\]
is canonically isomorphic to \(\End_\C(E^\vee)\).  For \(w\in W_E\), the element
\([w]\) normalizes this full matrix algebra.  Since every automorphism of a full
matrix algebra is inner, we may choose elements
\[
u_w\in \bigl(e_E\cA_{\theta^\circ}e_E\bigr)^\times
\qquad (w\in W_E)
\]
such that
\[
[w]\,a\,[w]^{-1}=u_w a u_w^{-1}
\qquad
\bigl(a\in e_E\cA_{\theta^\circ}e_E\bigr).
\]
Set
\[
T_w:=u_w^{-1}e_E[w]e_E
\qquad (w\in W_E).
\]
Then each \(T_w\) commutes with \(e_E\cA_{\theta^\circ}e_E\).  Consequently the
products \(T_wT_{w'}\) are scalar multiples of \(T_{ww'}\), and there is a
uniquely determined function
\[
\beta_E:W_E\times W_E\longrightarrow \C^\times
\]
such that
\[
T_wT_{w'}=\beta_E(w,w')T_{ww'}.
\]
Associativity implies that \(\beta_E\) is a \(2\)-cocycle.  Thus
\[
e_E\bigl(\cA_{\theta^\circ}\rtimes W_{\theta^\circ}\bigr)e_E
\;\cong\;
e_E\cA_{\theta^\circ}e_E\otimes_\C \C_{\beta_E}[W_E]
\;\cong\;
\End_\C(E^\vee)\otimes_\C \C_{\beta_E}[W_E].
\]

If the projective action of \(W_E\) on \(E^\vee\) lifts to a genuine linear
action, then the elements \(T_w\) may be rescaled so that
\[
T_wT_{w'}=T_{ww'}.
\]
Equivalently, the cohomology class of \(\beta_E\) is trivial.  In this case the
corner is untwisted:
\[
e_E\bigl(\cA_{\theta^\circ}\rtimes W_{\theta^\circ}\bigr)e_E
\;\cong\;
\End_\C(E^\vee)\otimes_\C \C[W_E].
\]
\end{proof}

\begin{remark}\label{rmk:corner-cocycle}
The cocycle \(\beta_E\) measures only the possible projectivity of the stabilizer
action of \(W_E\) on the Clifford multiplicity space \(E^\vee\).  It is separate
from the Clifford-label cocycle \(\alpha\), which has already been absorbed into the
twisted group algebra \(\cA_{\theta^\circ}\).  In the abelian extension case used
below, where \(\theta^\circ\) extends to its inertia group, one has
\[
\cA_{\theta^\circ}=\C[\Omega_{\theta^\circ}],
\]
and the relevant Clifford labels are one-dimensional.  In that situation the
corner cocycle is trivial, and the corner is the ordinary group algebra of the
stabilizer.  Without this simplification, the same statements remain valid after
replacing \(\C[W_E]\) by \(\C_{\beta_E}[W_E]\).
\end{remark}
\medskip
\subsection{Disconnected Howlett--Lehrer theorem when the component group is abelian}
\label{subsec:disc-HL-abelian}

Let $\bX$ be a possibly disconnected reductive group over $\F_q$ with Frobenius
endomorphism $F$, and let $\bP$ be an $F$-stable parabolic subgroup of $\bX$ with
$F$-stable Levi factor $\bM$. Write
\[
X:=\bX^F,\qquad X^\circ:=(\bX^\circ)^F,\qquad
M:=\bM^F,\qquad M^\circ:=(\bM^\circ)^F,\qquad P^\circ:=(\bP^\circ)^F.
\]
Assume that \(X/X^\circ\) is abelian. Let
\[
\theta^\circ\in \Irr(M^\circ)
\]
be cuspidal, and put
\[
I:=I_M(\theta^\circ)
=
\{\,m\in M \mid {}^{\,m}\theta^\circ\simeq \theta^\circ\,\},
\qquad
\Omega_{\theta^\circ}:=I/M^\circ.
\]
Since \(X/X^\circ\) is abelian, the quotient \(\Omega_{\theta^\circ}\) is abelian.
Assume that \(\theta^\circ\) extends to \(I\), and fix such an extension
\[
\widetilde\theta\in \Irr(I).
\]
For \(\eta\in \Irr(\Omega_{\theta^\circ})\), inflated to \(I\), define
\[
\widetilde\theta_\eta:=\widetilde\theta\otimes \eta,
\qquad
\sigma_\eta:=\Ind_I^M(\widetilde\theta_\eta)\in \Irr(M).
\]
By Clifford theory, the characters \(\sigma_\eta\), for \(\eta\in \Irr(\Omega_{\theta^\circ})\),
are precisely the irreducible characters of \(M\) lying above \(\theta^\circ\), and
\[
\Ind_{M^\circ}^{M}\theta^\circ
\;\cong\;
\bigoplus_{\eta\in \Irr(\Omega_{\theta^\circ})}\sigma_\eta.
\]

Set
\[
M_{X^\circ}:=R^{X^\circ}_{M^\circ\subset P^\circ}(\theta^\circ),
\qquad
N_X:=R^X_{M\subset P}\!\bigl(\Ind_{M^\circ}^{M}\theta^\circ\bigr).
\]
Let
\[
W_{\theta^\circ}
:=
\Bigl\{
w\in N_{X^\circ}(A_{\bM})/M^\circ
\;\Bigm|\;
{}^{\,w}\theta^\circ\simeq \theta^\circ
\Bigr\}
\]
be the connected relative Weyl group attached to \((M^\circ,\theta^\circ)\).
As in \S\ref{subsec:corner-lemma}, \(W_{\theta^\circ}\) acts on
\(\Irr(\Omega_{\theta^\circ})\), characterized by
\[
{}^{\,w}\sigma_\eta \;\cong\; \sigma_{\,w\cdot\eta}.
\]
For \(\eta\in \Irr(\Omega_{\theta^\circ})\), put
\[
W_{\sigma_\eta}
:=
\Bigl\{
w\in W_{\theta^\circ}\;\Bigm|\;{}^{\,w}\sigma_\eta\simeq \sigma_\eta
\Bigr\}
=
\Stab_{W_{\theta^\circ}}(\eta).
\]
%
\begin{theorem}[Disconnected Howlett--Lehrer theorem in the abelian extension case]
\label{thm:disc-HL-abelian}
Assume that \(X/X^\circ\) is abelian and that
\(\theta^\circ\in \Irr(M^\circ)\) extends to its inertia group
\[
I=I_M(\theta^\circ).
\]
Fix an extension \(\widetilde\theta\in\Irr(I)\). For
\(\eta\in\Irr(\Omega_{\theta^\circ})\), where
\(\Omega_{\theta^\circ}=I/M^\circ\), set
\[
\widetilde\theta_\eta=\widetilde\theta\otimes\eta,
\qquad
\sigma_\eta=\Ind_I^M(\widetilde\theta_\eta).
\]
Then, for every \(\eta\in\Irr(\Omega_{\theta^\circ})\), there is an algebra
isomorphism
\[
\End_X\!\bigl(R^X_{M\subset P}(\sigma_\eta)\bigr)^{\op}
\cong
\C[W_{\sigma_\eta}],
\]
where
\[
W_{\sigma_\eta}
=
\{\,w\in W_{\theta^\circ}\mid {}^w\sigma_\eta\simeq\sigma_\eta\,\}.
\]
Consequently,
\[
\Irr\bigl(X,(M,\sigma_\eta)\bigr)
\xrightarrow{\sim}
\Irr(W_{\sigma_\eta}).
\]
Equivalently, for any \(\sigma\in\Irr(M\mid\theta^\circ)\), one has
\[
\End_X\!\bigl(R^X_{M\subset P}(\sigma)\bigr)^{\op}
\cong
\C[W_\sigma].
\]
\end{theorem}
\begin{proof}
Since \(\theta^\circ\) extends to \(I\), the Clifford-label cocycle attached to
\(\theta^\circ\) is trivial, and so the algebra \(\cA\) of
\S\ref{subsec:corner-lemma} is the ordinary group algebra
\[
\cA=\C[\Omega_{\theta^\circ}].
\]
By the decomposition of \(\Ind_{M^\circ}^{M}\theta^\circ\) above, we have
\[
N_X
=
R^X_{M\subset P}\!\bigl(\Ind_{M^\circ}^{M}\theta^\circ\bigr)
\;\cong\;
\bigoplus_{\eta\in \Irr(\Omega_{\theta^\circ})}
R^X_{M\subset P}(\sigma_\eta)
\]
as an \((X\times \cA)\)-module.

For \(\eta\in \Irr(\Omega_{\theta^\circ})\), let \(e_\eta\in Z(\cA)\) be the primitive
central idempotent corresponding to \(\eta\). Since \(\eta\) is one-dimensional, the
\(\eta\)-isotypic component of \(N_X\) is exactly
\[
e_\eta N_X \;\cong\; R^X_{M\subset P}(\sigma_\eta).
\]
Hence
\begin{equation}\label{eq:corner-end-disc-HL}
e_\eta \End_X(N_X)^{\op} e_\eta
\;\cong\;
\End_X\!\bigl(R^X_{M\subset P}(\sigma_\eta)\bigr)^{\op}.
\end{equation}

By Lemma~\ref{lem:corner-lemma}, we have
\[
\End_X(N_X)^{\op}\xrightarrow{\sim} \cA\rtimes W_{\theta^\circ},
\]
where the action of \(W_{\theta^\circ}\) on \(\cA=\C[\Omega_{\theta^\circ}]\) is the one
induced by its action on \(\Irr(\Omega_{\theta^\circ})\).
We now compute the corner \(e_\eta(\cA\rtimes W_{\theta^\circ})e_\eta\).

Let \([w]\) denote the standard basis element attached to \(w\in W_{\theta^\circ}\).
Since \(w\) acts on the primitive central idempotents by
\[
[w]\,e_\eta = e_{w\cdot \eta}\,[w],
\]
it follows that
\[
e_\eta [w] e_\eta = 0
\qquad\text{if } w\notin W_{\sigma_\eta}.
\]
On the other hand, if \(w\in W_{\sigma_\eta}\), then \(w\cdot\eta=\eta\), so
\[
[w]\,e_\eta=e_\eta\,[w]
\]
and therefore
\[
e_\eta [w] e_\eta = e_\eta [w].
\]
Hence
\[
e_\eta(\cA\rtimes W_{\theta^\circ})e_\eta
=
\bigoplus_{w\in W_{\sigma_\eta}} \C\, e_\eta [w].
\]
Moreover, for \(w,w'\in W_{\sigma_\eta}\),
\[
(e_\eta [w])(e_\eta [w'])
=
e_\eta [w] e_\eta [w']
=
e_\eta [w][w']
=
e_\eta [ww'].
\]
Thus the assignment
\[
T_w\longmapsto e_\eta [w]
\qquad (w\in W_{\sigma_\eta})
\]
extends to an algebra isomorphism
\begin{equation}\label{eq:corner-group-algebra-disc-HL}
\C[W_{\sigma_\eta}]
\xrightarrow{\ \sim\ }
e_\eta(\cA\rtimes W_{\theta^\circ})e_\eta.
\end{equation}

Combining \eqref{eq:corner-end-disc-HL} and \eqref{eq:corner-group-algebra-disc-HL},
we obtain
\[
\End_X\!\bigl(R^X_{M\subset P}(\sigma_\eta)\bigr)^{\op}
\;\cong\;
\C[W_{\sigma_\eta}],
\]
as required.

Finally, \(R^X_{M\subset P}(\sigma_\eta)\) is semisimple over \(\C\), so its irreducible
constituents are in canonical bijection with the simple right modules of its
endomorphism algebra. The above isomorphism therefore yields a canonical bijection
\[
\Irr\bigl(X,(M,\sigma_\eta)\bigr)\xrightarrow{\sim}\Irr(W_{\sigma_\eta}).
\]
\end{proof}

\begin{remark}\label{rmk:disc-HL-abelian-application}
Theorem~\ref{thm:disc-HL-abelian} is the form needed later in the construction of
Jordan decomposition for connected groups. In that application one takes
\[
X=H=C_{G^*}(s),\qquad M=H_L=C_{L^*}(s),
\]
and \(\theta^\circ=u_\tau^\circ\) is a cuspidal unipotent character of \(H_L^\circ\).
By the preferred-extension results proved earlier, \(u_\tau\) is of the form
\[
u_\tau = \Ind_{I_\tau}^{H_L}(\widetilde u_\tau),
\qquad
I_\tau=\Stab_{H_L}(u_\tau^\circ),
\]
with \(\widetilde u_\tau\) an extension of \(u_\tau^\circ\) to \(I_\tau\). Since
\(H/H^\circ\) is abelian by Lemma~\ref{lem:H-component-group}, the theorem applies.
\end{remark}
\medskip
\section{Construction of unique Jordan decomposition for connected groups}
\label{sec:construction-unique-JD-connected-groups}\medskip
\subsection{Weyl equivariance of the cuspidal correspondence}
\begin{proposition}[Weyl transport of the cuspidal Jordan decomposition]
\label{prop:weyl-equivariance}
Let $\bG$ be a connected reductive group defined over $\fq$, and let $\bL$ be an
$F$-stable Levi factor of an $F$-stable parabolic subgroup of $\bG$.  Let
\[
        \tau\in \cE(L,s)_{\cusp},
\]
and let \(w\in W_{\bG}(A_{\bL})\).  Denote by \(w^*\in
W_{\bG^*}(A_{\bL^*})\) the dual Weyl element.  Choose representatives, still
denoted \(w\) and \(w^*\), inducing dual conjugation automorphisms of \(\bL\)
and \(\bL^*\).  Put
\[
        s^w:=\ad(w^*)(s),
        \qquad
        \tau^w:=\tau\circ\ad(w).
\]
Then \(\tau^w\in \cE(L,s^w)_{\cusp}\), and
\begin{equation}\label{eq:weyl-transport-cuspidal-JD}
        \J_{s^w}^{L}(\tau^w)
        =
        \J_s^{L}(\tau)\circ\ad((w^*)^{-1}).
\end{equation}
Here the right hand side is viewed as a unipotent character of
\(C_{\bL^*}(s^w)^{F^*}\) through the isomorphism
\[
        \ad((w^*)^{-1}):C_{\bL^*}(s^w)^{F^*}
        \xrightarrow{\sim}
        C_{\bL^*}(s)^{F^*}.
\]
In particular, if \(w\) stabilizes the isomorphism class of \(\tau\), then
\(s^w\) is \(L^*\)-conjugate to \(s\), and after the usual identification of
\(C_{\bL^*}(s^w)\) with \(C_{\bL^*}(s)\), the formula becomes the expected
stabilizer-equivariance statement.
\end{proposition}

\begin{proof}
The assertion that \(\tau^w\in\cE(L,s^w)\) is the standard compatibility of
Lusztig series with dual automorphisms.  The equality
\eqref{eq:weyl-transport-cuspidal-JD} follows from the way the cuspidal map
\(\J_s^L\) was constructed.  The map \(\ad(w)\) transports the source
diagonal-automorphism orbit and its Clifford data to those attached to
\(\tau^w\), while \(\ad(w^*)\) transports the connected centralizer
\(C_{\bL^*}(s)\) to \(C_{\bL^*}(s^w)\).  The homomorphism
\(L_{\ad}/L\to (H_L/H_L^\circ)^\vee\) used to match the two Clifford pictures
commutes with these Weyl actions, as recalled in
Remark~\ref{rmk:regular-embedding-adjoint-action}.  Finally, the
preferred-extension normalization is stable under transport by pinned
automorphisms: Lusztig's preferred extension is defined from the based root
datum, and Malle's matching is used with this transported normalization.
Therefore the extension, Clifford label, and cuspidal unipotent character
obtained from \(\tau^w\) are precisely the \(w^*\)-transports of those obtained
from \(\tau\), which is exactly \eqref{eq:weyl-transport-cuspidal-JD}.
\end{proof}

\begin{corollary}\label{coro:isoweylgp}
Let \(\tau\in \cE(L,s)_{\cusp}\), and let
\[
        u_\tau:=\J_s^L(\tau)\in \Uch(H_L)_{\cusp}.
\]
Then the duality isomorphism identifies the stabilizer of the cuspidal pair
\((L,\tau)\) with the stabilizer of the corresponding unipotent cuspidal pair
\((H_L,u_\tau)\):
\[
        W_{\tau}\cong W_{u_\tau}.
\]
\end{corollary}

\begin{proof}
Let \(w\in W_\tau\).  Then \(\tau^w\simeq\tau\), so the preceding proposition
implies that \(s^w\) is \(L^*\)-conjugate to \(s\).  Replacing the representative
\(w^*\) by the corresponding element in \(W_{\bH}(A_{\bL^*})\), where
\(\bH=C_{\bG^*}(s)\), the transport formula
\eqref{eq:weyl-transport-cuspidal-JD} gives
\[
        u_\tau\circ\ad((w^*)^{-1})\simeq u_\tau.
\]
Thus \(w^*\in W_{u_\tau}\).  The same argument applied to the inverse Weyl
element gives the converse inclusion.  Hence duality identifies the two
stabilizers.
\end{proof}
\subsection{Passage from cuspidal pairs to Harish--Chandra series}
\label{subsec:passage-cusp-to-HC}

Fix a semisimple element \(s\in G^*\), and put
\[
\bH:=C_{\bG^*}(s),\qquad H:=\bH^{F^*}.
\]
For every \(F\)-stable Levi subgroup \(\bL\leq \bG\) and every
\[
\tau\in \cE(L,s)_{\cusp},
\]
write
\[
\bH_L:=C_{\bL^*}(s),\qquad H_L:=\bH_L^{F^*},\qquad
u_\tau:=\J_s^L(\tau)\in \Uch(H_L)_{\cusp},
\]
where \(\J_s^L\) is the cuspidal Jordan decomposition for the connected group \(\bL\)
constructed in Lemma~\ref{lemma:bijection-cuspidal}.

\medskip

\begin{proposition}\label{prop:passage-cusp-to-HC}
Let \(\bL\) be an \(F\)-stable Levi factor of an \(F\)-stable parabolic subgroup of
\(\bG\), and let
\[
\tau\in \cE(L,s)_{\cusp}.
\]
Then there is a canonically defined bijection
\[
\J_{s,(L,\tau)}:\Irr\bigl(G,(L,\tau)\bigr)\xrightarrow{\ \sim\ }
\Irr\bigl(H,(H_L,u_\tau)\bigr).
\]

More precisely, let
\[
\Irr(W_\tau)=\{\phi\}
\]
and let
\[
\rho_\phi\in \Irr\bigl(G,(L,\tau)\bigr)
\]
be the constituent corresponding to \(\phi\) under the connected
Howlett--Lehrer comparison theorem. By Corollary~\ref{coro:isoweylgp}, we have a
canonical identification
\[
W_\tau \xrightarrow{\sim} W_{u_\tau}.
\]
Let
\[
u_\phi\in \Irr\bigl(H,(H_L,u_\tau)\bigr)
\]
be the constituent corresponding to the same \(\phi\), now viewed as an irreducible
character of \(W_{u_\tau}\) via the above identification and
Theorem~\ref{thm:disc-HL-abelian}. Then
\[
\J_{s,(L,\tau)}(\rho_\phi):=u_\phi
\]
defines the required bijection.

Furthermore, if \(\bT^*\subseteq \bL^*\) is an \(F^*\)-stable maximal torus containing
\(s\), then
\[
\bigl\langle R_{\bT^*}^{\bG}(s),\rho_\phi\bigr\rangle_G
=
\epsilon_{\bG}\epsilon_{\bH}\,
\bigl\langle R_{\bT^*}^{\bH}(1),u_\phi\bigr\rangle_H.
\]
\end{proposition}

\begin{proof}
By the connected Howlett--Lehrer theorem
(Theorem~\ref{thm:Howlett-Lehrer}), there is an algebra isomorphism
\[
\End_G\!\bigl(R_{\bL}^{\bG}(\tau)\bigr)^{\op}\xrightarrow{\sim}\C[W_\tau],
\]
hence a canonical bijection
\[
\Irr\bigl(G,(L,\tau)\bigr)\xrightarrow{\sim}\Irr(W_\tau),
\qquad
\rho_\phi \longleftrightarrow \phi.
\]
Likewise, by
Theorem~\ref{thm:disc-HL-abelian}, there is an algebra isomorphism
\[
\End_H\!\bigl(R_{\bH_L}^{\bH}(u_\tau)\bigr)^{\op}\xrightarrow{\sim}\C[W_{u_\tau}],
\]
hence a canonical bijection
\[
\Irr\bigl(H,(H_L,u_\tau)\bigr)\xrightarrow{\sim}\Irr(W_{u_\tau}),
\qquad
u_\phi \longleftrightarrow \phi.
\]
Using Corollary~\ref{coro:isoweylgp} to identify \(W_\tau\) with \(W_{u_\tau}\), we
obtain the bijection \(\J_{s,(L,\tau)}\).

We now prove the scalar-product identity. Since \(\rho_\phi\) corresponds to \(\phi\),
the regular-module decomposition of
\(\End_G(R_{\bL}^{\bG}(\tau))\cong \C[W_\tau]\) yields
\[
R_{\bL}^{\bG}(\tau)
=
\sum_{\phi\in \Irr(W_\tau)} \phi(1)\rho_\phi.
\]
By Frobenius reciprocity, this implies
\[
\bigl\langle {}^*R_{\bL}^{\bG}(\rho_\phi),\tau\bigr\rangle_L
=
\phi(1).
\]
Since \(\tau\) is cuspidal, the Harish--Chandra series of \(L\) attached to
\((L,\tau)\) is the singleton \(\{\tau\}\). Therefore
\[
{}^*R_{\bL}^{\bG}(\rho_\phi)=\phi(1)\tau.
\]
Using transitivity of Lusztig induction and Frobenius reciprocity, we obtain
\[
\bigl\langle R_{\bT^*}^{\bG}(s),\rho_\phi\bigr\rangle_G
=
\bigl\langle R_{\bT^*}^{\bL}(s),{}^*R_{\bL}^{\bG}(\rho_\phi)\bigr\rangle_L
=
\phi(1)\bigl\langle R_{\bT^*}^{\bL}(s),\tau\bigr\rangle_L.
\]

On the \(H\)-side, the same argument gives
\[
R_{\bH_L}^{\bH}(u_\tau)
=
\sum_{\phi\in \Irr(W_{u_\tau})} \phi(1)u_\phi,
\]
hence
\[
{}^*R_{\bH_L}^{\bH}(u_\phi)=\phi(1)u_\tau,
\]
because \(u_\tau\) is cuspidal in \(H_L\). Therefore
\[
\bigl\langle R_{\bT^*}^{\bH}(1),u_\phi\bigr\rangle_H
=
\bigl\langle R_{\bT^*}^{\bH_L}(1),{}^*R_{\bH_L}^{\bH}(u_\phi)\bigr\rangle_{H_L}
=
\phi(1)\bigl\langle R_{\bT^*}^{\bH_L}(1),u_\tau\bigr\rangle_{H_L}.
\]

By Lemma~\ref{lemma:bijection-cuspidal}, applied to the connected Levi subgroup
\(\bL\), we have
\[
\bigl\langle R_{\bT^*}^{\bL}(s),\tau\bigr\rangle_L
=
\epsilon_{\bL}\epsilon_{\bH_L}\,
\bigl\langle R_{\bT^*}^{\bH_L}(1),u_\tau\bigr\rangle_{H_L}.
\]
Since \(\bH_L^\circ\) is the Levi subgroup of \(\bH^\circ\) corresponding to
\(\bL\subseteq \bG\), one has
\[
\epsilon_{\bL}\epsilon_{\bH_L}=\epsilon_{\bG}\epsilon_{\bH}.
\]
Combining the preceding equalities yields
\[
\bigl\langle R_{\bT^*}^{\bG}(s),\rho_\phi\bigr\rangle_G
=
\epsilon_{\bG}\epsilon_{\bH}\,
\bigl\langle R_{\bT^*}^{\bH}(1),u_\phi\bigr\rangle_H.
\]
This proves the proposition.
\end{proof}

\subsection{Construction of $\J_s$}
\label{subsec:construction-Js}

\begin{theorem}\label{thm:JD-connected-groups}
Let \(\bG\) be a connected reductive group defined over \(\F_q\), equipped with a
pinning \(\cP\). Let \(s\in G^*\) be semisimple, and put
\[
\bH:=C_{\bG^*}(s),\qquad H:=\bH^{F^*}.
\]
Then there exists a uniquely determined bijection
\[
\J_s:\cE(G,s)\xrightarrow{\ \sim\ }\Uch(H)
\]
such that the following conditions hold.

\begin{enumerate}[label={\normalfont(\arabic*)}]
\item The restriction of \(\J_s\) to \(\cE(G,s)_{\cusp}\) is the cuspidal bijection
\[
\J_s^G:\cE(G,s)_{\cusp}\xrightarrow{\sim}\Uch(H)_{\cusp}
\]
constructed in Lemma~\ref{lemma:bijection-cuspidal}.

\item For every \(\rho\in \cE(G,s)\) and every \(F^*\)-stable maximal torus
\(\bT^*\subseteq \bG^*\) containing \(s\), one has
\[
\bigl\langle R_{\bT^*}^{\bG}(s),\rho\bigr\rangle_G
=
\epsilon_{\bG}\epsilon_{\bH}\,
\bigl\langle R_{\bT^*}^{\bH}(1),\J_s(\rho)\bigr\rangle_H.
\]

\item If
\[
\cE(G,s)=\bigsqcup_{[(L,\tau)]\in \Sigma_G(s)} \Irr\bigl(G,(L,\tau)\bigr)
\]
is the Harish--Chandra decomposition of the Lusztig series, then
\(\J_s\) restricts on each Harish--Chandra series to the bijection
\[
\J_{s,(L,\tau)}:\Irr\bigl(G,(L,\tau)\bigr)\xrightarrow{\sim}
\Irr\bigl(H,(H_L,u_\tau)\bigr)
\]
of Proposition~\ref{prop:passage-cusp-to-HC}, where
\[
u_\tau:=\J_s^L(\tau)\in \Uch(H_L)_{\cusp}.
\]
\end{enumerate}
\end{theorem}

\begin{proof}
By Proposition~\ref{prop:Lusztig-series-union-HC}, we have a disjoint decomposition
\[
\cE(G,s)
=
\bigsqcup_{[(L,\tau)]\in \Sigma_G(s)}
\Irr\bigl(G,(L,\tau)\bigr).
\]
On the \(H\)-side, Harish--Chandra theory for the possibly disconnected group \(H\)
gives a disjoint decomposition
\[
\Uch(H)
=
\bigsqcup_{[(M,u)]}
\Irr\bigl(H,(M,u)\bigr),
\]
where \((M,u)\) runs over the \(H\)-conjugacy classes of cuspidal unipotent pairs in
\(H\).

By Lemma~\ref{lemma:bijection-cuspidal}, applied to all connected Levi subgroups of
\(\bG\), together with Proposition~\ref{prop:weyl-equivariance}, the assignment
\[
(L,\tau)\longmapsto \bigl(H_L,\J_s^L(\tau)\bigr)
\]
induces a bijection from \(\Sigma_G(s)\) onto the set of \(H\)-conjugacy classes of
cuspidal unipotent pairs in \(H\). Therefore, for each
\([(L,\tau)]\in \Sigma_G(s)\), Proposition~\ref{prop:passage-cusp-to-HC} provides a
canonical bijection
\[
\J_{s,(L,\tau)}:\Irr\bigl(G,(L,\tau)\bigr)\xrightarrow{\sim}
\Irr\bigl(H,(H_L,u_\tau)\bigr).
\]
Taking the disjoint union of these bijections over all
\([(L,\tau)]\in \Sigma_G(s)\), we obtain a well-defined bijection
\[
\J_s:\cE(G,s)\xrightarrow{\sim}\Uch(H).
\]

Property \((3)\) is built into the definition. Property \((1)\) follows because when
\(L=G\), the Harish--Chandra series \(\Irr(G,(G,\rho))\) is the singleton
\(\{\rho\}\), and the above construction recovers the cuspidal bijection
\(\J_s^G\). Finally, Property \((2)\) is exactly the scalar-product identity proved in
Proposition~\ref{prop:passage-cusp-to-HC} on each Harish--Chandra series.

The uniqueness of \(\J_s\) is immediate from the disjoint decomposition of
\(\cE(G,s)\) into Harish--Chandra series and the fact that its restriction to each
series is prescribed by Proposition~\ref{prop:passage-cusp-to-HC}.
\end{proof}

\section{Cuspidal Jordan decomposition for disconnected $\bG$}\label{sec:cuspidal-JD-disconnected}
The preceding sections establish the connected Jordan decomposition.  We now
pass to the Clifford-theoretic extension needed when the group whose characters
are being decomposed is itself disconnected.  From this point onward the
following rational condition is imposed on the finite rational components of
that disconnected group.

\begin{hypothesis}[Rational pinned-component condition]
\label{hyp:rational-pinned-component-condition}
Let $\mathbf G$ be a possibly disconnected reductive group over $\F_q$, with
identity component $\mathbf G^\circ$, and put
\[
        G:=\mathbf G^F,\qquad G^\circ:=\mathbf G^{\circ F},\qquad
        \Omega_G:=G/G^\circ.
\]
Fix an $F$-stable pinning $\cP$ of $\mathbf G^\circ$.  Define
\[
        G_{\cP}
        :=
        \{g\in G:\Ad(g)|_{\mathbf G^\circ}
        \text{ preserves }\cP\}.
\]
We assume that the natural map
\[
        G_{\cP}\longrightarrow \Omega_G
\]
is surjective.  Equivalently,
\[
        G=G^\circ\cdot G_{\cP}.
\]
Whenever the disconnected Jordan decomposition is applied to an $F$-stable
regular Levi subgroup $\mathbf L\subset \mathbf G$, we impose the same condition
on $L:=\mathbf L^F$ with the pinning induced on $\mathbf L^\circ$.
\end{hypothesis}

\begin{remark}[Scope of the disconnected theory]
\label{rmk:scope-rational-pinned-components}
Hypothesis~\ref{hyp:rational-pinned-component-condition} is imposed only in the
disconnected part of the paper, where the group whose characters are being
decomposed may be disconnected.  It is not a hypothesis in the connected Jordan
decomposition theorem: there the character side is the finite group of points of
a connected reductive group, while the possibly disconnected group
$C_{\mathbf G^*}(s)^{F^*}$ occurs on the unipotent side.

The disconnected results below are not intended as a Jordan decomposition for
every disconnected reductive finite group with abelian component group.  The
extra rational pinned-component condition is imposed because the primary
application is to depth-zero local Langlands constructions (see \cite{Mishra2026PinnedLLCDepthZero}), where the relevant
finite groups are full reductive quotients of parahoric stabilizers.  In that
setting the components come from vertex stabilizers, and the required input is
that their action on the connected reductive quotient is pinning-preserving
modulo inner conjugacy by the connected finite quotient.  Thus the disconnected
theory used here is tailored to the finite groups that occur in the local
Langlands application, and it excludes the diagonal-automorphism pathology in
which an algebraically inner, but not rationally inner, automorphism acts
nontrivially on characters of the finite identity component.
\end{remark}

Throughout the rest of this section, \(\bG\) is a possibly disconnected reductive
group over \(\F_q\), equipped with an \(F\)-stable pinning \(\cP\) of
\(\bG^\circ\).  Set \(G:=\bG^F\), \(G^\circ:=\bG^{\circ F}\), and
\(\Omega:=G/G^\circ\).  We assume that \(\Omega\) is abelian and that
\((\bG,\cP)\) satisfies Hypothesis~\ref{hyp:rational-pinned-component-condition}.

Fix \(s\in \bG^{\circ*\,F^*}=(\bG^*)^{\circ F^*}\) semisimple and put
\[
\bH:=C_{\bG^*}(s),\qquad
\bH_0:=C_{\bG^{\circ*}}(s)=C_{(\bG^*)^\circ}(s),\qquad
\bH^\circ:=\bH_0^\circ=C_{\bG^{\circ*}}(s)^\circ.
\]
Set
\[
H:=\bH^{F^*},\qquad
H_0:=\bH_0^{F^*},\qquad
H^\circ:=(\bH^\circ)^{F^*}.
\]
Thus \(H_0\) is, in general, disconnected, while \(H^\circ\) is its identity-component finite subgroup.

For each \(\omega\in\Omega\), let
\[
\sigma_\omega\in \Aut(\bG^\circ,\cP)
\]
be the pinned representative of the outer action from Definition~\ref{def:pinned-Omega-action},
and let
\[
\sigma_\omega^*\in \Aut(\bG^{\circ*},\cP^*)
\]
be the dual pinned automorphism. We shall use these pinned automorphisms below only up to rational conjugacy of
semisimple parameters.  Thus, for a semisimple element
\(s\in \bG^{\circ *F^*}\), the relevant component group is the stabilizer of
the \(\bG^{\circ *F^*}\)-conjugacy class of \(s\), not necessarily the
pointwise stabilizer of \(s\) under the automorphisms \(\sigma_\omega^*\).


\subsubsection{Stabilizer matching and the source Clifford class}
\label{subsubsec:stabilizer-matching-source-clifford}

Keep the notation fixed at the beginning of this section:
\[
\bH=C_{\bG^*}(s),\qquad
\bH_0=C_{\bG^{\circ *}}(s)=C_{(\bG^*)^\circ}(s),
\qquad
H=\bH^{F^*},\qquad H_0=\bH_0^{F^*}.
\]
Thus \(H_0\lhd H\).  The point to keep in mind is that
\(\bH\) is not, in general, the semidirect product
\(\bH_0\rtimes\Omega_s\), where
\(\Omega_s=\{\omega:\sigma_\omega^*(s)=s\}\).  Rather, the relevant
component group is the stabilizer of the finite rational conjugacy class of
\(s\).

Put
\[
\Omega_{[s]}
:=
\left\{
\omega\in\Omega:
\sigma_\omega^*(s)
\text{ is } \bG^{\circ *F^*}\text{-conjugate to }s
\right\}.
\]
For \(\omega\in\Omega_{[s]}\), choose
\[
x_\omega\in \bG^{\circ *F^*},
\qquad x_1=1,
\]
such that
\begin{equation}\label{eq:xomega-conjugates-s}
x_\omega\,\sigma_\omega^*(s)\,x_\omega^{-1}=s.
\end{equation}
Then
\[
h_\omega:=(x_\omega,\omega)\in
\bG^{\circ *F^*}\rtimes\Omega
\]
belongs to \(H\).  Moreover, the projection \(H\to\Omega\) has image
\(\Omega_{[s]}\), and one has an exact sequence
\begin{equation}\label{eq:H-H0-Omega-class}
1\longrightarrow H_0
\longrightarrow H
\longrightarrow \Omega_{[s]}
\longrightarrow 1.
\end{equation}
For \(\omega,\eta\in\Omega_{[s]}\), the chosen section
\(\omega\mapsto h_\omega\) has factor set
\begin{equation}\label{eq:centralizer-factor-set}
c_s(\omega,\eta)
:=
h_\omega h_\eta h_{\omega\eta}^{-1}
=
x_\omega\,\sigma_\omega^*(x_\eta)\,x_{\omega\eta}^{-1}
\in H_0 .
\end{equation}

Let
\[
\rho^\circ\in \cE(G^\circ,s)_{\cusp}
\]
and set
\[
u_0:=\J_s^{G^\circ}(\rho^\circ)\in \Uch(H_0)_{\cusp},
\]
where \(\J_s^{G^\circ}\) denotes the pinned cuspidal Jordan decomposition for
the connected group \(\bG^\circ\).  Define
\[
\Omega_{\rho^\circ}
:=
\operatorname{im}\bigl(I_G(\rho^\circ)\to G/G^\circ=\Omega\bigr),
\]
and
\[
\Omega_{u_0}
:=
\operatorname{im}\bigl(I_H(u_0)\to H/H_0\hookrightarrow\Omega\bigr).
\]

\begin{lemma}[Stabilizer matching]
\label{lem:disc-stabilizer-match}
With the notation above, one has
\[
\Omega_{\rho^\circ}=\Omega_{u_0}\subseteq \Omega_{[s]}.
\]
Denote this common subgroup by \(\Omega_0\).  Via the projection to
\(\Omega_0\), there are canonical identifications
\[
I_G(\rho^\circ)/G^\circ\cong \Omega_0,
\qquad
I_H(u_0)/H_0\cong \Omega_0.
\]
In particular, the Clifford class attached to the source normal inclusion
\[
        G^\circ\lhd I_G(\rho^\circ)
\]
may be regarded as a canonical class
\[
        [\alpha_{\rho^\circ}]\in H^2(\Omega_0,\C^\times).
\]
\end{lemma}

\begin{proof}
We first record the centralizer calculation.  Since
\[
\bG^*=\bG^{\circ *}\rtimes\Omega,
\]
an element \((g,\omega)\in \bG^{\circ *}\rtimes\Omega\) centralizes \(s\) if
and only if
\[
(g,\omega)(s,1)(g,\omega)^{-1}
=
(g\,\sigma_\omega^*(s)\,g^{-1},1)
=
(s,1).
\]
Equivalently,
\[
g\,\sigma_\omega^*(s)\,g^{-1}=s.
\]
Thus \(\omega\) need not fix \(s\) itself; it is enough and necessary that
\(\sigma_\omega^*(s)\) be \(\bG^{\circ *}\)-conjugate to \(s\).  On rational
points this gives \eqref{eq:H-H0-Omega-class}.

We next compare the two stabilizers.  Let \(\omega\in\Omega_{\rho^\circ}\).
Choose first any representative of \(\omega\) in \(I_G(\rho^\circ)\).  By
Hypothesis~\ref{hyp:rational-pinned-component-condition} we may replace it,
modulo left multiplication by an element of \(G^\circ\), by a representative
\[
        \dot\omega\in I_G(\rho^\circ)\cap G_{\cP}.
\]
Inner conjugation by \(G^\circ\) is invisible on characters of \(G^\circ\).  For
this representative, conjugation on \(G^\circ\) is the pinned automorphism
\(\sigma_\omega\).  Since pullback by \(\sigma_\omega^{-1}\) sends
\[
\cE(G^\circ,s)
\quad\text{onto}\quad
\cE(G^\circ,\sigma_\omega^*(s)),
\]
and since \({}^{\dot\omega}\rho^\circ\simeq\rho^\circ\), the two connected
Lusztig series
\[
\cE(G^\circ,s)
\quad\text{and}\quad
\cE(G^\circ,\sigma_\omega^*(s))
\]
coincide.  Hence \(s\) and \(\sigma_\omega^*(s)\) are
\(\bG^{\circ *F^*}\)-conjugate.  Thus \(\omega\in\Omega_{[s]}\), and we may
choose \(x_\omega\in\bG^{\circ *F^*}\) satisfying
\eqref{eq:xomega-conjugates-s}.  Then
\[
h_\omega=(x_\omega,\omega)\in H.
\]

The connected pinned Jordan decomposition is equivariant for the following
transport operation: first apply the pinned automorphism \(\sigma_\omega\) on
the \(G^\circ\)-side, which sends \(s\) to \(\sigma_\omega^*(s)\), and then
conjugate on the dual side by \(x_\omega\), which brings
\(\sigma_\omega^*(s)\) back to \(s\).  Therefore, for every
\(\chi\in\cE(G^\circ,s)_{\cusp}\),
\begin{equation}\label{eq:JD-equivariance-with-xomega}
\J_s^{G^\circ}({}^{\dot\omega}\chi)
=
{}^{h_\omega}\J_s^{G^\circ}(\chi).
\end{equation}
Applying \eqref{eq:JD-equivariance-with-xomega} to \(\chi=\rho^\circ\), and
using \({}^{\dot\omega}\rho^\circ\simeq\rho^\circ\), gives
\[
{}^{h_\omega}u_0
=
{}^{h_\omega}\J_s^{G^\circ}(\rho^\circ)
=
\J_s^{G^\circ}({}^{\dot\omega}\rho^\circ)
\simeq
\J_s^{G^\circ}(\rho^\circ)
=
u_0.
\]
Thus \(h_\omega\in I_H(u_0)\), and hence
\[
\Omega_{\rho^\circ}\subseteq\Omega_{u_0}.
\]

Conversely, let \(\omega\in\Omega_{u_0}\), and choose
\(h=(x,\omega)\in I_H(u_0)\).  Since \(h\in H=C_{\bG^*}(s)^{F^*}\), we have
\[
x\,\sigma_\omega^*(s)\,x^{-1}=s.
\]
By Hypothesis~\ref{hyp:rational-pinned-component-condition}, choose a
pinning-preserving representative \(\dot\omega\in G_{\cP}\) of \(\omega\).
The same equivariance argument gives
\[
\J_s^{G^\circ}({}^{\dot\omega}\rho^\circ)
=
{}^{h}\J_s^{G^\circ}(\rho^\circ)
=
{}^h u_0
\simeq
u_0
=
\J_s^{G^\circ}(\rho^\circ).
\]
Since \(\J_s^{G^\circ}\) is injective on \(\cE(G^\circ,s)_{\cusp}\), it follows
that
\[
{}^{\dot\omega}\rho^\circ\simeq\rho^\circ.
\]
Hence \(\omega\in\Omega_{\rho^\circ}\).  Therefore
\[
\Omega_{\rho^\circ}=\Omega_{u_0}=:\Omega_0.
\]

The projection maps now give the announced identifications
\[
I_G(\rho^\circ)/G^\circ\cong\Omega_0,
\qquad
I_H(u_0)/H_0\cong\Omega_0,
\]
because their kernels are respectively \(G^\circ\) and \(H_0\).  Clifford
theory for the normal inclusion \(G^\circ\lhd I_G(\rho^\circ)\) therefore gives
the stated canonical cohomology class on \(\Omega_0\).
\end{proof}

\begin{definition}[Projective labels with prescribed multiplier]
\label{def:projective-labels-fixed-multiplier}
Let \(A\) be a finite group and let \([\alpha]\in H^2(A,\C^\times)\).  We write
\[
        \Irr_{[\alpha]}(A)
\]
for the set of irreducible projective representations of \(A\) with multiplier
in the class \([\alpha]\).  Equivalently, after choosing a normalized cocycle
representative \(\alpha\in Z^2(A,\C^\times)\), it is the set
\(\Irr(\C_\alpha[A])\), with the usual identification under replacing
\(\alpha\) by a cohomologous cocycle.
\end{definition}

\begin{definition}[Cuspidal enriched unipotent target]
\label{def:enh-cusp-target}
With the above notation, define
\[
\Uch^{\mathrm{enh}}_{G,s}(H)_{\cusp}
:=
\left\{
(u_0,[\alpha_{\rho^\circ}],E)
:
\begin{array}{l}
u_0=\J_s^{G^\circ}(\rho^\circ),\quad
\rho^\circ\in\cE(G^\circ,s)_{\cusp},\\[2mm]
E\in \Irr_{[\alpha_{\rho^\circ}]}(\Omega_{\rho^\circ})
\end{array}
\right\}/\Omega_{[s]}.
\]
Here \(\Omega_{[s]}\) acts by conjugating \(\rho^\circ\), hence also the
corresponding unipotent character \(u_0\), the stabilizer
\(\Omega_{\rho^\circ}\), the Clifford class \([\alpha_{\rho^\circ}]\), and the
projective label \(E\).  Thus the notation records precisely the source
Clifford class and source Clifford label, transported to the unipotent side as
part of the enhanced datum.
\end{definition}
\subsection{The cuspidal disconnected JD}
\subsubsection{Lusztig series for $G^F$ attached to $s\in (G^*)^{\circ F^*}$}
\label{subsubsec:disc-Lusztig-series}

We keep the standing assumptions and notation from \S9: $G$ is a (possibly disconnected) reductive
$\F_q$--group with Frobenius $F$, $G^\circ\lhd G$ is the identity component, and $\Omega:=G^F/G^{\circ F}$
is abelian. Fix an $F$--stable pinning $\mathbf P$ of $G^\circ$, and assume the rational
pinned-component condition of Hypothesis~\ref{hyp:rational-pinned-component-condition}.  As in Definition~\ref{def:dual-semidirect-pinning}, this yields a semidirect-product model
\[
G^* \;=\; G^{\circ *}\rtimes \Omega,
\]
together with pinned automorphisms $\sigma_\omega\in \Aut(G^\circ,\mathbf P)$ commuting with $F$ and their
duals $\sigma_\omega^*\in \Aut(G^{\circ *})$ for $\omega\in \Omega$.

\begin{definition}[Disconnected Lusztig series in the identity-component case]
\label{def:disc-Lusztig-series}
Let $s\in (G^*)^{\circ F^*}=G^{\circ *F^*}$ be semisimple. Define
\[
\cE(G^F,s)\;:=\;\Bigl\{\;\chi\in \Irr(G^F)\ \Bigm|\
\Res^{G^F}_{G^{\circ F}}\chi\ \text{has some constituent }\chi^\circ\in \cE(G^{\circ F},s)\Bigr\}.
\]
Equivalently,
\[
\cE(G^F,s)\;=\;\bigcup_{\chi^\circ\in \cE(G^{\circ F},s)} \Irr(G^F\mid \chi^\circ),
\]
where $\Irr(G^F\mid \chi^\circ)$ denotes the set of irreducible characters of $G^F$ lying above $\chi^\circ$.
\end{definition}

\begin{proposition}[Partition of $\Irr(G^F)$ by identity-component series]
\label{prop:disc-series-partition}
As $s$ varies over the $G^{*F^*}$--conjugacy classes of semisimple elements of $(G^*)^{\circ F^*}$, the sets
$\cE(G^F,s)$ form a partition of $\Irr(G^F)$:
Let
\[
\mathcal S(G^*)
:=
\{\, s\in (G^*)^{\circ F^*}\mid s \text{ semisimple}\,\}/G^{*F^*}.
\]
Then
\[
\Irr(G^F)
=
\bigsqcup_{[s]\in \mathcal S(G^*)}
\cE(G^F,s).
\]

\end{proposition}

\begin{proof}
Let first $s,s'\in (G^*)^{\circ F^*}$ be semisimple elements whose images
$(s,1)$ and $(s',1)$ are $G^{*F^*}$--conjugate.  Choose
$(g,\omega)\in G^{*F^*}=G^{\circ *F^*}\rtimes \Omega$ such that
\[
(s',1)=(g,\omega)(s,1)(g,\omega)^{-1}=(g\,\sigma_\omega^*(s)\,g^{-1},1).
\]
Since connected Lusztig series depend only on the
$G^{\circ *F^*}$--conjugacy class of the semisimple parameter, this gives
\begin{equation}\label{eq:disc-series-conjugacy}
\cE(G^{\circ F},s')=
\cE(G^{\circ F},g\,\sigma_\omega^*(s)\,g^{-1})=
\cE(G^{\circ F},\sigma_\omega^*(s)).
\end{equation}
By functoriality of Lusztig series under pinned automorphisms, pullback along
$\sigma_\omega^{-1}$ induces a bijection
\[
\omega:\Irr(G^{\circ F})\to \Irr(G^{\circ F}),
\qquad
\chi^\circ\longmapsto \omega\chi^\circ:=\chi^\circ\circ\sigma_\omega^{-1},
\]
satisfying
\begin{equation}\label{eq:pinned-action-on-disc-series}
\omega\bigl(\cE(G^{\circ F},s)\bigr)
=
\cE(G^{\circ F},\sigma_\omega^*(s)).
\end{equation}
On the other hand, by Hypothesis~\ref{hyp:rational-pinned-component-condition}
we may choose a lift $\dot\omega\in G^F$ of $\omega$ whose conjugation action
on $G^{\circ F}$ is the pinned action corresponding to pullback by
$\sigma_\omega^{-1}$.  Conjugation by this lift permutes the irreducible
constituents of $\Res^{G^F}_{G^{\circ F}}\chi$.  An arbitrary lift differs
from it by an element of $G^{\circ F}$, hence gives the same action on
characters.  Therefore
\[
\Res^{G^F}_{G^{\circ F}}\chi \text{ has a constituent in }\cE(G^{\circ F},s)
\]
if and only if
\[
\Res^{G^F}_{G^{\circ F}}\chi \text{ has a constituent in }
\cE(G^{\circ F},\sigma_\omega^*(s)).
\]
Together with \eqref{eq:disc-series-conjugacy}, this shows that
\[
\chi\in \cE(G^F,s)
\quad\Longleftrightarrow\quad
\chi\in \cE(G^F,s'),
\]
and hence $\cE(G^F,s)$ depends only on the $G^{*F^*}$--conjugacy class of
$(s,1)$.

We next prove that the resulting subsets exhaust $\Irr(G^F)$.  Let
$\chi\in\Irr(G^F)$, and choose an irreducible constituent $\chi^\circ$ of
$\Res^{G^F}_{G^{\circ F}}\chi$.  For the connected group $G^\circ$, the
series $\cE(G^{\circ F},t)$, with $t$ running over semisimple elements of
$G^{\circ *F^*}$ up to $G^{\circ *F^*}$--conjugacy, form a partition of
$\Irr(G^{\circ F})$.  Thus $\chi^\circ\in \cE(G^{\circ F},s)$ for some
semisimple $s\in G^{\circ *F^*}=(G^*)^{\circ F^*}$, and
Definition~\ref{def:disc-Lusztig-series} gives $\chi\in \cE(G^F,s)$.

It remains to prove disjointness after indexing by $G^{*F^*}$--conjugacy
classes.  Suppose
\[
\chi\in \cE(G^F,s_1)\cap \cE(G^F,s_2)
\]
for semisimple $s_1,s_2\in (G^*)^{\circ F^*}$.  Choose constituents
$\chi_1^\circ,\chi_2^\circ$ of $\Res^{G^F}_{G^{\circ F}}\chi$ such that
$\chi_i^\circ\in \cE(G^{\circ F},s_i)$.  By Clifford theory, the constituents
of $\Res^{G^F}_{G^{\circ F}}\chi$ form a single $G^F$--orbit.  Hence there is
an element $g\in G^F$ such that
$\chi_2^\circ\simeq {}^g\chi_1^\circ$.  If $\omega$ denotes the image of $g$
in $\Omega$, then ${}^g\chi_1^\circ$ is the same constituent as
$\omega\chi_1^\circ$ in the notation above.  By
\eqref{eq:pinned-action-on-disc-series},
\[
\omega\chi_1^\circ\in
\cE(G^{\circ F},\sigma_\omega^*(s_1)).
\]
Thus
\[
\chi_2^\circ\in
\cE(G^{\circ F},s_2)
\cap
\cE(G^{\circ F},\sigma_\omega^*(s_1)).
\]
Since the connected Lusztig series form a partition of
$\Irr(G^{\circ F})$, the elements $s_2$ and $\sigma_\omega^*(s_1)$ are
$G^{\circ *F^*}$--conjugate.  Equivalently, $(s_2,1)$ and $(s_1,1)$ are
$G^{*F^*}$--conjugate.  Therefore two distinct $G^{*F^*}$--conjugacy classes
of semisimple elements in $(G^*)^{\circ F^*}$ give disjoint subsets.  Combining
exhaustion and disjointness proves the asserted partition.
\end{proof}

\begin{theorem}[Cuspidal enriched Jordan decomposition for disconnected $\bG$]
\label{thm:disc-cusp-JD}
Let \(\bG\) be a possibly disconnected reductive \(\F_q\)-group with abelian
component group, equipped with an \(F\)-stable pinning \(\cP\) of
\(\bG^\circ\).  Assume that \((\bG,\cP)\) satisfies
Hypothesis~\ref{hyp:rational-pinned-component-condition}.  For
\(s\in(\bG^*)^{\circ F^*}\) semisimple, put
\[
        \bH:=C_{\bG^*}(s),\qquad H:=\bH^{F^*}.
\]
There is a canonical bijection
\[
J^{\mathrm{enh},\cusp}_{G,s}:
\cE(G,s)_{\cusp}
\xrightarrow{\ \sim\ }
\Uch^{\mathrm{enh}}_{G,s}(H)_{\cusp},
\]
depending only on the fixed pinning.  If \(\rho\in\cE(G,s)_{\cusp}\) and
\(\rho^\circ\prec\Res^G_{G^\circ}\rho\), then
\[
J^{\mathrm{enh},\cusp}_{G,s}(\rho)
=
\left[
\J_s^{G^\circ}(\rho^\circ),
[\alpha_{\rho^\circ}],
E_\rho
\right],
\]
where \(E_\rho\in\Irr_{[\alpha_{\rho^\circ}]}(\Omega_{\rho^\circ})\) is the
Clifford label of \(\rho\) for the normal inclusion
\[
        G^\circ\lhd I_G(\rho^\circ).
\]
\end{theorem}

\begin{proof}
Let \(\rho\in\cE(G,s)_{\cusp}\), and choose an irreducible constituent
\[
        \rho^\circ\prec\Res^G_{G^\circ}\rho .
\]
Then \(\rho^\circ\in\cE(G^\circ,s)_{\cusp}\).  Put
\[
        u_0:=\J_s^{G^\circ}(\rho^\circ)
        \in \Uch\bigl(C_{\bG^{\circ *}}(s)^{F^*}\bigr)_{\cusp}.
\]
By Lemma~\ref{lem:disc-stabilizer-match}, the stabilizers of \(\rho^\circ\) and
\(u_0\) have the same image \(\Omega_{\rho^\circ}=\Omega_{u_0}\) in the pinned
component group.  Clifford theory for
\[
        G^\circ\lhd I_G(\rho^\circ)
\]
attaches to \(\rho^\circ\) a canonical class
\([\alpha_{\rho^\circ}]\in H^2(\Omega_{\rho^\circ},\C^\times)\), and, after
choosing any cocycle representative, parametrizes the irreducible characters of
\(G\) lying above \(\rho^\circ\) by
\[
        \Irr_{[\alpha_{\rho^\circ}]}(\Omega_{\rho^\circ}).
\]
Let \(E_\rho\) be the resulting Clifford label of \(\rho\).  We define
\[
        J^{\mathrm{enh},\cusp}_{G,s}(\rho)
        :=
        [u_0,[\alpha_{\rho^\circ}],E_\rho]
        \in \Uch^{\mathrm{enh}}_{G,s}(H)_{\cusp}.
\]

This is independent of the choice of \(\rho^\circ\).  Indeed, any other
irreducible constituent of \(\Res^G_{G^\circ}\rho\) is a \(G\)-conjugate of
\(\rho^\circ\).  Conjugation transports the inertia group, the quotient
\(\Omega_{\rho^\circ}\), the Clifford class, and the projective label, and the
quotient in Definition~\ref{def:enh-cusp-target} identifies the resulting data.

The inverse map is obtained by reversing this construction.  Given an enhanced
datum represented by \((u_0,[\alpha],E)\), take
\[
        \rho^\circ=(\J_s^{G^\circ})^{-1}(u_0).
\]
By the definition of the target, \([\alpha]=[\alpha_{\rho^\circ}]\), and Clifford
theory reconstructs a unique irreducible character
\[
        \rho_E\in\Irr(G\mid\rho^\circ)
\]
with Clifford label \(E\).  This construction is again invariant under replacing
the representative of the enhanced datum by a conjugate one.  Hence
\(J^{\mathrm{enh},\cusp}_{G,s}\) is a canonical bijection.
\end{proof}


\section{An enriched Jordan decomposition for disconnected $G$ with rationally pinned abelian component group}
\label{sec:disc-full-JD}

\noindent
Throughout this section we keep the standing notation of the preceding disconnected sections:
\(G\) is a possibly disconnected reductive \(\mathbb F_q\)-group with Frobenius \(F\),
\(G^\circ\lhd G\) is the identity component, we write \(G:=G^F\),
\(G^\circ:=(G^\circ)^F\), and we assume the component group
\[
        \Omega:=G/G^\circ
\]
is abelian.  Fix an \(F\)-stable pinning \(\mathbf P\) of \(G^\circ\), and assume
that the rational pinned-component condition of
Hypothesis~\ref{hyp:rational-pinned-component-condition} holds for \(G\), and for
each regular \(F\)-stable Levi subgroup to which the construction is applied.
We use the resulting semidirect-product model
\(G^*=G^{\circ *}\rtimes\Omega\).  For
\(s\in(G^*)^{\circ F^*}=G^{\circ *F^*}\) semisimple, put
\[
        H:=C_{G^*}(s)^{F^*}.
\]
We recall the disconnected Lusztig series in the identity-component case:
\[
\cE(G,s)
:=
\left\{
\chi\in\Irr(G):
\Res^G_{G^\circ}\chi\text{ has some constituent in }\cE(G^\circ,s)
\right\}.
\]

\medskip
We shall use the following standard consequence of the Mackey formula.

\begin{proposition}[Mackey criterion for disjoint Harish--Chandra series]
\label{prop:disjoint-HC-series-disconnected}
Let \(G\) be a possibly disconnected reductive \(\F_q\)-group.  Let
\((L,\tau)\) and \((M,\lambda)\) be cuspidal pairs of \(G\), with \(L,M\)
regular \(F\)-stable Levi subgroups, and let \(P,Q\) be \(F\)-stable
parabolic subgroups with Levi factors \(L,M\).  Assume that the Mackey
formula holds for the Lusztig functors attached to \((L,P)\) and
\((M,Q)\).  Then
\[
\left\langle
R^G_{L\subset P}(\tau),
R^G_{M\subset Q}(\lambda)
\right\rangle_{G^F}
=0
\]
unless \((L,\tau)\) and \((M,\lambda)\) are \(G^F\)-conjugate.  Consequently,
the sets
\[
\Irr(G,(L,\tau))
\]
are pairwise disjoint as \((L,\tau)\) ranges over \(G^F\)-conjugacy classes
of cuspidal pairs for which the relevant Mackey formula holds.
\end{proposition}

\subsection{Disconnected Lusztig series are unions of Harish--Chandra series}
\begin{proposition}
\label{prop:disc-Lusztig-union-HC}
Let \(s\in (G^*)^{\circ F^*}\) be semisimple.
For a regular \(F\)-stable Levi subgroup \(L\leq G\), choose an \(F\)-stable
parabolic \(P\) with Levi decomposition \(P=L\ltimes U\).  For a cuspidal pair
\((L,\tau)\) set
\[
\Irr(G,(L,\tau))
:=
\{\text{irreducible constituents of }R^G_{L\subset P}(\tau)\}.
\]
Let \(\Sigma_G(s)\) be the set of \(G\)-conjugacy classes of cuspidal pairs
\((L,\tau)\) such that \(\tau\in\cE(L,s)\).  Then
\[
\cE(G,s)
=
\bigcup_{[(L,\tau)]\in\Sigma_G(s)}\Irr(G,(L,\tau)).
\]
Moreover, this union is disjoint once the indexing is taken modulo
\(G\)-conjugacy of cuspidal pairs.
\end{proposition}

\begin{proof}
The proof is the same as in the ordinary disconnected setting, and uses only the
definition of the disconnected Lusztig series by restriction to \(G^\circ\).  If
\((L,\tau)\) is a cuspidal pair with \(\tau\in\cE(L,s)\), choose
\(\tau^\circ\prec\Res^L_{L^\circ}\tau\) with
\(\tau^\circ\in\cE(L^\circ,s)\).  For any irreducible constituent
\(\chi\) of \(R^G_{L\subset P}(\tau)\), compatibility of Harish--Chandra
restriction with restriction to identity components gives a constituent
\(\chi^\circ\prec\Res^G_{G^\circ}\chi\) which occurs in
\(R^{G^\circ}_{L^\circ\subset P^\circ}(\tau^\circ)\).  The connected
Harish--Chandra decomposition of Lusztig series gives
\(\chi^\circ\in\cE(G^\circ,s)\), hence \(\chi\in\cE(G,s)\).

Conversely, if \(\chi\in\cE(G,s)\), choose
\(\chi^\circ\prec\Res^G_{G^\circ}\chi\) with
\(\chi^\circ\in\cE(G^\circ,s)\).  In the connected group \(G^\circ\), the
character \(\chi^\circ\) belongs to a Harish--Chandra series attached to some
cuspidal pair \((L^\circ,\tau^\circ)\) with \(\tau^\circ\in\cE(L^\circ,s)\).
Extending \(L^\circ\) to the corresponding regular Levi subgroup \(L\) of
\(G\), adjointness and the compatibility of Harish--Chandra restriction with
restriction to identity components produce an irreducible constituent
\(\tau\) of \({}^*R^G_{L\subset P}(\chi)\) lying above \(\tau^\circ\).  The same
restriction argument shows that \(\tau\) is cuspidal; hence
\(\tau\in\cE(L,s)\), and \(\chi\in\Irr(G,(L,\tau))\).

For Lusztig functors attached to \(F\)-stable parabolics in the present
disconnected setting, the Mackey formula holds by
\cite[Theorem~3.2]{DigneMichel94}.  Proposition~\ref{prop:disjoint-HC-series-disconnected}
therefore gives disjointness modulo \(G\)-conjugacy of cuspidal pairs.
\end{proof}

\medskip
\subsection{Enriched Harish--Chandra labels}
\label{subsec:enriched-HC-labels-disconnected}

Let \(L\leq G\) be a regular \(F\)-stable Levi subgroup, and let
\[
        \tau\in\cE(L,s)_{\cusp}.
\]
Choose an irreducible constituent
\[
        \tau^\circ\prec\Res^L_{L^\circ}\tau .
\]
Put
\[
        H_L:=C_{L^*}(s)^{F^*},
        \qquad
        K_0:=C_{(L^*)^\circ}(s)^{F^*},
\]
and
\[
        u_{L,0}:=\J_{s}^{L^\circ}(\tau^\circ)
        \in \Uch(K_0)_{\cusp}.
\]
By Lemma~\ref{lem:disc-stabilizer-match}, applied to the disconnected group
\(L\), the source inertia quotient
\[
        \Omega_{\tau^\circ}:=I_L(\tau^\circ)/L^\circ
\]
is identified with the corresponding stabilizer quotient of \(u_{L,0}\) on the
centralizer side.  Clifford theory on the source extension
\[
        L^\circ\lhd I_L(\tau^\circ)
\]
gives a canonical class
\[
        [\alpha_{\tau^\circ}]\in H^2(\Omega_{\tau^\circ},\C^\times).
\]
Let
\[
        E_\tau\in\Irr_{[\alpha_{\tau^\circ}]}(\Omega_{\tau^\circ})
\]
be the Clifford label of \(\tau\) above \(\tau^\circ\).  We write
\begin{equation}\label{eq:enh-cuspidal-datum-tau}
        \mathfrak u^{\mathrm{enh}}_\tau
        :=
        [u_{L,0},[\alpha_{\tau^\circ}],E_\tau]
        \in \Uch^{\mathrm{enh}}_{L,s}(H_L)_{\cusp}.
\end{equation}

Let \(W_{\tau^\circ}\) be the connected relative Weyl group attached to the
cuspidal pair \((L^\circ,\tau^\circ)\) in \(G^\circ\).  The corner lemma
\ref{lem:corner-lemma}, applied with \(X=G\), \(M=L\), and
\(\theta^\circ=\tau^\circ\), shows that \(W_{\tau^\circ}\) acts on the set of
Clifford labels \(\Irr_{[\alpha_{\tau^\circ}]}(\Omega_{\tau^\circ})\).  Let
\[
        W_\tau:=\Stab_{W_{\tau^\circ}}(E_\tau).
\]
The same corner lemma supplies a canonically defined cohomology class
\[
        [\beta_\tau]\in H^2(W_\tau,\C^\times)
\]
which measures the possible projectivity of the action of \(W_\tau\) on the
Clifford multiplicity space.  After choosing a representative
\(\beta_\tau\in Z^2(W_\tau,\C^\times)\), the irreducible constituents of
Harish--Chandra induction from \((L,\tau)\) are parametrized by
\[
        \mathfrak R_\tau
        :=
        \Irr\bigl(\C_{\beta_\tau}[W_\tau]\bigr),
\]
with the usual coboundary indeterminacy in \(\beta_\tau\).  Thus each
\[
        \rho\in\Irr(G,(L,\tau))
\]
has a canonical relative Harish--Chandra label
\[
        \varphi_\rho\in\mathfrak R_\tau .
\]
When the corner cocycle \(\beta_\tau\) is trivial, this label set is simply
\(\Irr(W_\tau)\), which is the untwisted situation used in the ordinary
Howlett--Lehrer formulation.

\begin{definition}[Full enriched unipotent target]
\label{def:full-enriched-target}
Define
\[
\Uch^{\mathrm{enh}}_{G,s}(H)
:=
\bigsqcup_{[(L,\tau)]\in\Sigma_G(s)}
\left\{
        [H_L,\mathfrak u^{\mathrm{enh}}_\tau,\varphi]
        :
        \varphi\in\mathfrak R_\tau
\right\}.
\]
Here the brackets mean that replacing \((L,\tau,\tau^\circ)\) by a
\(G\)-conjugate replaces \(H_L\), the enriched cuspidal datum, and the relative
label by the conjugate data.
\end{definition}

\subsection{Jordan decomposition for disconnected $G$ with rationally pinned abelian component group: enriched bijection part}
\begin{theorem}[Enriched pinned disconnected Jordan decomposition]
\label{thm:disc-JD-bijection}
Let \(G\) be as above; in particular, assume that the finite component group is
abelian and that the rational pinned-component condition holds.  Let
\(s\in (G^*)^{\circ F^*}\) be semisimple and put \(H=C_{G^*}(s)^{F^*}\).  There
is a canonical bijection
\[
J^{\mathrm{enh}}_{G,s}:
\cE(G,s)
\xrightarrow{\ \sim\ }
\Uch^{\mathrm{enh}}_{G,s}(H),
\]
depending only on the fixed pinning.  On the cuspidal part it is the bijection
of Theorem~\ref{thm:disc-cusp-JD}.  More generally, if
\[
        \rho\in\Irr(G,(L,\tau))
\]
for a cuspidal pair \((L,\tau)\in\Sigma_G(s)\), then
\[
        J^{\mathrm{enh}}_{G,s}(\rho)
        =
        [H_L,\mathfrak u^{\mathrm{enh}}_\tau,\varphi_\rho],
\]
where \(\mathfrak u^{\mathrm{enh}}_\tau\) is defined by
\eqref{eq:enh-cuspidal-datum-tau} and \(\varphi_\rho\) is the relative
Harish--Chandra label of \(\rho\).
\end{theorem}

\begin{proof}
By Proposition~\ref{prop:disc-Lusztig-union-HC},
\[
        \cE(G,s)
        =
        \bigsqcup_{[(L,\tau)]\in\Sigma_G(s)}
        \Irr(G,(L,\tau)).
\]
For each summand, the source-side disconnected Howlett--Lehrer analysis recalled
above gives a canonical parametrization
\[
        \Irr(G,(L,\tau))
        \xrightarrow{\ \sim\ }
        \mathfrak R_\tau,
        \qquad
        \rho\longmapsto\varphi_\rho .
\]
Combining this label with the enriched cuspidal datum
\(\mathfrak u^{\mathrm{enh}}_\tau\) gives a map
\[
        \rho
        \longmapsto
        [H_L,\mathfrak u^{\mathrm{enh}}_\tau,\varphi_\rho]
\]
from the Harish--Chandra series \(\Irr(G,(L,\tau))\) onto the corresponding
summand of \(\Uch^{\mathrm{enh}}_{G,s}(H)\).  These maps are bijective on each
summand, and the summands are disjoint on both sides.  Taking their disjoint
union gives the asserted bijection.

If \(\rho\) is cuspidal, then the relevant cuspidal pair is \((G,\rho)\), the
relative Weyl group is trivial, and the construction reduces exactly to the
cuspidal enriched bijection of Theorem~\ref{thm:disc-cusp-JD}.  Conjugating the
choice of \(\tau^\circ\) conjugates the inertia quotient, Clifford class,
Clifford label, and relative Harish--Chandra label compatibly, so the resulting
element of the quotient target is independent of all auxiliary choices.
\end{proof}

\begin{remark}[Recovery of an ordinary disconnected target]
\label{rmk:ordinary-disc-JD-recovery}
The theorem above does not assert a bijection
\[
        \cE(G,s)\longrightarrow \Uch(C_{G^*}(s)^{F^*})
\]
under Hypothesis~\ref{hyp:rational-pinned-component-condition} alone.  Such an
ordinary target is recovered only after adding the extra assertion that, for
every relevant cuspidal constituent \(\rho^\circ\), the transported source
Clifford class \([\alpha_{\rho^\circ}]\) agrees with the ordinary Clifford class
of the unipotent centralizer side, and that the corresponding relative
Howlett--Lehrer corner labels are identified.  Under that additional condition,
the enriched datum \([u_0,[\alpha_{\rho^\circ}],E]\) is realized by the usual
irreducible unipotent character of \(C_{G^*}(s)^{F^*}\) with the same Clifford
label, and \(J^{\mathrm{enh}}_{G,s}\) descends to the ordinary disconnected
pinned Jordan decomposition.
\end{remark}

\subsection{Conditional ordinary packet sums}
\label{subsec:conditional-ordinary-packet-sums}

The preceding construction deliberately keeps the source Clifford class as part
of the target.  Therefore the regular packet-sum statements that would be phrased
in terms of ordinary irreducible unipotent characters of
\(C_{G^*}(s)^{F^*}\) require the additional ordinary-recovery hypothesis of
Remark~\ref{rmk:ordinary-disc-JD-recovery}.  Under that extra hypothesis, the
standard Clifford regular-sum argument applies: summing over all projective
labels in a fixed Clifford packet gives the induced character from the identity
component, and hence the sum depends only on the underlying orbit-valued Jordan
decomposition.  Without that extra hypothesis, the correct invariant object is
the enriched packet
\[
        \{[u_0,[\alpha_{\rho^\circ}],E]:
        E\in\Irr_{[\alpha_{\rho^\circ}]}(\Omega_{\rho^\circ})\},
\]
not an ordinary packet inside \(\Uch(C_{G^*}(s)^{F^*})\).



\subsection{Stable finite regular sums}
\label{subsec:stable-finite-regular-sums}

The enriched disconnected Jordan decomposition records more information than the
orbit-valued Jordan decomposition: it remembers the source Clifford class and
the projective Clifford label.  Stability, however, is insensitive to the
individual projective label.  The stable finite objects used in the local
application are obtained by summing regularly over the Clifford packet.  We
record this compatibility in a form which does not require an ordinary
unipotent-character target for the disconnected centralizer.

\begin{lemma}[Independence of the Clifford-regular sum]
\label{lem:clifford-regular-sum-independent}
Let \(N\lhd M\) be finite groups, and let \(\theta\in\Irr(N)\).  Put
\[
        I_M(\theta):=\{m\in M\mid {}^m\theta\simeq \theta\},
        \qquad
        \Omega_\theta:=I_M(\theta)/N .
\]
Choose a projective extension \(\widetilde\theta\) of \(\theta\) to
\(I_M(\theta)\), and let
\[
        \alpha_\theta\in Z^2(\Omega_\theta,\C^\times)
\]
be the Clifford-label cocycle in the convention of
\S\ref{subsec:clifford-labels}.  For
\[
        E\in\Irr(\C_{\alpha_\theta}[\Omega_\theta]),
\]
let
\[
        \chi_{\theta,E}
        :=
        \Ind_{I_M(\theta)}^M
        \bigl(\widetilde\theta\otimes
        \infl_{\Omega_\theta}^{I_M(\theta)}E\bigr)
        \in \Irr(M\mid \theta).
\]
Then
\[
        \sum_{E\in\Irr(\C_{\alpha_\theta}[\Omega_\theta])}
        \dim(E)\,\chi_{\theta,E}
        =
        \Ind_N^M\theta
\]
as characters of \(M\).  In particular, this regular sum is independent of the
chosen projective extension, of the cocycle representative, and of the Clifford
parametrization.  It depends only on the \(M\)-orbit of \(\theta\).
\end{lemma}

\begin{proof}
Let \(\Reg_{\alpha_\theta}\) be the left regular module of the twisted group
algebra \(\C_{\alpha_\theta}[\Omega_\theta]\).  Since this algebra is semisimple,
its regular character decomposes as
\[
        \Reg_{\alpha_\theta}
        =
        \sum_{E\in\Irr(\C_{\alpha_\theta}[\Omega_\theta])}
        \dim(E)\,E .
\]
Therefore
\[
        \sum_E \dim(E)
        \bigl(\widetilde\theta\otimes
        \infl_{\Omega_\theta}^{I_M(\theta)}E\bigr)
\]
is the character of
\[
        \widetilde\theta\otimes
        \infl_{\Omega_\theta}^{I_M(\theta)}\Reg_{\alpha_\theta}.
\]
If \(i\in I_M(\theta)\) maps to a nontrivial element of \(\Omega_\theta\), then
left multiplication by its image on the regular module has trace zero.  If
\(i=n\in N\), then the image of \(i\) in \(\Omega_\theta\) is trivial and the
trace on the regular module is \(|\Omega_\theta|\).  Since
\(\widetilde\theta|_N=\theta\), the above character has value zero on
\(I_M(\theta)\setminus N\) and value
\[
        |\Omega_\theta|\,\theta(n)
\]
on \(n\in N\).  These are precisely the values of
\(\Ind_N^{I_M(\theta)}\theta\), because \(\theta\) is
\(I_M(\theta)\)-stable.  Hence
\[
        \sum_E \dim(E)
        \bigl(\widetilde\theta\otimes
        \infl_{\Omega_\theta}^{I_M(\theta)}E\bigr)
        =
        \Ind_N^{I_M(\theta)}\theta .
\]
Inducing from \(I_M(\theta)\) to \(M\) and using transitivity of induction gives
\[
        \sum_E \dim(E)\,\chi_{\theta,E}
        =
        \Ind_N^M\theta .
\]
The right-hand side is intrinsic.  Replacing \(\theta\) by an \(M\)-conjugate
also leaves \(\Ind_N^M\theta\) unchanged.  This proves the asserted
independence.
\end{proof}

\begin{theorem}[Stable finite regular sums and enriched pinned Jordan decomposition]
\label{thm:pinned-JD-stable-sum-compatible}
Keep the notation of Theorem~\ref{thm:disc-JD-bijection}.  Thus \(G\) is a
possibly disconnected finite reductive group satisfying the rational
pinned-component hypothesis, \(G^\circ\lhd G\), and
\[
        s\in (G^*)^{\circ F^*}.
\]
Put
\[
        H=C_{G^*}(s)^{F^*},
        \qquad
        H_0=C_{(G^*)^\circ}(s)^{F^*}.
\]
Let
\[
        J^{\mathrm{enh}}_{G,s}:
        \cE(G,s)
        \xrightarrow{\ \sim\ }
        \Uch^{\mathrm{enh}}_{G,s}(H)
\]
be the enriched pinned Jordan decomposition, and let
\(J^{\mathrm{enh},\cusp}_{G,s}\) denote its cuspidal restriction from
Theorem~\ref{thm:disc-cusp-JD}.  Let
\[
        O\subset \Uch(H_0)_{\cusp}
\]
be an \(H/H_0\)-orbit representing a stable connected finite unipotent datum,
and choose \(u_0\in O\).  Let
\[
        \rho^\circ=(\J_s^{G^\circ})^{-1}(u_0)
        \in \cE(G^\circ,s)_{\cusp}.
\]
For
\[
        E\in\Irr_{[\alpha_{\rho^\circ}]}(\Omega_{\rho^\circ}),
\]
let \(\rho_E\in\cE(G,s)_{\cusp}\) be the unique character satisfying
\[
        J^{\mathrm{enh},\cusp}_{G,s}(\rho_E)
        =
        [u_0,[\alpha_{\rho^\circ}],E].
\]
Then the Clifford-regularized finite sum
\[
        \Theta^{\mathrm{enh}}_{s,O}
        :=
        \sum_{E\in\Irr_{[\alpha_{\rho^\circ}]}(\Omega_{\rho^\circ})}
        \dim(E)\,\rho_E
\]
satisfies
\[
        \Theta^{\mathrm{enh}}_{s,O}
        =
        \Ind_{G^\circ}^{G}\rho^\circ .
\]
Consequently \(\Theta^{\mathrm{enh}}_{s,O}\) is independent of the representative
\(u_0\in O\), of the projective extension used to define the Clifford class, of
the cocycle representative, and of the pinned refinement.  It depends only on
the orbit-valued connected Jordan datum \(O\).

More generally, if a connected finite stable packet coefficient is a finite
linear combination of such orbit data, then its enriched Clifford-regularization
is obtained by applying the preceding construction orbit by orbit.  Hence the
enriched regularized coefficient has the same finite stable-transport invariance
as the connected finite stable coefficient.
\end{theorem}

\begin{proof}
By the construction of the cuspidal enriched bijection in
Theorem~\ref{thm:disc-cusp-JD}, the characters of \(G\) lying above
\(\rho^\circ\) are precisely the characters \(\rho_E\) obtained by Clifford
theory for the normal inclusion
\[
        G^\circ\lhd G
\]
with projective label
\[
        E\in\Irr_{[\alpha_{\rho^\circ}]}(\Omega_{\rho^\circ}).
\]
Applying Lemma~\ref{lem:clifford-regular-sum-independent} with
\[
        N=G^\circ,
        \qquad
        M=G,
        \qquad
        \theta=\rho^\circ
\]
gives
\[
        \sum_E \dim(E)\,\rho_E
        =
        \Ind_{G^\circ}^{G}\rho^\circ .
\]
This proves the displayed identity and shows at once that the sum is independent
of all choices entering the Clifford parametrization.

If \(u_1={}^h u_0\) is another representative of the \(H/H_0\)-orbit \(O\), then
Lemma~\ref{lem:disc-stabilizer-match} and the equivariance of the connected
Jordan decomposition identify
\[
        (\J_s^{G^\circ})^{-1}(u_1)
\]
with a \(G\)-conjugate of \(\rho^\circ\).  The induced character
\(\Ind_{G^\circ}^{G}\rho^\circ\) is unchanged by replacing \(\rho^\circ\) by a
\(G\)-conjugate.  Therefore \(\Theta^{\mathrm{enh}}_{s,O}\) depends only on the
orbit \(O\), not on its representative.  Since the right-hand side only sees the
orbit-valued connected Jordan datum, the regularized sum is also independent of
the additional pinned refinement which selects the individual enriched labels.

Finally, Lusztig's finite character-sheaf theory identifies the connected
orbit-valued packet coefficients attached to stable finite unipotent data with
stable almost characters; see \cite[Ch.~13]{Lus84}, \cite{LusztigCS5}, and, for
the disconnected-centre form of the Jordan decomposition, \cite{Lus88,DM90}.
The preceding paragraphs show that the enriched disconnected coefficient is the
Clifford regularization of this connected stable coefficient.  Clifford
regularization is functorial for conjugacy transport, because it is expressed as
\(\Ind_{G^\circ}^{G}\rho^\circ\) and induction across the normal inclusion
\(G^\circ\lhd G\) commutes with such transport.  Therefore any finite stable
transport preserving the connected coefficient also preserves the enriched
regularized coefficient.  This is the finite stability property used in the
local depth-zero application.
\end{proof}


\section{Finite-field preservation of rank-one Hecke parameters}
\label{sec:finite-rank-one-parameter-comparison}

In this section, we reprove \cite[Proposition 3.2.3]{ohara2025} by reducing the comparison of $q$-parameters to a comparison of Schur elements arising from the endomorphism algebras of Harish--Chandra induced representations.

Let $\bG$ be a connected reductive group over a finite field $\fq$, and let $\bL \subseteq \bG$ be an $F$-stable Levi subgroup of an $F$-stable parabolic subgroup of $\bG$.
Let $\tau$ be an irreducible cuspidal representation of $L$, and assume $\tau \in \cE(L, s)$ for a semisimple element $s \in L^*$, where $\bL^*$ is the dual group of $\bL$ and $\bL^* \subset \bG^*$ is an $F^*$-stable Levi subgroup of an $F^*$-stable parabolic subgroup of $\bG^*$.

Let $u_{\tau} \in \Uch(H_L)$ denote the unipotent representation of the centralizer $H_L := C_{L^*}(s)$ associated with $\tau$ through the Jordan decomposition of  characters $\J_s^L$ as in Theorem \ref{thm:JD-connected-groups}.
Consider the Harish-Chandra inductions:
\[
\pi := R_{\bL}^{\bG}(\tau), \qquad \pi^* := R_{\bH_L}^{\bH}(u_{\tau}).
\]
Assume that \(R_{\bL}^{\bG}(\tau)\) has length at most two.  If
\[
R_{\bL}^{\bG}(\tau)=\pi_1\oplus\pi_2
\]
has two irreducible constituents, ordered so that
\(\dim\pi_1\geq \dim\pi_2\), set
\[
q\bigl(R_{\bL}^{\bG}(\tau)\bigr)
=
\frac{\dim\pi_1}{\dim\pi_2}.
\]
If \(R_{\bL}^{\bG}(\tau)\) is irreducible, set
\[
q\bigl(R_{\bL}^{\bG}(\tau)\bigr)=1.
\]
We aim to prove the following.
\[
q(\pi) = q(\pi^*),
\]
\subsection{Schur element}
Associated with the cuspidal pair $(L, \tau)$ is the relative Hecke algebra
\[
\mathcal{H}_G(M, \tau) := \End_{G}(R_{\bL}^{\bG}(\tau)),
\]
which is a finite type symmetric algebra over a ring $A = \mathbb{Z}[\mathbf{x}^{\pm1}]$ of Hecke parameters.
Let $W_{\tau}$ denote the relative Weyl group associated with the pair $(L,\tau)$,
and let $K$ be a splitting field for $\mathcal{H}_G(L, \tau)$.

By \cite[Theorem~3.2.5]{Book:GeckandMalle}, there is a well-known bijection:
\[
\mathrm{Irr}(\mathcal{H}_G(L, \tau)) \longleftrightarrow \Irr(G, (L, \tau)),
\]
where each irreducible character $\phi \in \Irr(W_{\tau})$ corresponds to a unique irreducible constituent $\pi_\phi$ of $\pi = R_{\bL}^{\bG}(\tau)$.
Let $c_\phi \in A'$ denote the \textit{Schur element} associated with $\phi$ with respect to the canonical symmetrizing form $\lambda$ on $\mathcal{H}_G(L, \tau)$, where $A'$ denotes the integral closure of $A$ in $K$. Then the character of $\lambda$ decomposes as:
\[
\lambda = \sum_{\phi \in \Irr(W_{\tau})} \frac{1}{c_\phi} \chi_\phi,
\]
where $\chi_\phi$ is the character of the simple module corresponding to $\phi$.
Upon specialization $\mathbf{x} \mapsto 1$, the Schur elements $c_\phi(q)$ become rational numbers.
By \cite[Theorem~3.2.18]{Book:GeckandMalle}, the dimension of the irreducible representation $\pi_{\phi}\in \Irr(G,(L,\tau))$ corresponding to the irreducible representation $\phi\in\Irr(W_{\tau})$ is given by:
\[
\dim(\pi_\phi) = \dim(\pi) \cdot c_\phi(q)^{-1}.
\]
In particular, if $\pi = \pi_1 \oplus \pi_2$ corresponds to the characters $\phi_1, \phi_2 \in \Irr(W_{\tau})$, then:
\[
q(\pi)= \frac{\dim(\pi_1)}{\dim(\pi_2)} = \frac{c_{\phi_2}(q)}{c_{\phi_1}(q)}.
\]

\subsection{Compatibility Under Jordan Decomposition of characters}

By Theorem \ref{thm:JD-connected-groups}, the Jordan decomposition of characters can be chosen to commute with Harish-Chandra induction.
Moreover, for any Levi subgroup $\bL \subseteq \bG$ as before, and any $\tau \in \cE(L, s)$, we have a canonical isomorphism of relative Hecke algebras:
\[
\mathcal{H}_G(L, \tau) \cong \mathcal{H}_{H}(H_L, u_{\tau}),
\]
which preserves the indexing of irreducible characters and the associated Schur elements.
Therefore, the constituents $u_1, u_2$ of $\pi^* := R_{\bH_L}^{\bH}(u_{\tau})$ correspond to the same $\phi_1, \phi_2 \in \Irr(W_{\tau})=\Irr(W_{u_\tau})$, and satisfy:
\[
\dim(u_i) = \dim(\pi^*) \cdot c_{\phi_i}(q)^{-1}, \quad \text{for } i = 1, 2.
\]
Consequently,
\[
q(\pi^*) = \frac{\dim(u_1)}{\dim(u_2)} = \frac{c_{\phi_2}(q)}{c_{\phi_1}(q)} = q(\pi).
\]
As a consequence, we get that
\begin{theorem}[Preservation of rank-one finite-field parameters under Jordan decomposition]
\label{thm:finite-JD-preserves-q-parameter}
With notation as above, assume that \(R_{\bL}^{\bG}(\tau)\) has length at most two.
Then \(R_{\bH_L}^{\bH}(u_\tau)\) also has length at most two, and
\[
q\left( R_{\bL}^{\bG}(\tau) \right)
=
q\left( R_{\bH_L}^{\bH}(u_{\tau}) \right).
\]
\end{theorem}


\section{Application: depth-zero to unipotent affine Hecke algebra factors}
\label{sec:padic-depth0-to-unipotent-hecke}
\label{sec:finite-depthzero-to-unipotent-hecke}

In this section we explain how the finite-field comparison proved above gives a
pinned canonical form of Ohara's comparison between depth-zero affine Hecke
algebra parameters and unipotent affine Hecke algebra parameters.  The conclusion
is deliberately restricted to the affine Hecke algebra factor.  The full
AFMO--Morris Hecke algebra usually has the form of an affine Hecke algebra
crossed with a finite twisted group algebra, and Ohara's theorem does not identify
that finite crossed-product factor with the corresponding finite crossed-product
factor for the unipotent type.

Thus the result below should be read as a canonical version of the affine-factor
comparison in \cite[Theorem~4.4.1]{ohara2025}.  It does not assert a general
isomorphism from a depth-zero Hecke algebra to the Hecke algebra of a unipotent
type.  Instead, it says that, after fixing the pinning, the affine Hecke algebra
factor of the depth-zero algebra may be replaced by the corresponding unipotent
affine Hecke algebra factor, while the finite group \(\Omega(\rho_M)\) and the
cocycle \(\mu\) in the crossed product remain the original source-side ones.

The only additional point which must be made explicit for our use of the finite
Jordan decomposition is that the finite groups occurring in the depth-zero
construction are not always the connected reductive quotients
\(G(F)_{x,0}/G(F)_{x,0+}\).  They are often the full quotients
\(G(F)_x/G(F)_{x,0+}\), and these may be disconnected.  We therefore first record
why these full quotients satisfy the hypotheses imposed in the enriched
disconnected Jordan decomposition theorem.

\subsection{Full parahoric quotients and the pinned-component condition}
\label{subsec:parahoric-quotients-pinned-condition}

Let \(F\) be a non-archimedean local field with residue field \(\ff\), and let
\(G\) be a connected reductive group over \(F\).  For a vertex
\(x\in\mathcal B(G,F)\), put
\[
        K_x:=G(F)_x,\qquad
        K_x^0:=G(F)_{x,0},\qquad
        K_x^+:=G(F)_{x,0+}.
\]
We write
\[
        \overline G_x:=K_x/K_x^+,
        \qquad
        \overline G_x^\circ:=K_x^0/K_x^+ .
\]
Thus \(\overline G_x^\circ\) is the group of \(\ff\)-points of the connected
reductive quotient of the parahoric group scheme attached to \(K_x^0\).  To avoid
confusing this finite group with an integral model, we denote that connected
reductive \(\ff\)-group by
\[
        \underline{\mathbf G}_x^\circ,
        \qquad
        \underline{\mathbf G}_x^\circ(\ff)=\overline G_x^\circ .
\]
The full quotient \(\overline G_x\) will be regarded as the group of \(\ff\)-points
of the corresponding possibly disconnected reductive quotient
\(\underline{\mathbf G}_x\), with identity component
\(\underline{\mathbf G}_x^\circ\).

Fix the global pinned datum \(\cP\) used throughout the paper.  For each
vertex \(x\) occurring below, we denote by
\[
        \cP_x^{\cP}
\]
the \(\ff\)-pinning of \(\underline{\mathbf G}_x^\circ\) induced by the global
pinning \(\cP\).  Concretely, the construction may be described by taking the
\(\cP\)-compatible alcove adjacent to \(x\), using it to determine the Borel
subgroup of \(\underline{\mathbf G}_x^\circ\), and reducing the root-group
parametrizations of \(\cP\) through the Bruhat--Tits root-group quotients at
\(x\).  We shall sometimes denote this local alcove by \(\mathfrak a\) inside
proofs, but \(\mathfrak a\) is not an additional datum: the finite-quotient
pinning used below is \(\cP_x^{\cP}\), determined by the fixed global pinning
\(\cP\).

\begin{lemma}[Component group of the full parahoric quotient]
\label{lem:parahoric-component-abelian}
With notation as above, the component group of the full quotient is abelian.  More
precisely, the Kottwitz homomorphism induces an injective homomorphism
\[
        K_x/K_x^0
        \hookrightarrow
        X^*(Z(\widehat G)^I)^\Phi .
\]
Consequently
\[
        \overline G_x/\overline G_x^\circ
        \cong
        K_x/K_x^0
\]
is abelian.
\end{lemma}

\begin{proof}
The parahoric subgroup \(K_x^0\) is the kernel, inside the stabilizer \(K_x\), of
the Kottwitz invariant.  Equivalently, the Kottwitz map is trivial on
\(G(F)_{x,0}\) and separates the components of the stabilizer of the vertex
\(x\).  Hence it induces the displayed injection
\[
        K_x/K_x^0\hookrightarrow X^*(Z(\widehat G)^I)^\Phi .
\]
The last assertion follows because
\[
        \overline G_x/\overline G_x^\circ
        =(K_x/K_x^+)/(K_x^0/K_x^+)
        \cong K_x/K_x^0 .
\]
\end{proof}

\begin{lemma}[Parahoric component actions are pinned modulo the connected quotient]
\label{lem:parahoric-component-action-pinned}
Let \(\bar k\in\overline G_x\), and let
\[
        \alpha_{\bar k}\in\Aut_\ff(\underline{\mathbf G}_x^\circ)
\]
be the automorphism induced by conjugation on
\(\overline G_x^\circ=K_x^0/K_x^+\).  Then there exist
\[
        \bar g_{\bar k}\in \overline G_x^\circ
        \qquad\text{and}\qquad
        \sigma_{\bar k}\in
        \Aut_\ff(\underline{\mathbf G}_x^\circ,\cP_x^{\cP})
\]
such that
\[
        \alpha_{\bar k}
        =
        \operatorname{Int}(\bar g_{\bar k})\circ\sigma_{\bar k}.
\]
Equivalently, if
\[
        \overline G_{x,\cP}
        :=
        \{\bar k\in\overline G_x:
        \Ad(\bar k)(\cP_x^{\cP})=
        \cP_x^{\cP}\},
\]
then
\[
        \overline G_x
        =
        \overline G_x^\circ\cdot \overline G_{x,\cP}.
\]
In particular, the possibly disconnected finite reductive quotient
\(\underline{\mathbf G}_x\), together with the pinning
\(\cP_x^{\cP}\) of its identity component, satisfies
Hypothesis~\ref{hyp:rational-pinned-component-condition}.
\end{lemma}

\begin{proof}
We recall the standard normalizer description of parahoric stabilizers.  Let
\(\mathfrak a=\mathfrak a_{\cP,x}\) denote the \(\cP\)-compatible local alcove
used to describe the induced pinning \(\cP_x^{\cP}\).  Choose a maximal
\(F\)-split torus \(S\) whose apartment contains \(x\) and this alcove.  Let
\[
        N=N_G(S)(F),
        \qquad
        \widetilde W=N/S(F)_0
\]
be the Iwahori--Weyl group.  The stabilizer \(K_x\) is generated by \(K_x^0\) and
the stabilizer of \(x\) in \(N\).  Hence every class in
\(K_x/K_x^0\) has a representative \(n\in N\cap K_x\).

Let \(k\in K_x\) be a lift of \(\bar k\), and choose such an \(n\) in the same
class modulo \(K_x^0\).  Then \(kn^{-1}\in K_x^0\).  If \(\bar h\) denotes the
image of \(kn^{-1}\) in \(\overline G_x^\circ\), then on
\(\underline{\mathbf G}_x^\circ\)
\[
        \alpha_{\bar k}
        =
        \operatorname{Int}(\bar h)\circ\alpha_{\bar n},
\]
where \(\bar n\) is the image of \(n\) in \(\overline G_x\).  It is therefore
enough to treat representatives coming from \(N\cap K_x\).

For such an \(n\), conjugation by \(n\) stabilizes \(x\) and sends the alcove
\(\mathfrak a\) to another alcove whose closure contains \(x\).  On the reductive
quotient, it sends the Borel, torus, and simple root-group quotients defining
\(\cP_x^{\cP}\) to the corresponding pinned data attached to
\(n\mathfrak a\), with the root-group parametrizations transported by the fixed
Tits normalizations.  The two alcoves \(\mathfrak a\) and \(n\mathfrak a\)
determine two \(\ff\)-rational pinnings of the connected reductive quotient
\(\underline{\mathbf G}_x^\circ\) arising from the same Bruhat--Tits root-group
parametrizations.  The two chamber descriptions \(\mathfrak a\) and \(n\mathfrak a\) determine
two normalized pinnings of \(\underline{\mathbf G}_x^\circ\), both obtained
from the fixed global pinning \(\cP\) by the Bruhat--Tits root-group
filtrations and the Tits normalizations.  In the finite BN-pair of
\(\underline{\mathbf G}_x^\circ(\ff)\), the chambers corresponding to
\(\mathfrak a\) and \(n\mathfrak a\) differ by an element of the finite Weyl
group.  Choosing the corresponding Tits representative gives an element
\(\bar b_n\in \overline G_x^\circ\) which carries this induced normalized
pinning attached to \(n\mathfrak a\) back to the induced normalized pinning
attached to \(\mathfrak a\).
Equivalently,
\[
        \operatorname{Int}(\bar b_n)\circ\alpha_{\bar n}
        \in
        \Aut_\ff(\underline{\mathbf G}_x^\circ,\cP_x^{\cP}).
\]
Set
\[
        \sigma_{\bar n}:=
        \operatorname{Int}(\bar b_n)\circ\alpha_{\bar n}.
\]
Then
\[
        \alpha_{\bar n}
        =
        \operatorname{Int}(\bar b_n^{-1})\circ\sigma_{\bar n}.
\]
Combining this with the previous decomposition of \(\alpha_{\bar k}\), we obtain
\[
        \alpha_{\bar k}
        =
        \operatorname{Int}(\bar h\bar b_n^{-1})\circ\sigma_{\bar n}.
\]
This is the required formula, with
\(\bar g_{\bar k}=\bar h\bar b_n^{-1}\) and
\(\sigma_{\bar k}=\sigma_{\bar n}\).

The equivalent factorization of \(\overline G_x\) follows immediately: if
\(\alpha_{\bar k}=\operatorname{Int}(\bar g_{\bar k})\circ\sigma_{\bar k}\), then
\(\bar g_{\bar k}^{-1}\bar k\) acts on
\(\underline{\mathbf G}_x^\circ\) by the pinning-preserving automorphism
\(\sigma_{\bar k}\).  Hence
\(\bar g_{\bar k}^{-1}\bar k\in\overline G_{x,\cP}\), and so
\(\bar k\in\overline G_x^\circ\cdot\overline G_{x,\cP}\).  The reverse inclusion
is tautological.  This proves the rational pinned-component condition for the
full quotient.
\end{proof}

\begin{remark}
\label{rmk:full-quotient-not-connected-quotient}
In the application below, the enriched disconnected Jordan decomposition theorem
is used only to make the finite Jordan-decomposition input canonical for the
full finite quotient \(\overline G_x=G(F)_x/G(F)_{x,0+}\).  The connected quotient
\(\underline{\mathbf G}_x^\circ\) enters because the pinning is a pinning of the
identity component, as in Hypothesis~\ref{hyp:rational-pinned-component-condition}.
Lemmas~\ref{lem:parahoric-component-abelian} and
\ref{lem:parahoric-component-action-pinned} verify precisely the two additional
inputs needed to use Theorem~\ref{thm:disc-JD-bijection} for these full parahoric
quotients.  This use of the disconnected theorem concerns the finite
Jordan-decomposition labels and the finite-field parameter comparison; it does
not identify the finite crossed-product factor in the AFMO Hecke algebra with a
unipotent crossed-product factor.
\end{remark}

\subsection{Set-up and statement}

Let \(F\) be a non-archimedean local field with residue field \(\ff\) of
characteristic \(p\), and fix a coefficient field \(\C\) of characteristic
\(0\).  Let \(G\) be a connected reductive group defined over \(F\).  We assume
the standard tameness hypotheses under which the depth-zero type theory of
\cite{AFMO24a,AFMO24b} applies, as recalled in \cite[\S1]{ohara2025}:
\begin{equation}\label{eq:tameness}
G \text{ splits over a tamely ramified extension of }F,
\qquad
p\nmid |W_{\mathrm{abs}}(G)|.
\end{equation}
Fix once and for all a global pinned datum \(\cP\) for \(G\) in the sense used
above.  Whenever a parahoric quotient occurs in the construction, we use the
induced pinning \(\cP_x^{\cP}\) on its connected reductive quotient.

Let \(\mathfrak s=[M,\sigma]_G\) be a depth-zero inertial equivalence class for
\(G(F)\).  Let \((K_{x_0},\rho_{x_0})\) be the depth-zero \(\mathfrak s\)-type of
\cite{AFMO24a}, as recalled in \cite[\S4.2]{ohara2025}, and denote its Hecke
algebra by
\[
\cH_{\mathfrak s}(G)
:=
\cH\bigl(G(F),(K_{x_0},\rho_{x_0})\bigr).
\]
By \cite[Theorem~4.2.2]{ohara2025}, equivalently by
\cite[Theorem~5.3.6]{AFMO24a}, this algebra has the AFMO--Morris presentation
\begin{equation}\label{eq:AFMO-depth0-structure}
\cH_{\mathfrak s}(G)
\simeq
\C[\Omega(\rho_M),\mu]\rtimes
\cH_{\C}\!\bigl(W(\rho_M)_{\mathrm{aff}},q\bigr).
\end{equation}
Here \(W(\rho_M)_{\mathrm{aff}}\) is the affine Weyl group attached to the
depth-zero datum, \(q\) is the corresponding parameter function,
\(\Omega(\rho_M)\) is the stabilizer of the distinguished chamber in the relevant
extended symmetry group, and \(\mu\) is the AFMO \(2\)-cocycle appearing in the
twisted group algebra.

From \((K_{x_0},\rho_{x_0})\), Ohara constructs in \cite[\S4.3]{ohara2025} a
connected reductive group \(G_\theta\) splitting over an unramified extension of
\(F\), together with a unipotent type \((K_{\theta,x_0},u_{x_0})\) for
\(G_\theta(F)\).  Ohara's corresponding unipotent type has Hecke algebra
\begin{equation}\label{eq:AFMO-unip-structure}
\cH_{\mathrm{unip}}(G_\theta)
\simeq
\cH_{\C}\!\bigl(W(u_{\rho_M})_{\mathrm{aff}},q_\theta\bigr),
\end{equation}
as in \cite[Theorem~4.3.7]{ohara2025}.  In particular, the unipotent
Hecke algebra appearing in Ohara's comparison is the affine Hecke algebra
factor; it is not enlarged by the source-side twisted group algebra.

\begin{theorem}[Pinned comparison of affine Hecke algebra factors]
\label{thm:padic-depth0-to-unipotent-hecke}
Assume \eqref{eq:tameness}.  With notation as above, the fixed pinning \(\cP\)
canonically determines an isomorphism of affine Hecke algebras
\[
\Psi_{\cP}:
\cH_{\C}\!\bigl(W(\rho_M)_{\mathrm{aff}},q\bigr)
\xrightarrow{\ \sim\ }
\cH_{\C}\!\bigl(W(u_{\rho_M})_{\mathrm{aff}},q_\theta\bigr).
\]
Equivalently, using the AFMO--Morris presentation of \(\cH_{\mathfrak s}(G)\), one
may rewrite the depth-zero Hecke algebra as
\[
\cH_{\mathfrak s}(G)
\simeq
\C[\Omega(\rho_M),\mu]\rtimes_{\Psi_{\cP}}
\cH_{\C}\!\bigl(W(u_{\rho_M})_{\mathrm{aff}},q_\theta\bigr),
\]
where the finite group \(\Omega(\rho_M)\), the cocycle \(\mu\), and the crossed-product
structure are the original source-side ones, with the action transported across
\(\Psi_{\cP}\).  This statement does not assert, and in general should not be
read as asserting, an isomorphism
\[
\cH_{\mathfrak s}(G)
\cong
\cH_{\mathrm{unip}}(G_\theta).
\]
The affine-factor isomorphism preserves the standard anti-involutions on the two
affine Hecke algebra factors.
\end{theorem}

\subsection{Proof}
\begin{proof}
The two crossed-product presentations \eqref{eq:AFMO-depth0-structure} and
\eqref{eq:AFMO-unip-structure} show where the comparison can be made.  Ohara's
comparison theorem \cite[Theorem~4.4.1]{ohara2025} identifies the affine
reflection data on the two sides:
\[
        H_{\mathcal K\text{-rel}}=H_{\mathcal K_\theta\text{-rel}},
        \qquad
        W(\rho_M)_{\mathrm{aff}}
        \cong
        W(u_{\rho_M})_{\mathrm{aff}},
\]
and the parameter functions agree under this identification,
\[
        q=q_\theta .
\]
Consequently one obtains an isomorphism of affine Hecke algebra factors
\begin{equation}\label{eq:aff-hecke-iso}
\Psi_{\cP}:
\cH_{\C}\!\bigl(W(\rho_M)_{\mathrm{aff}},q\bigr)
\xrightarrow{\ \sim\ }
\cH_{\C}\!\bigl(W(u_{\rho_M})_{\mathrm{aff}},q_\theta\bigr),
\end{equation}
characterized by sending the standard generator attached to a simple reflection
to the standard generator attached to the corresponding simple reflection.

We now explain the role of the pinned Jordan decomposition in this comparison.
The proof of \cite[Theorem~4.4.1]{ohara2025} reduces the equality of the affine
Hecke parameters to finite-field equalities for the parameters attached to
rank-one Harish--Chandra inductions.  In the notation of the present manuscript,
this is exactly the kind of equality supplied by
Theorem~\ref{thm:finite-JD-preserves-q-parameter}: if \(\tau\) is a cuspidal
character in a Lusztig series and \(u_\tau\) is its unipotent Jordan correspondent,
then
\[
q\!\left(R_L^G(\tau)\right)
=
q\!\left(R_{H_L}^{H}(u_\tau)\right).
\]
When the finite group arising from a parahoric quotient is the full quotient
\(\overline G_x=G(F)_x/G(F)_{x,0+}\), rather than only its identity component,
Lemmas~\ref{lem:parahoric-component-abelian} and
\ref{lem:parahoric-component-action-pinned} show that the hypotheses of the
enriched disconnected theorem, Theorem~\ref{thm:disc-JD-bijection}, are satisfied.
Thus the finite Jordan-decomposition label used in the finite-field parameter
comparison is fixed by the chosen pinning \(\cP\), with the source Clifford class
retained as part of the enriched label whenever the full quotient is disconnected.
This removes the non-canonical finite Jordan-decomposition choice from the affine
parameter comparison.

This proves the claimed pinned affine-factor isomorphism \eqref{eq:aff-hecke-iso}.
Combining \eqref{eq:aff-hecke-iso} with the source-side AFMO--Morris presentation
\eqref{eq:AFMO-depth0-structure} gives
\[
\cH_{\mathfrak s}(G)
\simeq
\C[\Omega(\rho_M),\mu]\rtimes_{\Psi_{\cP}}
\cH_{\C}\!\bigl(W(u_{\rho_M})_{\mathrm{aff}},q_\theta\bigr).
\]
Here \(\rtimes_{\Psi_{\cP}}\) means that the original action of
\(\Omega(\rho_M)\) on
\(\cH_{\C}(W(\rho_M)_{\mathrm{aff}},q)\) is transported to the identified affine
factor by \(\Psi_{\cP}\).  The twisted group algebra is still
\(\C[\Omega(\rho_M),\mu]\).  No comparison with
\(\C[\Omega(u_{\rho_M}),\mu_\theta]\) is being asserted.

Finally, the compatibility with anti-involutions follows on the affine factors
because \(\Psi_{\cP}\) matches the standard generators attached to corresponding
simple reflections, and the structural AFMO--Morris isomorphism
\eqref{eq:AFMO-depth0-structure} preserves the standard anti-involution by
\cite[Theorem~4.2.2]{ohara2025}.  This proves the theorem.
\end{proof}

\begin{remark}[No full unipotent Hecke algebra isomorphism]
The theorem above is intentionally weaker than the statement it replaces.  In
general one should not identify
\[
\C[\Omega(\rho_M),\mu]\rtimes
\cH_{\C}\!\bigl(W(\rho_M)_{\mathrm{aff}},q\bigr)
\]
with
\[
\C[\Omega(u_{\rho_M}),\mu_\theta]\rtimes
\cH_{\C}\!\bigl(W(u_{\rho_M})_{\mathrm{aff}},q_\theta\bigr).
\]
Ohara's theorem identifies the second, affine Hecke algebra factors.  It does not
identify the finite stabilizer groups \(\Omega(\rho_M)\), \(\Omega(u_{\rho_M})\),
nor the cocycles \(\mu\), \(\mu_\theta\).  Therefore the correct consequence is the
source-side crossed product with the affine factor rewritten in unipotent terms,
not an isomorphism with the full Hecke algebra of the unipotent type.  If in a
special situation one has an additional compatible identification of the finite
stabilizer groups and cocycle classes, then the affine-factor comparison may
upgrade to a full crossed-product comparison; such an additional hypothesis is
not part of the present theorem.
\end{remark}

\begin{remark}
One may view Theorem~\ref{thm:padic-depth0-to-unipotent-hecke} as a canonical-up-to-pinning
refinement of the affine-Hecke-algebra comparison of \cite[Theorem~4.4.1]{ohara2025}.
The finite-field ingredient in \cite{ohara2025} uses Jordan decomposition in
Lusztig series.  Here that input is replaced, for the finite quotients relevant
to the depth-zero construction, by the pinned canonical bijection of the present
paper, and by the enriched pinned bijection of Theorem~\ref{thm:disc-JD-bijection}
when the full quotient is disconnected.
\end{remark}
\section*{Acknowledgements}

The authors are grateful to Maarten Solleveld and Amoru Fujii for their
careful reading of earlier versions of this paper.  Their questions,
comments, and corrections pointed out several errors and gaps in the
arguments, and led to significant improvements in the formulation and
proofs of the results presented here.

\bibliographystyle{alpha}
\bibliography{JD-references}
\end{document}
